% Generated by scripts/build_guide.py from reviewed-section paragraphs.
\def\EditionMode{3}
\def\EditionLabel{A guide for the reader}
\def\EditionStatus{WORKING GUIDE: 99 completed section/front-matter paragraphs; 374 of 706 source pages reviewed}
\documentclass[10pt,oneside]{book}
\usepackage{fontspec}
\setmainfont{Latin Modern Roman}
\setsansfont{Latin Modern Sans}
\usepackage{amsmath,amssymb,mathtools}
\usepackage{unicode-math}
\setmathfont{Latin Modern Math}
\usepackage{xcolor,tikz-cd,lua-ul,marginnote,pdflscape}
\usetikzlibrary{arrows.meta,decorations.pathmorphing,calc,shapes.geometric}
\tikzcdset{ampersand replacement=\&}
\tikzcdset{every label/.append style={text=black}}
\usepackage[a4paper,left=18mm,right=45mm,top=18mm,bottom=22mm,marginparwidth=28mm,marginparsep=7mm]{geometry}
\usepackage[unicode,colorlinks=true,urlcolor=black,linkcolor=black]{hyperref}
\usepackage{bookmark}
\setlength{\parindent}{0pt}
\setlength{\parskip}{4pt plus 1pt minus 1pt}
\setlength{\emergencystretch}{1.5em}
\raggedbottom
\definecolor{PSHand}{RGB}{25,58,132}
\definecolor{PSDelete}{RGB}{146,37,37}
\definecolor{PSTypedOld}{RGB}{95,95,95}
\ifnum\EditionMode=2\colorlet{blue}{black}\else\colorlet{blue}{PSHand}\fi
\RenewDocumentCommand{\textcircled}{m}{\tikz[baseline=(PSCircle.base)]\node[draw,circle,inner sep=.5pt,line width=.4pt](PSCircle){#1};}
\newif\ifPSInApparatus
\newcounter{psnote}
\newcounter{psvariant}
\newcounter{pslinkmark}
\newif\ifPSInLink
\newsavebox{\PSNoteBox}
\newdimen\PSNoteLimit
{\footnotesize\global\PSNoteLimit=3\baselineskip}
\newcommand{\CurrentScanURL}{}
\newcommand{\CurrentKey}{}
\newcommand{\CurrentScan}{}
\newcommand{\pssource}[3]{\psflushmargins\clearpage\gdef\CurrentKey{#1}\gdef\CurrentScan{#2}\gdef\CurrentScanURL{#3}\setcounter{psvariant}{0}\typeout{PS-SOURCE: #1}}
\newcommand{\psdivision}[2]{\pdfbookmark[0]{#1}{#2}\par\noindent{\large\bfseries #1\par}\medskip}
\newcommand{\pslinkmarker}[2]{\stepcounter{pslinkmark}\latelua{texio.write_nl('log','PS-LINK-SHIP|#1|\CurrentKey|\arabic{pslinkmark}|#2')}}
\newcommand{\pslink}[2]{\ifPSInLink\pslinkmarker{#1}{nested}#2\else\begingroup\PSInLinktrue\href{\CurrentScanURL}{\pslinkmarker{#1}{outer}#2}\endgroup\fi}
\newcommand{\scanpage}[1]{\ifnum\EditionMode=2\leavevmode\marginpar{\scriptsize\sffamily\pslink{scanpage}{scan #1}}\else\par\noindent{\small\sffamily\pslink{scanpage}{[scan #1]}}\par\medskip\fi}
\newcommand{\mspage}[1]{\ifnum\EditionMode=2\leavevmode\marginpar{\scriptsize\sffamily p.\,#1}\else{\small\sffamily[p.\,#1]}\fi}
\newcommand{\hand}[1]{{\ifnum\EditionMode=2\color{black}\else\color{PSHand}\fi #1}}
\let\over\relax
\NewDocumentCommand{\over}{O{} m m}{{\textcolor{PSTypedOld}{#2}\ensuremath{^{\!\nearrow}}\if\relax\detokenize{#1}\relax\hand{#3}\else #3\fi}}
\NewDocumentCommand{\inserted}{O{} m}{{\if\relax\detokenize{#1}\relax\hand{\ensuremath{\langle}#2\ensuremath{\rangle}}\else\ensuremath{\langle}#2\ensuremath{\rangle}\fi}}
\newcommand{\deleted}[1]{{\color{PSDelete}\ifmmode\text{[del.: }#1\text{]}\else[\textsf{del.:} #1]\fi}}
\newcommand{\xout}[1]{{\color{PSTypedOld}\ifmmode\text{[x: }#1\text{]}\else[\textsf{x:} #1]\fi}}
\let\PSMathUnderline\underline
\RenewDocumentCommand{\underline}{+m}{\ifmmode\PSMathUnderline{#1}\else\underLine{#1}\fi}
\NewDocumentCommand{\uline}{+m}{\underline{#1}}
\NewDocumentCommand{\uuline}{+m}{\ifmmode\PSMathUnderline{\PSMathUnderline{#1}}\else\underLine{\underLine{#1}}\fi}
\newunderlinetype\PSWaveLine{\leaders\hbox{\lower.7ex\hbox{\tikz[x=1pt,y=1pt,baseline=0pt]{\draw[line width=.3pt] (0,0) .. controls (1,1) and (2,1) .. (3,0) .. controls (4,-1) and (5,-1) .. (6,0);}}}}
\NewDocumentCommand{\uwave}{+m}{{\PSWaveLine #1}}
\newcommand{\unsure}[1]{\pslink{unsure}{[#1\ifmmode\,?\else\,?\fi]}}
\newcommand{\illegible}[1]{\ifmmode\text{\pslink{illegible}{[illeg.\ #1]}}\else\pslink{illegible}{[illeg.\ #1]}\fi}
\NewDocumentCommand{\diagram}{+m}{\ifmmode\vcenter{\hbox{#1}}\else\par\medskip{\centering #1\par}\medskip\psflushmargins\fi}
\newenvironment{diagramblock}{\par\begin{center}}{\end{center}\par}
\newcommand{\PSMathNotes}{}
\makeatletter
\NewDocumentCommand{\psplacednote}{m +m}{%
      \sbox{\PSNoteBox}{\parbox{\marginparwidth}{\footnotesize\raggedright\hand{\textsf{m#1.} #2}}}%
      \ifdim\dimexpr\ht\PSNoteBox+\dp\PSNoteBox\relax<\PSNoteLimit
        \leavevmode\marginpar{\usebox{\PSNoteBox}}%
      \else
        \par\nobreak\smallskip\hrule\nobreak\smallskip
        {\footnotesize\textsf{Marginal note m#1.}\quad\hand{#2}\par}
        \nobreak\smallskip\hrule\smallskip
      \fi}
\newcommand{\psflushmargins}{\ifmmode\else\ifinner\else\ifx\PSMathNotes\@empty\else\let\PSTempNotes\PSMathNotes\global\let\PSMathNotes\@empty\PSTempNotes\fi\fi\fi}
\NewDocumentCommand{\margin}{+m}{%
  \ifPSInApparatus\ifmmode\hand{\text{[marg.: }#1\text{]}}\else\hand{\textsf{[marg.: }#1\textsf{]}}\fi
  \else\stepcounter{psnote}%
    \ifmmode
      \text{\textsuperscript{m\arabic{psnote}}}%
      \begingroup\edef\PSQueue{\endgroup\noexpand\g@addto@macro\noexpand\PSMathNotes{\noexpand\psplacednote{\arabic{psnote}}{\unexpanded{$#1$}}}}\PSQueue
    \else
      \textsuperscript{m\arabic{psnote}}\psplacednote{\arabic{psnote}}{#1}%
    \fi
  \fi}
\AtBeginDocument{\everymath\expandafter{\the\everymath\aftergroup\psflushmargins}\everydisplay\expandafter{\the\everydisplay\aftergroup\psflushmargins}}
\makeatother
\newcommand{\psnumber}[2]{\expandafter\gdef\csname psnum@#1\endcsname{#2}}
\NewDocumentCommand{\var}{+m +m m m}{\ifnum\EditionMode=0 #1\else\ifnum\EditionMode=1\ifmmode\text{\textsuperscript{\tiny\csname psnum@#4\endcsname}}\else\textsuperscript{\tiny\csname psnum@#4\endcsname}\fi\fi #2\fi}
\newcommand{\psapparatusstart}{\par\bigskip\hrule\smallskip\PSInApparatustrue\begingroup\footnotesize\raggedright\textsf{Apparatus for scan \CurrentScan}\par}
\newcommand{\psapparatusend}{\par\endgroup\PSInApparatusfalse}
\NewDocumentCommand{\psentry}{m m m +m +m m}{\par\noindent\textbf{#2. #3.}\quad\textit{Typescript:} #4\quad\textit{Corrected:} \if\relax\detokenize{#5}\relax\textnormal{[omitted]}\else #5\fi\quad #6\par}
\hyphenation{Gro-then-dieck abel-ian-iza-tion}

\begin{document}
\hypersetup{pageanchor=false}
\begin{titlepage}\raggedright
{\Large Alexander Grothendieck\par}\vspace{12mm}
{\Huge Pursuing Stacks\par}\medskip{\LARGE 1983\par}\vspace{12mm}
{\Large A guide for the reader\par}\vspace{8mm}
{\large\EditionStatus\par}\vfill
Université de Montpellier, Grothendieck archive, cotes 134-2--134-8.\par
Including the appendix translation by R. Brown, cote 134-2, sites 115--143.\par
\end{titlepage}\hypersetup{pageanchor=true}
\chapter*{Editorial statement}

AI-transcribed by Claude Opus 5.5 (Anthropic, effort high); checked against the scans, and this edition built, by GPT-6 Astra (OpenAI, ChatGPT Pro mode); prompts and records by Claude (Anthropic).

The witness consulted and not reproduced is the LaTeX edition of Mateo Carmona with the collaboration of Ulrik Buchholtz (arXiv 2111.01000 v2, 2021), extended from the 2015 edition by ``the scrivener'', CC0 1.0.

The translation in the appendix is R. Brown's, prepared with assistance from L. Breen and J-L. Loday, as his note says; the rights in it belong to its translator or his heirs.

The rights in the original work belong to its author or the author's representatives or heirs. Where we hold rights, in the AI-produced transcription, translation, corrections and typesetting, they are dedicated to the public domain under CC0 1.0. This edition is published for educational, non-profit purposes, and we claim no rights we do not hold. If you hold rights in this work and object to its publication here, please open a GitHub issue and we will take it down.


\section*{Using this guide}
Each paragraph describes a section already checked against both supplied scan halves. It identifies the section's task, terminology, dependencies, revisions and consequential witness differences. Mathematical cautions are editorial observations, distinguished from the author's statements; they do not alter the readers. Archive sites and internal typescript folios are kept distinct, especially where the initial letter has been reordered or renumbered. The account is provisional wherever the relevant later section has not yet been reviewed.

The guide is incomplete while source collation continues. No paragraph for an unreviewed section has been supplied from the arXiv witness alone. The completed paragraphs follow the original section order; gaps in the current working guide are intentional and will close as their sections are checked.
\chapter*{Section guide}
\subsection*{Front matter: plans, contents, and dedication}
The front matter (cote 134-2, sites 1--13) should be read as evidence of a developing book: the title page and chapter-dependency diagram precede several layers of contents, including cancelled headings, renumbered sections, and an entirely handwritten replacement page. The diagram uses the author's chapter labels and places derivators within the projected sequence; its heavily retraced connections should not be treated as a settled logical dependency chart. The cancelled chapter-IV block on site 7 and the replacement handwritten contents on site 8 show the reorganization of the asphericity discussion. Sites 9--12 then project the progression through homology and abelianization, schematization, and linearization: the proposed comparison between schematic and discrete homotopy theory leads to the ``breakdown of an idyllic picture'' of section 119, then to revised programs and questions about the six operations in sections 124 and 134--139. These titles announce later developments; they do not themselves establish the claims. In particular, the pair of arrows listed for section 121 points from homotopy types to abelian homotopy types above and back below, with handwritten arrows and accents over typed objects and letters. Several continuations on site 11---\emph{topos}, \emph{bundles}, and \emph{valences}---are typing guided by handwritten curves, a material distinction lost in the initial transcription. The doubled \emph{ff} in the title for section 131 is visibly typed. The French dedication on site 13 retains a cancelled address to former pupils above the dedication to Dodek Shapiro, together with its cancelled and unresolved words. It remains French in all readers. The contents' cancellations and the final handwritten ``to be completed'' are particularly useful reminders that the source records successive intentions, including abandoned formulations, rather than a retrospectively regularized exposition (Grothendieck, 1983, cote 134-2, sites 1--13).

\subsection*{The letter to Quillen as a whole (cote 134-2, sites 14--36)}
The letter and its dated continuations supply the motivating problem for the subsequent notes: describe homotopy types by higher groupoid structures, relativize those structures over topoi to obtain stacks, and organize their operations through a universal categorical pattern. The progression is from the need for coherent higher composition, through globular source and target data and categories of standard amalgamations, to successive freely adjoined fillers and their comparison with topological or model-category realizations. The name ``coherator'' arrives only after substantial revisions to the proposed construction. The letter's mathematical force lies partly in those revisions: strict formulas give way to specified homotopies; the supposed universality of a category of actual disc maps becomes doubtful; economical choices of operations are treated as secondary to higher comparison; and unrestricted closure under fiber products is retracted in the postscript. The main text later returns to the program in sections 18--19 and develops a different, extensive investigation through test categories, asphericity, localizers, and integration. The earlier Breen correspondence reproduced in the appendix explains the topos and coefficient background invoked at the start and near the end of the letter. The archive places centered folios 4 and 3 in the reverse order at sites 16--17; the readers retain that archive order, while the section paragraphs identify the mathematical continuation. The witness's expanded marginal notes, silent index corrections, omitted signatures, and retained cancelled passages must therefore be distinguished from the source's own stages. The guide records the local formula cautions in the relevant section paragraphs, without replacing the author's provisional mathematics by a reconstructed finished theory.

\subsection*{1. Foundational questions and tame topology (letter pp. 1--1')}
The letter dated 19 February 1983 starts from the papers sent by Ronnie Brown and recalls the 1975 correspondence with Larry Breen as a program joining homotopical and homological algebra over topoi. Before pursuing that program, Grothendieck argues for an even more urgent foundation for ``tame'' topology and stratified tame spaces, contrasting the conceptual difficulty with the inertia of established habits. The next section returns to the homotopical program; section 14 later resumes the criticism of the category of topological spaces and its inadequacy for isotopy and stratification. The opening is therefore a statement of priorities, not a theorem about a proposed category. Its cancellation history matters: on cote 134-2, site 14, the addressee is crossed out, and the long passage about the recipient's axiomatic homotopical algebra and lecture notes is cancelled and replaced by a shorter handwritten phrase. The arXiv witness prints that cancelled continuation as live text. It also substitutes ``as for schemes'' for the clearly handwritten ``or for schemes''. The section's descriptive title in the witness is absent from the scanned opening, which gives only its number. The continuation on site 15 concludes the comparison with programs of homological, homotopical, or topological algebra.

\subsection*{2. Higher groupoids and coherent composition (letter, sites 15--18)}
Returning to the Breen program recalled in section1, Grothendieck proposes that $n$-truncated homotopy types be described by $n$-groupoids, with a fundamental $n$-groupoid, geometric realization, and an $n$-equivalence comparison; relativization over a topos should yield $n$-gerbs and $n$-stacks as noncommutative cohomological coefficients. The assertion that every $n$-groupoid should be $n$-equivalent to a fundamental one is present on site15 but omitted from the arXiv witness, which joins the surrounding clauses. The attempted globular description then introduces sets $F_i$, source, target, identity and inversion maps, and composition at different levels. Strict associativity is rejected as insufficient for the intended homotopy types; equalities are to become specified homotopies, whose own coherence leads from ``secondary'' to ``ternary'' and further structure, illustrated by the pentagon at the start of site18. Section3 asks whether this can support a model-category description; sections18--19 later revise the coherator and realization proposals. The internal reading order here is site15, then site17 (centered folio3), then site16 (folio4), then site18 (folio5); the readers preserve the archive's different order. Several formulas require caution. On site17 the first factor of the composition domain has an overprinted $t$, with the underlying letter still doubtful, where the witness gives $s$; the marginal identity notation is $id_a$, not $d_a$. The handwritten target formula ends with $s^i_1(u)$ in its second right-hand factor, an apparent source/target slip. On site16 the inversion formula places separate inversion marks on both left-hand factors and appears tautological, although the first right-hand letter remains overwritten and doubtful; the witness substitutes a question mark for that right side. The handwritten interchange range also permits equal levels, $1\leq\alpha\leq\nu$: for a one-object group at equal levels, the displayed interchange would force the two middle factors to commute, so this range is too broad for arbitrary fundamental groups. These source readings remain unchanged. The author replaces ``implying'' by ``expressing'' beneath the compatibility conditions, a correction the witness does not apply. The resulting discussion is a developing proposal for coherence, and these local defects should not be treated as established definitions or silently repaired.

\subsection*{3. The model-category question (letter, sites 18--19)}
Grothendieck asks Quillen whether an appropriate notion of higher category already exists, recalling a description by multiple semisimplicial sets but distinguishing strict associativity from the coherent structure sought in section 2. He proposes two connections with homotopical algebra: the conjectural category of infinity-groupoids should model ordinary homotopy types, and a fundamental infinity-groupoid construction should be available in any model category from a cylinder object of its final object. The latter is presented as informative chiefly in the nonpointed case, and the handwritten ending explicitly identifies the cylinder as playing the role of the unit segment. Section 4 immediately reopens the skeletal data of that proposed construction; section 14 returns to the limitations of topological-space models, and the later interval discussion develops the cylinder intuition. These are proposals and questions, without a completed definition or asserted model-structure theorem. On site 18, the handwritten replacement of ``your'' by ``Quillen's'' in the reference to axiomatics is not applied by the arXiv witness, and its ``contents'' misreads the typed ``comments''. Site 19 closes the first letter with a handwritten signature and is numbered 5-prime on the scan, where the witness labels it 3-prime. The ensuing dated continuation on site 20 begins a substantive revision rather than simply continuing an established definition.

\subsection*{4. Skeletal data and successive structure (letter, sites 20--21)}
In the continuation dated 20 February, Grothendieck revises the undifferentiated ``primary'' structure proposed in section 2: the skeletal datum is only the family of sets with source and target maps, while inversion, degeneration and first composition belong to structure added subsequently. Equivalences of ordinary groupoids motivate replacing strict involutions by pseudoinvolutions, and choices of interval maps explain why degeneracies and path composition should also be nonprimitive. Godement's next composition is then secondary, and the further composition levels successively ternary and higher. This ordering prepares the tower of universal categories in section 5 and makes explicit a change of mind: the text says that the preceding day's operations were hastily put into one pot. Two details deserve care. On site 20 the formal maps $s_o,t_o$ are described as coming from $F_0$ to $F_{-1}$ but are assigned ``$i=1$''; the evident index mismatch is retained. On site 21 the handwritten relocation yields ``are rather'', while the witness retains ``rather are''. The replacement ``the real'' on that page is typed above the line, not handwritten as the initial transcription classified it. The source says ``all these operations'', where the witness shortens this to ``all the operations'', and retains the two handwritten marginal composition symbols absent from that witness. The overwritten folio at the head of site 21 remains uncertain between 7 and 4; the archive order and all scan links are unaffected.

\subsection*{5. A tower of universal categories (letter, sites 21--24)}
Grothendieck proposes a sequence of categories $C_0,C_1,\ldots$ with common objects and progressively added arrows, whose direct limit should classify the proposed infinity-groupoid structures by functors preserving specified limits. The discussion first compares universal structures in categories with finite products, finite inverse limits, and the stronger exactness available in topoi; it then discards the separate finite-product requirement while writing site 23, restricting the common objects to ``standard'' iterated fiber products determined by source and target. This is the categorical organizing device promised in section 4. Section 6 gives it a topological model; the later coherator discussion revises how the successive operations should be adjoined. The text expressly presents the strict increase at every stage as a conviction still needing proof. On site 23 the letter in ``exist in F'' is visibly F, although the ambient category is C and the witness prints C; the source letter is retained as an apparent local slip. The iterated-source indices on that page contain typed 1-shaped bases with handwritten ell/chi-like corrections and remain marked, rather than being settled solely by the witness's ell. The live insertion says that the word ``iterated'' will henceforth be dropped, while the cancelled marginal version says ``qualification''; the witness combines the latter wording with the live text. On site 24 the marginal note assigns order ell to the composition bearing lower index ell, where the witness gives that composition lower index 2. Some of this note's lettering and cancellation status remain doubtful, but the difference of indices should not be silently harmonized. The section's revisions distinguish a proposed universal construction from a finished coherence theorem.

\subsection*{6. Discs and a proposed universal co-groupoid (letter, sites 24--25)}
The tower of section 5 receives its motivating topological interpretation: a space $X$ supplies sets of maps from discs $D_i$, and the source and target maps are induced contravariantly by the two boundary embeddings. Compositions and further operations are to arise from maps between discs and their prescribed amalgamated sums, producing a co-infinity-groupoid whose transform under $h_X$ is the desired infinity-groupoid. Grothendieck then proposes that the category $B_\infty$ generated by these cells, sums and arrows should be dual to the universal category $C_\infty$. The expected universality would make the category independent of the successive arbitrary choices up to unique isomorphism, while he insists that the formal construction of $C_n$ remains meaningful without that expectation. Section 7 begins to specify the disc maps and their coglobular identities. Site 25 records a significant change of mind: the confident parenthesis ``I believe it is'' acquires handwritten exclamation marks, and a marginal addition dated 23 February says ``But no longer now'' and refers to page 10, section 11. The crossing pencil hatching leaves the note's cancellation status uncertain; the readers retain that doubt. The arXiv witness paraphrases this retraction, omits its explicit page reference, and prints the cancelled article before ``various'' amalgamated cells. On site 24 the unit-segment formula retains the typed interlinear $I=$ and the live $D_1$, an identification shortened in the witness. These distinctions preserve both the proposed model and the author's later withdrawal of confidence in its universality.

\subsection*{7. Coglobular discs and standard amalgamations (letter, sites 25--27)}
The two hemisphere maps between successive discs give the cellular skeleton needed in section 6. Grothendieck describes the boundary identifications and the decomposition into open cells, then interprets composition at level ell as a choice of extension from a boundary sphere into an amalgamated sum. Contractibility supplies the proposed existence argument; its inductive dependence on lower-level operations explains the hierarchy introduced in section 4. He next seeks a combinatorial account of ``standard'' amalgamations by attaching cells along prescribed subcells, dualized into relations among iterated sources and targets in a product of the objects $F_i$. Section 8 recasts these data as finite ordered sets. Two scan readings materially alter the comparison witness: on site 26 the range is strictly $0\leq j<i$, and the open-cell rings are present on $D_n$, both $S_j$ symbols, and the argument of the positive map. The corresponding negative-map argument has no visible ring; this apparent local omission is retained. The same page draws undirected circle examples, where the witness adds arrowheads. On site 27 the proposed restriction on pairwise cell intersections receives a marginal counterexample. Its three regions are cut out by two arcs sharing their endpoints, so $f\cap f''=\{s,s'\}=S_0$; a central three-sector diagram would destroy that counterexample and has been corrected in the checked transcription. The note says ``see this'', with its first word and hatching status still uncertain, whereas the witness substitutes explanatory prose. The product formula uses $F_{n(i)}$, not the witness's $D_{n(i)}$, and the handwritten sign variables are epsilon and epsilon-prime throughout. These qualifications keep the proposed combinatorial scheme and its own marginal objection visible together.

\subsection*{8. The finite ordered sets underlying $B_o$ and $C_o$ (letter, site 28)}
The dated continuation of 21 February replaces each amalgamated cell complex by its finite ordered set of cells, with distinguished subposets isomorphic to the standard $2n+1$-element posets $J_n$. The incidence relation and the specified identifications encode the combinatorial structure sought in section 7; their opposite category is proposed as $C_o$, the starting point for the inductive construction in section 9. The intersection $K_i\cap K_{i'}$ is expressly a closed subposet isomorphic to some $J_n$; the arXiv witness misprints its intersection sign as an arrow. Grothendieck then excludes the empty structure and claims stability under binary amalgamated sums, extending the claim to all $C_n$ and describing the desired functors by preservation of those sums. The page itself objects: the all-categories phrase has a wavering underline and a marginal note saying ``This is obviously false, see P.S. p.12'', with repeated calls to note 13. The pencil hatching leaves that note's status uncertain, but its wording is clear. The distinction between amalgamated sums in $B_o$ and limits in its opposite is also essential to reading this developing assertion; the readers preserve the claim rather than repairing its variance. The witness removes the all-categories phrase and rewrites the note with later qualifications referring to sections 13 and 18. The guide therefore treats the closure claim as a proposal already challenged on its own page, pending the later revisions.

\subsection*{9. Universal coherence and comparison between definitions (letter, sites 28--30)}
Conditions (A) and (B) propose that suitable parallel arrows in $C_n$ acquire a filler at the next stage and that $C_{n+1}$ be freely generated from such fillers subject to their prescribed boundaries. This turns the informal coherence program of sections 2 and 4 into a proposed universal construction, while section 10 will consider how economically to choose the new arrows. The phrase ``an isomorphism'' is visibly cancelled beneath the typed replacement ``a homotopy'' on site 29: the intended coherence is itself higher-dimensional, rather than an identification of all operations. Grothendieck nevertheless says that these conditions do not characterize $C_\infty$ up to equivalence. His sequence of seven mathematicians on site 30 explains the intended weaker independence: competing definitions should have comparison functors, the composites should be infinity-equivalent to the originals, different comparison choices should themselves admit equivalences, and comparisons between those equivalences should continue indefinitely. The promised finite mathematical statement of this story remains explicitly unwritten. Two local cautions concern the construction itself. The filler has target $F_{i+1}$ on site 28 but $F_i$ on site 29; the source mismatch is preserved. The marginal correction on site 29 says that amalgamated sums exist only in $C_o$, not the initial transcription's $C_\infty$, and another note calls the asserted stability of all $C_n$ inaccurate. The underlined demonstrative and abbreviated ``should'' in the first note, as well as its hatching status, remain uncertain. These source objections continue the revisions noted in section 8 and prevent the asserted universal construction from being read as a settled theorem. The arXiv witness retains the cancelled ``can'' instead of the live ``may'', omits the source article in ``at all a sense'', and supplies articles before ``same'' that the scan does not contain. None of those changes is imported into the readers.

\subsection*{10. Redundant operations and choices of coherence data (letter, sites 30--31)}
Grothendieck considers progressively narrower choices of fillers for the construction of section 9: one for every admissible pair, one for every isomorphism class of such situations, one only when no previous filler exists, and finally only for pairs that do not factor nontrivially through a smaller object. Even these restrictions leave many operations serving the same purpose. In particular, binary composition and all simultaneous $n$-ary compositions appear at the first stage; the second stage relates them by homotopies and the third relates composites of those homotopies. The comparison with two descriptions of a category with associative tensor composition gives a finite-dimensional analogue of the seven-mathematician story. Its claim about nonequivalence of the ordinary categories of algebraic structures is explicitly qualified by the typed insertion ``I believe''; it should not be read as denying the higher comparison the preceding section seeks. The marginal qualification ``even not in the grou-poid context'' remains uncertain in status because of pencil hatching, and the arXiv witness expands it into a different footnote mentioning Gr-categories. Grothendieck concludes that economy of the presentation is secondary and leaves the notion partly indeterminate, subject to (A), (B), and the description of $C_0$. Section 11 tests the alternative of selecting actual maps of topological cells. On site 31 the first variable is visibly capital $U_1$, while the witness prints $u_1$; this local notational mismatch is preserved. The initial transcription also mistook the live overprinted letters of ``operation'' for a deletion, which would have produced ``opeion'' in a corrected reader; the image check restores their proper typed-overprint status.

\subsection*{11. Allowable maps and the faithfulness question (letter, site 32)}
Grothendieck returns to the topological realization proposed in section 6, asking whether the universal category can be described as the opposite of a category of actual maps between amalgamated cells. He proposes subsets of allowable maps from discs, stable under composition and admitting extensions from their boundary spheres. These conditions model closure and the filler requirement, and elementary tameness conditions seem capable of enforcing them. They do not by themselves secure the freeness required by condition (B): actual composition of continuous maps might impose extra relations not forced by the universal algebra. Thus the stepwise universal construction gives a contravariant realization functor, but its faithfulness is neither evident nor, in the closing judgment, essential. This is the substantive withdrawal of confidence anticipated by the dated marginal retraction on site 25. The later comparison of competing realizations can therefore proceed without identifying the abstract category with its image in topological spaces. The arXiv witness drops K from ``every K a subset'' and changes the source index $C_{n-1}$ to $C_{n+1}$ when describing the universal properties of $C_n$; the latter source reading is clearly minus 1 and suits the stated inductive construction. The overprinted A/a in the reference following condition b) remains visually doubtful, but the reference concerns the earlier filler requirement; it should not be silently interpreted as the local closure condition a). The superscript mark on $S_n$ also remains unresolved, and the handwritten subscript in the phrase about a piecewise-linear structure is uncertain between i and n. These uncertainties are preserved in the readers rather than decided by the witness.

\subsection*{12. A co-groupoid in a model category (letter, site 33)}
The proposed disc construction is transferred from topological spaces to an arbitrary model category, as anticipated in section 3. Grothendieck chooses trivial fibrant objects $D_i$ and cofibrations from the successive boundaries, then seeks to extend the resulting realization functor from $B_o$ through the universal stages. He expects standard gluings of those objects to remain fibrant and trivial and proposes localization of the associated infinity-groupoids to recover the homotopy category. For semisimplicial sheaves over a topos, this is related to the operation called ``integration'' or ``sections'' of gerbs and stacks in the earlier Breen program. These expectations are explicitly presented as things still to be checked. In particular, the asserted stability of fibrancy under the indicated pushouts is not a consequence of the model axioms: in simplicial sets, glue the nerves of two contractible two-object groupoids along a common vertex. All three inputs are contractible Kan complexes and the inclusions are cofibrations, but a composable pair crossing the shared vertex lacks a filling 2-simplex in the pushout, so that pushout is not Kan. The mathematical assertion therefore remains as written and should not be silently converted into a cofibrancy assertion. The subsequent terminology changes in section 13 retain the conceptual program while abandoning an unqualified name for its choice-dependent objects. On this page the witness changes the source name (Ho) to Hot and the handwritten lower-case pi to capital Pi. The final two interlinear replacements of ``groupoids'' by ``gerbs'' are typed and are retained as such. The overprinted disc indices and the supposed equals correction after $S_o$ remain marked where the image does not securely separate their layers.

\subsection*{13. Coherators, intrinsic structure, and retractions (letter, sites 34--36)}
Grothendieck replaces the potentially misleading name ``infinity-groupoid'' for the auxiliary algebraic pattern by ``coherator'', reserving ``stacks'' for structures modelled on such a pattern. Building on the universal constructions of sections 9--12, he asks how much of the chosen globular objects, source and target maps, and filtration can be recovered from the abstract category. The successive answers on site 35 must be read with their retractions: the proposed automorphism groups $(\pm1)^i$ are challenged by the observation that the inversion operations need be neither involutions nor automorphisms; an irredundancy condition is then offered to recover rigidity of the whole globular system. The marginal objection dated 25 February says that the proposed intrinsic construction can give ``two large a subcategory'' $C_{n+1}$, and the closing marginal note on site 36 suggests supplying the primitive subcategory $C_o$ as extra structure. These are changes within a developing definition, not an established characterization theorem. The continuation in sections 18--19 revisits this program after the intervening reflections on topology. The final postscript, visibly dated 23.2 rather than the witness's 25.2., explicitly withdraws the claim that all the categories $C_n$ admit the fiber products previously asserted: the surviving statement concerns the fiber products already present in $C_o$, their preservation in $C_n$, and functors whose restrictions to $C_o$ preserve them. This resolves the objections on sites 28--29 without deleting the earlier claims. Three local mathematical slips remain in the text. On site 35, the intrinsic characterization of the source--target pair refers to chosen $F_i$ and $F_{i+1}$, although the preceding arrows have codomain $F_{i-1}$; the witness silently substitutes $i-1$. Later the second boundary composite uses an ununderlined $t$, and the final endomorphism is named $h$ while its boundary equalities use $f$. Neither notation is silently harmonized. The source's marginal notes have shorter, sometimes differently located wording than the witness footnotes: call 13 follows ``amalgamation'', call 19 rejects the automorphism claim, and call 20 qualifies the filtration. Two of those notes were hidden in comments in the initial LaTeX and have been restored to printable anchors. Their readable wording remains visible with uncertainty about crossing pencil hatching. The extra cluster after ``Maybe'' on site 36 and the second numeral in its body call remain unresolved, while both calls in the postscript read 13 followed by an exclamation mark. The typed interlinear ``later'' gives ``a later elaboration'', where the witness rearranges the phrase. The letter thus ends by stating the intended role of stacks as noncommutative cohomological coefficients while expressly leaving the coherator axioms open to further testing.

\subsection*{Introduction to the Breen letters (cote 134-2, site 37)}
The appendix introduces letters dated 5 February, 17 February, and 17--19 July 1975 as the earlier program recalled in the Quillen letter. Grothendieck explains that the first two are reproduced in French and the third in Ronnie Brown's English translation, prepared with assistance from Larry Breen and J.-L. Loday; he identifies his present interventions as comments, numbered subtitles, and corrections chiefly occasioned by the handwriting of the original. The first two letters develop the proposed connection between higher categories, stacks, and truncated homotopy types. The third broadens the program to duality, geometric local and global class field theory, tame topology, and the inverse system of indexed homotopy types associated with equisingular stratifications. These themes supply the historical setting for the later investigations of integration and tame topology without asserting that the 1983 formalism has already been established. Three witness differences materially affect that setting: the typed copy is dated 1975 rather than 1945; the intended phrase is ``fine homotopy type'', not ``finite homotopy type''; and the handwritten replacement of ``ordered set'' is ``inverse system'', not ``inverted system''. The source also specifies notes 22--31 and subtitles App.1--App.18, details shortened in the witness. The stepped pen cue begins a new paragraph for the third letter, while the retained cancellations and insertions distinguish the 1983 editorial introduction from the letters it presents.

\subsection*{Appendix 1: Picard categories and truncated complexes (cote 134-2, sites 38--40)}
The first part of the letter of 5 February 1975 asks what geometric objects, rather than their isomorphism classes alone, realize cohomology groups. Its four examples move from gerbes and their higher torsors to Picard stacks, stacks of $(\mathrm{Gr})$-categories, and Azumaya algebras. For a band $L$ with center $Z$, the $Z$-gerbes form a strict Picard 2-category governed by the degree-two truncation of $R\Gamma_X(Z)$; the $L$-gerbes form a pseudo-2-torsor, whose nonemptiness has an obstruction in $H^3(X,Z)$. Passing from global objects to the corresponding stacks shifts the description toward 2-gerbes and a degree-three truncation, anticipating Appendix 2's proposed interpretation of noncommutative complexes by higher categories. Picard stacks with invariants $M,N$ introduce the unusual $\mathrm{Ext}^2$ attached to $R\mathrm{Hom}(X(M),N)$, distinguished from the ordinary lower Ext groups and from the strict case governed by $R\mathrm{Hom}(M,N)$; the source range for those lower groups is $0\leq i\leq1$, whereas the witness extends it to 2. Likewise, the source repeatedly has $Z$-1-gerbes or $Z$-1-champ, not the witness's $Z$-$L$ expressions. For a group $G$ acting on $N$, the proposed invariants are relative cohomology groups in degrees 3 and 2 together with $Z^1(G,N)$, explicitly the cocycle group, which must not be replaced by $H^1$. The Brauer example assigns to an Azumaya algebra its gerbe of trivializations and distinguishes the Brauer subgroup of $H^2(X,\underline G_m)$ from the full cohomology group; this leads Grothendieck to question effective 2-descent for the proposed fibered 2-category. The middle displayed invariant on site 40 is literally $H^1(X,\underline H_m)$: the surrounding $\underline G_m$-gerbe construction and the adjacent degree-two and degree-zero terms indicate the intended coefficient $\underline G_m$, which the witness silently supplies, but the source $\underline H_m$ is retained in every reader. The geometric trivialization variable on that page is handwritten $\phi$, not an empty-set symbol. The change of emphasis from classes to structured objects, and then from global objects to stacks, provides the dependency needed for the ensuing Lefschetz discussion; the later notes 22--23 qualify the band and the complex $X(M)$ already cited here.

\subsection*{Appendix 2: Lefschetz theorems as a test for higher stacks (cote 134-2, sites 40--42)}
The preceding Picard examples suggest interpreting truncated complexes, explicitly as objects of derived categories, by strict Picard $n$-categories and then forgetting the additive structure to obtain candidates for noncommutative complexes. Grothendieck motivates this proposal through weak Lefschetz theorems in étale cohomology: depth conditions and the strongest available statements in degrees zero and one are naturally formulated for stacks, rather than merely for cohomology sets or fundamental groups. Site 41 distinguishes the prime-to-characteristic restrictions accompanying duality and smooth or generic base change from the unrestricted proper-case statements, and identifies proper base change, smooth base change, and generic cohomological properness as the indispensable tests of a useful formalism. The witness omits the constructibility hypothesis in the generic case, together with the transition explaining that smooth and generic base change serve the generic Lefschetz version; the omission materially weakens the stated assumptions. The handwritten étale-cohomology phrase belongs after ``coéfficients discrets'', and the cancelled phrase before ``les résultats les meilleurs possibles'' reads ``en fait'', not the initial transcription's tentative ``aussi''. Site 42 extends the proposed framework to higher stacks and complex analytic variants, while insisting that serious geometric theorems, rather than formal manipulation alone, must justify it. The concluding dissatisfaction with the definition of triangulated category prepares note 24 and the later concern with retaining structure lost in localization. Appendix 3 then supplies the promised heuristic link with simplicial homotopy theory; the Quillen letter and the main body return to that link through test categories and higher groupoids. The present passage states a program and its tests, not a constructed general theory of noncommutative derived categories.

\subsection*{Appendix 3: Higher groupoids and truncated homotopy types (cote 134-2, sites 42--45)}
The postscript proposes the homotopy hypothesis in its finite-dimensional form: an $n$-category in groupoids, up to $n$-equivalence, should encode a simplicial homotopy type with homotopy groups above degree $n$ discarded. Paths, homotopies of paths, and successive higher homotopies motivate $C_n(K.)$, with the top-dimensional arrows taken modulo homotopy; the failure of strict associativity is explicitly a reason to seek weak higher categories and coherent higher functors. The source inequality on site 42 is $i\geq n+1$, while the witness reverses it, thereby changing which homotopy groups are discarded. Site 43 defines the invariants through equivalence classes of iterated endomorphisms; the witness omits the full clause separating the set of components from the next endomorphism-class construction, and also misbrackets the displayed functor $C_n(K.)\to C_n(K'.)$. The capital $N$ in one sentence and the omitted basepoint in $\pi_i(C_n(K.))$ remain literal source features; the retraced subscript in the earlier definition is genuinely doubtful. Site 44 considers localization, truncation, and adjoining identity arrows, then passes toward infinity and relates strict Picard objects to the Dold--Puppe picture. His suggestion that localization might additionally require surjectivity on $\pi_{n+1}$ is a hesitation, not a needed condition for the category of $n$-types: for nonzero $A$, the map from a point to $K(A,n+1)$ is already an isomorphism on all groups through degree $n$, although it is not surjective in degree $n+1$. The higher categories in this comparison must retain the groupoid qualification: ordinary categories localized only at categorical equivalences do not identify the contractible interval category with the terminal category. Site 45 tests the proposal on loop objects, mapping objects, homotopy fibers, Eilenberg--Mac Lane objects, and Postnikov decomposition, leaving suspension without a satisfactory categorical candidate. These constructions feed Appendix 4's relative version and the later Quillen-letter program. The proposed classifying topos of a higher groupoid is expressly reconsidered in the following section, so the optimistic Postnikov description here should be read together with that later retraction.

\subsection*{Appendix 4: Relativization and the classifying-topos difficulty (cote 134-2, sites 45--46)}
The absolute comparison is relativized over a topos $X$, pairing simplicial sheaves with higher stacks in groupoids. The full source specification of the simplicial-sheaf side is omitted in the witness. More seriously, the source compares the evident inverse-image operation for simplicial sheaves with the subtler inverse image for higher stacks, whereas the witness substitutes direct image at that point; the next sentence then reverses the perspective, comparing direct image of stacks with a derived direct image of simplicial sheaves. Grothendieck explicitly hesitates between $Lf_*$ and $Rf_*$, and the readers retain both symbols and the hesitation. The Postnikov analogy is linked to the decomposition of a stack in groupoids into its component sheaf and a gerbe over the induced topos. On site 46, the paragraph headed ``Remords'' retracts the preceding expectation that extending classifying topoi to higher stacks should be straightforward: he cannot construct the requisite object in higher degree, and reduces the difficulty to obtaining a functorial topos of homotopy type $K(\pi,n)$ for an abelian group $\pi$. The proposed nerve construction, compatibility of homotopy and cohomology invariants, and the possibility of needing higher topoi are presented as unresolved requirements. His stated low-dimensional exception and the surrounding inequalities remain as written. This correction of perspective explains why Appendix 5 treats the proposed theory as a coordinated program involving several kinds of objects, rather than a completed identification of them; it also anticipates the later use of simplicial and categorical models to construct classifying spaces and homotopy types.

\subsection*{Appendix 5: The proposed topological algebra (cote 134-2, sites 46--47)}
The closing section of the first letter gathers four families of objects into one proposed theory: topological spaces and topoi; simplicial sets and simplicial sheaves; higher categories, especially higher groupoids, and higher stacks; and complexes of abelian groups or abelian sheaves. Their relations and possible additional structures, rather than one favored presentation alone, define the program. This synthesis depends on the Picard examples, the Lefschetz test, the homotopy hypothesis, and the relative difficulties developed in Appendices 1--4. Grothendieck proposes ``algèbre topologique'' as a shorter name than noncommutative homological algebra, argues that the familiar collection of topological rings, fields and groups does not constitute a sufficiently unified doctrine to monopolize that name, and invokes Serre's usage of analytic geometry as an analogy for terminological reassignment. The witness repeats a phrase in the enumeration on site 46 and omits the handwritten word ``vocable''; the checked source retains the actual single enumeration and the recovered word. The closing initials on site 47 are also present. The unusual extra typed acute above the m of the final analytic-geometry phrase is preserved as a source feature, not converted into a standard spelling. The next letter, beginning Appendix 6, takes up the promised relationship between locally constant higher stacks and a topos's homotopy type rather than claiming that this proposed unified theory already exists.

\subsection*{14. Homotopical algebra and the limits of present topology (pp. 1--2)}
Grothendieck opens the continuation of the letter to Quillen by distinguishing the study of spaces of maps and their successive homotopies from the study of spaces of homeomorphisms, immersions, embeddings and fibrations. His terms ``wildness'', ``isotopy'', ``stratified structures'' and ``tubular neighbourhoods'' delimit problems for which the familiar category of topological spaces seems poorly adapted; the contrast with piecewise linear structures is that their maps are too sparse, notably because real-valued piecewise linear functions are not closed under multiplication. This develops the complaint recalled explicitly from the letter, p.~1, and supplies the contrast on which section~15 builds its more favorable account of localization and homotopy types. The interlinear revisions narrow the opening description to ``maps between spaces'' and replace ``numerical'' by $\mathbb R$-valued (p.~2). The witness differs in two mathematically relevant places: the typescript says ``pro-topoi'' on p.~1, where arXiv has ``pretopoi'', and arXiv retains ``numerical'' without applying its handwritten replacement on p.~2. These readings are established by the scan; the edition does not identify the two notions or silently repair either authorial formulation.

\subsection*{15. Models of homotopy types and the place of stacks (pp. 2--4)}
The argument passes from the failure of the Hauptvermutung as a literal assertion about triangulations to its much weaker homotopical analogue: appropriate localizations of topological spaces and semi-simplicial sets yield the same homotopy category. Grothendieck compares CW-spaces, Kan complexes and Serre's cubical complexes, then proposes stacks, previously called ``$\infty$-groupoïdes'', as a third description particularly close to the homotopy groups themselves. His analogy with chain complexes and the Dold--Puppe equivalence makes the intended role of noncommutative structure explicit; section~16 immediately generalizes the motivation from spaces to topoi. The distinction between describing and computing invariants matters: on p.~3 he twice strikes ``description'' and writes ``computation'', corrections which the arXiv witness does not apply. On p.~4 a marginal note dated 3.5 begins ``This is an error'' and retracts the claim that stack cohomology has no immediate description from the primitive structure, referring instead to the nerve and to forgetting structure. Broad pencil strokes across this note make its final cancellation status uncertain, and one verb remains doubtful; both uncertainties are preserved. The arXiv footnote is substantially longer than this marginal wording and must not be substituted for it as though it were a literal transcription.

\subsection*{16. Stacks over topoi and noncommutative cohomological algebra (pp. 4--6)}
The motivation for stacks is enlarged from section~15's models of homotopy types to coefficients for cohomology over an arbitrary topos. Starting with Giraud's $0$-, $1$- and $2$-dimensional theory, Grothendieck treats $n$-stacks as coefficients whose ``integration'' is evaluation at the final object, with the homotopy groups of the resulting stack supplying cohomology after a dimension shift. His letters to Breen and Mme Raynaud's work on Lefschetz comparisons provide the geometric motivation: higher noncommutative cohomology should unify homological and homotopical comparison statements. The passage to infinite-order stacks is presented as a qualitative enlargement of perspective, and the following section asks how much of the homotopy-theoretic side ordinary categories can already model. The scans matter for the formulas: p.~4 has $H^i$, not the witness's $K^i$; p.~5 has $R\Gamma F_*$ and complexes $F_*$, where the supplied arXiv comparison reports a missing $F$ and a $j$ in place of the star. The handwritten qualification ``[non commutative]'' on p.~5 must also remain visible. These are witness differences, not authorial mathematical corrections introduced by this edition.

\subsection*{16$'$. Categories and the ``impressive bunch'' of models (pp. 7--9)}
Grothendieck associates a homotopy type to a small category both through its presheaf topos and through its nerve, comparing the first construction with Verdier and Artin--Mazur pro-homotopy types. He recalls and rejects his former guess that passage to the enveloping groupoid preserves this type: Quillen had pointed out that ordinary categories realize arbitrary homotopy types. This corrects the initial groupoid intuition while continuing sections~15--16's search for models, and motivates the later study of the ``impressive bunch'' and of coherators. A second proposed shortcut is itself explicitly retracted: the margin of p.~7, dated 5.3, says that the asserted weak equivalence $K\longrightarrow NN'(K)$ is false and directs the reader to section~24, pp.~21--22. The note is shorter than its arXiv counterpart; its call number is (18), and the section has been renumbered $16'$ rather than left as $16$bis. Page~8 proposes cofibrations injective on objects and arrows with a further call (19), then contrasts general $n$-stacks with $n$-Gr-stacks and their truncated homotopy types. Its wording ``left derived functors'' for the inverse-limit cohomology construction should be treated as an authorial formulation requiring care, not silently changed to the usual right-derived terminology. The final paragraph on p.~9 introduces the dependence on a coherator and the universal-algebra question, despite professing little inclination to pursue it.

\subsection*{17. The Artin--Mazur cohomological criterion (p. 9)}
This short section supplements the unification program by recalling that weak equivalences can be detected through cohomology with twisted coefficients, including noncommutative $H^0$ and $H^1$ as well as the commutative higher groups. Grothendieck explains how this permits étale cohomological information to constrain homotopy types, while wondering whether its success discouraged the development of higher noncommutative cohomology. The reference back to section~15 is the author's own handwritten cross-reference and is retained even though the immediate unification discussion occupies section~16 in the present numbering. The section bridges that motivation and the concrete corrections of section~18. Its handwritten insertion reads ``(namely, unification)'', not the arXiv ``messy unification''; the final sentence deletes ``that'' and says ``their observation'', distinctions that are obscured by the witness. The section number itself is corrected from 18 to 17 on the scan.

\subsection*{18. Corrections to the coherator and higher coherence (pp. 10--11)}
The return to the letter (L) is explicitly corrective. Grothendieck withdraws the existence of arbitrary binary amalgamated sums in $\underline B_o$: the geometric intuition of gluing contractible cells along embedded subcells had concealed the presence of nonmonomorphic maps, including the codiagonal. He also revises the construction of $\underline C_{n+1}$ from $\underline C_n$ by insisting on the total set of new arrows, and rejects the assertion that intersections of maximal subcells must be subcells. The marginal drawing on p.~10 makes the latter failure concrete: three regions $f,f',f''$ meet so that $f\cap f''=\{s,t\}$. The scan therefore supplies a two-point intersection, not the quite different formula in the initial transcription. A separate correction explains why associativity beyond primary composition cannot be represented by a single homotopy arrow: the sources and targets are themselves only homotopic, so squares, cubes and higher coherence data intervene. Page~11 names the dual categories left and right coherators, then connects the proposed non-groupoidal stacks to Benabou's bicategories and their cellular, topological interpretation. This prepares section~19's realization and singular-stack adjunction. The source wording $L\to K$ and $M\to K$ beside $K\amalg_L M$ on p.~10 is retained as an apparent authorial slip; it does not give the usual span for that pushout. Calls (11) and (14) on p.~11 are present without the explanatory footnotes printed in arXiv.

\subsection*{19. A provisional itinerary and the sheaf temptation (pp. 12--13)}
The first two items of the itinerary are the definition of $\underline B_o$, its dual and the coherators, followed by a realization/singular-stack adjunction with topological spaces. Unlike arbitrary presheaves on a left coherator, Gr-stacks must carry the designated amalgamated sums to fiber products. The realization functor nevertheless extends to all presheaves by preserving colimits; its values on basic cells are prescribed, while extensions to the coherator require choices. The anticipated independence of those choices is therefore tied to an as-yet undeveloped notion of equivalence between Gr-stacks. On p.~13 Grothendieck asks whether the exactness conditions could instead be sheaf conditions for a topology, which would make Gr-stacks a topos, and proposes weak equivalences detected by realization, monomorphisms as cofibrations and a lifting-property definition of fibrations. These are programmatic proposals rather than established results. The ``quick glance'' suggesting a topology is retracted immediately in section~20; even the concluding sentence about the inclusion, revised to composite, functor is already struck on p.~13. The source has ``one reward'', not the witness's ``one ward'', and handwritten ``as cofibrations'', not ``for cofibrations''. Its call (13) contains no footnote text on this page.

\subsection*{20. Model categories, sites and a retracted canonical topology (pp. 14--15)}
The section opens by correcting section~19's ``quick glance'': the proposed topology on the coherator would make the component functors fiber functors and force componentwise colimits, which the author recognizes as false. The category of categories supplies a warning and a new candidate construction. Grothendieck considers families of functors whose nerves cover every finite composable string, proposes identifying them with universally effectively epimorphic families, and associates to a category $A$ the topos of sheaves on $(\mathrm{Cat})_{/A}$. The handwritten margin immediately marks the asserted equivalence as false and points to section~24; the digits are readable, while pencil hatching leaves the note's cancellation status uncertain. Page~15 broadens the proposal to a natural site structure on model categories, with the homotopy type of $X$ identified with that of the induced site or topos over $X$. Its further presumption that the topology should be canonical is itself contradicted by a dated margin note absent from the arXiv comparison. These local corrections are essential to reading the section as an evolving investigation. The four arrows in the diagram of effective, universal and ordinary epimorphisms are handwritten, as are the topos tildes in $(M_{/X})^{\sim}=M^{\sim}_{/X}$; the object labels are typed. The inquiry leads into section~21's wider inventory of models rather than establishing a general theorem about canonical sites.

\subsection*{21. Multicomplexes, symmetry and restricted coherence (pp. 15--16)}
The inventory now includes symmetric variants of semi-simplicial and cubical complexes and their multiple-complex analogues. A diagonal complex already describes the intended homotopy type, but Grothendieck explicitly pursues the variety of algebraic presentations, not a unique most convenient model. He compares the passage from complexes to multicomplexes with that from ordinary categories to higher categories or stacks, recalling Quillen's suggestion of defining higher categories through multiple semi-simplicial complexes. The return to the coherator is concrete: the coherence diagrams examined in Benabou's bicategory axioms use restricted amalgamated sums describable through multisimplices, suggesting a smaller starting category $\underline B_o$ than the generous one in the original letter. The handwritten cross-references on p.~16 revise the old page references to sections~8 and~9, and the scan preserves the tense and spelling ``This I ckecked'', where arXiv reads ``This I check''. The suggested nerve-like functor from $n$-stacks to $n$-semi-simplicial complexes is a goal for later investigation. Section~22 next reverses direction, seeking simpler special cases of categories, beginning with preorders and ordered sets. The uncertain final singular/plural in ``connection(s)'' on p.~15 remains an image-edge problem and has no mathematical consequence for this comparison.

\subsection*{22. Ordered sets, realizations and maquettes (pp. 16--19)}
Preorders and ordered sets offer simpler categorical presentations of homotopy types. He relates ordered strata in conically stratified spaces to a realization whose barycentric simplices are finite totally ordered subsets, or flags. The resulting realization is canonically compared with that of the nerve. On p.~17 the corrected chain is $(\mathrm{Ord})\hookrightarrow(\mathrm{Preord})\hookrightarrow(\mathrm{Cat})\hookrightarrow(\mathrm{Ss\ sets})$, not the reversed first pair in arXiv. He cautions that the adjoint passage from a category to its object preorder can lose homotopy information, as the one-object case shows; the handwritten qualification ``but not Fl(A)'' is indispensable to that example. Page~18 gives a second realization through a nonseparated topology on the underlying ordered set and identifies sheaves with covariant functors, hence presheaves on the opposite order. Reversing the order preserves the flags and yields homeomorphic geometric realizations, while the two associated presheaf topoi need not be equivalent. The scan confirms hats in the latter claim and the author's rounded inclusion sign for the initial order comparison. Finally, maquettes are described by vertices and finite simplices and embedded in ordered sets through their face posets. The body claims that this category lacks finite products, while a dated marginal note adds that the notion is not an algebraic structure species; the hatching and call number of that note remain uncertain. Section~23 consequently shifts attention from satisfying Quillen's full package to constructing and localizing a natural functor into $(\mathrm{Hot})$. The edition preserves these claims and reservations as stages of the investigation, without treating the later witness's explanatory footnotes as text written on these pages.

\subsection*{23. The functor to homotopy types and a cellular core (pp. 19--21)}
Grothendieck separates the initial construction of a natural functor $M\to(\mathrm{Hot})$ from the further demand that $M$ carry a Quillen model structure. The functor would define weak equivalences by inverse image of isomorphisms; only then would one choose monomorphisms as cofibrations and test the resulting lifting-property definition of fibrations. The proposed site-theoretic construction uses the homotopy type of the induced site over each object, provided its pro-homotopy type is essentially constant. The canonical-topology suggestion is already qualified by a handwritten retraction pointing to section~24. Page~20 adds a structural requirement, universal and disjoint coproducts, and presents an algebraic species by a category $\underline C$ with finite inverse limits, its dual $\underline B$, and the anticipated equivalence $\mathrm{Ind}(\underline B)\simeq M$. This connects localization to finite presentations and to the question of which exactness properties are already visible in $\underline B$. Page~21 then isolates a smaller category $\underline B_{\mathrm{oo}}$ of basic cells and monomorphic boundary maps: the ordered set of proper subobjects should realize the boundary sphere of the corresponding cell. That construction uses section~22's ordered-set realizations and develops section~18's correction of the initial gluing intuition. The distinction between the broad species-level proposal and the geometric cell data remains heuristic. In the scan, the long p.~20 insertion includes ``can be interpreted'', omitted by arXiv, and the identity-functor arrow is drawn with two shafts; p.~21 preserves the printed slip ``of cubical'' and the typed replacements ``define'', ``family'' and ``sphere''.

\subsection*{24. Truncated nerve topologies and the failure of uniqueness (pp. 21--24)}
This section supplies the correction repeatedly anticipated in the earlier margins. Covering all composable pairs does not cover arbitrary strings: Grothendieck defines a descending sequence $T_N$ by epimorphic families of $N$-truncated nerves, identifying the first two levels with universal epimorphy and universal effective epimorphy. The comparison lemma relates the associated sheaf topoi to presheaves on truncated simplex categories. The finite levels retain only skeletal information, so the full nerve topology $T_\infty$ is needed for the proposed homotopy-type construction. The claim on p.~22 that an $N$-skeleton has the same homotopy and cohomology invariants through dimension $N$ requires caution at the endpoint: the $N$-skeleton of an $(N+1)$-simplex is an $N$-sphere, although the full simplex is contractible. The edition preserves the author's inequality rather than silently strengthening or repairing this intermediate assertion. The same loss of information explains the failure of $K\to NN'(K)$ discussed in section~16$'$: $N'$ factors through the two-dimensional data, while the nerve itself can still induce the desired equivalence after localization. Page~23 then replaces simplicial diagrams by commutative cubical diagrams, producing a strictly coarser topology with a different associated topos but the same intended homotopy interpretation. Multiple complexes suggest infinitely many further suitable topologies; the conclusion is that no unique topology is privileged. This leads directly to section~25's search for an equivalence relation among them. The scan's slice notation on p.~22 places $/A$ outside the parentheses ending after $T_N$, followed by an additional handwritten closing parenthesis; the earlier transcription's placement inside the subscript was corrected, while the extra parenthesis remains as written.

\subsection*{25. Hot-equivalence of topologies and global asphericity (pp. 24--25)}
The lack of a privileged suitable topology leads to a proposed equivalence relation. For comparable topologies $T\geq T'$ on $M$, Grothendieck compares the induced morphism of sheaf topoi and asks that every induced morphism over an object $A$ preserve homotopy type. He calls the resulting property ``globally aspheric'', strengthening the familiar cohomological notion of global acyclicity by allowing noncommutative coefficients. The primitive relation between comparable topologies is not yet an equivalence relation; taking the relation it generates is a first proposal, and he asks whether a coarser useful notion could do better. The intended invariant is a canonically associated functor $M\to(\mathrm{Hot})$, perhaps strong enough to recover the topology class itself. The analogy with Giraud's recovery of a topology from its associated subtopos motivates, rather than proves, a hoped-for correspondence. Page~25 narrows the aspiration to a single natural Hot-equivalence class for suitable species, admitting that even $M=(\mathrm{Cat})$ has no intrinsic characterization yet. Section~26 then reformulates the problem through small subcategories. Two revisions on p.~24 are essential: ``inverse image'' becomes ``direct image'', and the prime is moved so that the inclusion is from $T$-sheaves into $T'$-sheaves. The vertical margin note actually says ``just means'' and ``namely contractible'' with the qualification ``in the Čech sense''; the initial transcription's ``weakly contractible'' was not a literal reading. Its pencil shading leaves cancellation status uncertain. The body ends with the typed slip ``he'' for ``be'', and its final ``this'' is visibly canceled before the continuation on p.~25.

\subsection*{26. Test categories, test functors and the passage back to $M$ (pp. 25--29)}
The proposed intrinsic characterization of suitable topologies is reformulated through small categories $B$ inside a model category $M$. The cellular cores of section~23 and the simplex and cube categories of section~24 suggest that the presheaf category $\widehat B$ should itself describe homotopy types: its natural functor to $(\mathrm{Hot})$, computed from categories of objects over a presheaf, should be a localization. Grothendieck separates this internal requirement on $B$ from the eligibility of a functor $B\to M$; marginal terminology on p.~27 calls them ``test categories'' and ``test functors''. The comparison lemma gives an equivalence between the induced sheaf topos and $\widehat B$ under either of two fullness hypotheses. The decisive qualification comes on p.~28: localization of that sheaf category is not yet localization of $M$ itself. Generation by strict epimorphisms is considered as a way of making $M\to\widehat B$ fully faithful, but the author then proposes dropping both fullness requirements and directly demanding that $M\to\widetilde M\to\widehat B\to(\mathrm{Hot})$ be a localization. For categories $\underline B$ of finite presentations he further hopes for a localization onto homotopy types of finite type. The questions whether equal weak equivalences determine the functor, and whether suitable choices form one natural class, feed sections~27--28. The margin of p.~27 actually says ``same question arises'' for Hot-equivalence of topologies; ``arises'' was absent from the initial transcription. Its wording is now readable, but pencil hatching still leaves the note's cancellation status uncertain. The call after objects ``over a'' on p.~26 is (23), not (3).

\subsection*{27. Geometric realization as a third construction (p. 29)}
Calling the preceding discussion ``thin air conjecturing'', Grothendieck warns that both test categories and test functors may need stronger conditions. He places geometric or ``topologic'' realization beside the two constructions through topologies and through test functors. Semi-simplicial and cubical complexes offer comparatively direct realizations, hence so do the categories embedded in semi-simplicial sets. For Gr-stacks, the prescribed realization of the primitive $\underline B_{\mathrm{oo}}$ structure has to extend to a left coherator, and that extension involves the arbitrary choices already encountered in sections~18--19. Rather than immediately eliminating them, he compares this indeterminacy with the variety of test topologies and functors developed in sections~24--26. The paragraph thus states the next intended problem, relating realization to the two test constructions; section~28 immediately postpones it in favor of rigidity of $(\mathrm{Hot})$. The scan preserves ``topologic'', where the witness expands it to ``topological'', and the word ``is'' in ``is that a suitable'', where the witness has ``it''. Its call (28) supplies no text of the witness's retrospective footnote. The typed insertions ``localization'', ``left'' and ``between'' sharpen the objects of the comparison; the two round letters in $\underline B_{\mathrm{oo}}$ and the primitive graph designation remain literal rather than silently replaced by a joined infinity glyph.

\subsection*{28. Rigidity of $(\mathrm{Hot})$ and the first modelizers (pp. 30--33)}
The intended return to geometric realization is suspended by an ``Assumption'': the category of self-equivalences of $(\mathrm{Hot})$ and their natural isomorphisms should be equivalent to a point. Grothendieck separates its two demands, every self-equivalence is isomorphic to the identity and the identity has no nontrivial automorphism, while wondering whether arbitrary natural endomorphisms could also be admitted. This is asserted as a plausible assumption, not proved; it supplies the desired uniqueness up to unique isomorphism of a localization functor once its weak equivalences are known. A ``modelizing category'' or ``modelizer'' is then a category $M$ with a saturated class $W$ whose localization is equivalent to $(\mathrm{Hot})$. A model-preserving functor carries weak equivalences to weak equivalences and induces an equivalence of localizations, equivalently commuting up to isomorphism with the type functors. This formalizes sections~26--27's questions without yet specifying cofibrations and fibrations. The diagram on p.~31 compares ordered sets, preorders, categories, semi-simplicial sets and cubical sets; the category-of-elements functors $i_A$ lead to right adjoints $j_A$. Page~32 carefully distinguishes those right adjoints from the familiar nerve functors, contrasting infinite categories of simplices with finite ordered face structures when degeneracies are omitted. Mixed multicomplexes then enlarge the diagram and motivate the test-category definition in section~29. The mixed-product display on p.~32 visibly has $\Delta$ in both factors, although the adjacent prose says $p$ times simplicial and $q$ times cubical; this apparent authorial symbol slip is preserved. The same page changes ``modelators'' to ``modelizers'' by hand, a correction the witness does not apply, and contains a small branching diagram omitted there. The p.~30 cross-reference is to section~26; calls (28$'$), (29) and (30) are present without the witness's expanded footnotes.

\subsection*{29. The basic modelizer and a stronger test condition (pp. 33--36)}
The category of small categories now becomes the foundational modelizer: $(\mathrm{Hot})$ is its localization, and weak equivalences are described through the induced morphisms of presheaf topoi and the Artin--Mazur cohomological criterion recalled from section~17. A small category $A$ is provisionally a test category when $i_A\colon\widehat A\to(\mathrm{Cat})$, sending a presheaf to its category of elements, is model preserving; its presheaf category is an ``elementary modelizer''. This replaces a special preference for semi-simplicial sets by a comparison among all such presentations, while acknowledging the computational convenience of the nerve and the simplicial construction of pro-homotopy types. The general adjunction calculation of pp.~34--35 asks when both adjoints induce inverse equivalences after localization. Applied to $i_A\dashv j_A$, it leads to condition (T): every counit $i_Aj_A(B)\to B$ should be a weak equivalence. Grothendieck explicitly leaves open whether this strengthens the preceding definition, and proposes strengthening it if necessary; the call (33) does not contain the witness's retrospective answer. Products of test categories and direct comparison of elementary models become the next questions. The closing heuristic on p.~36 would test (T) by a fibration with contractible fibers, either after taking nerves or directly in $(\mathrm{Cat})$; section~30 immediately attacks the model-category assumption behind the latter route. On p.~35 the source twice writes that $M'$ or $M$ is saturated, although the discussion concerns the arrow classes $W'$ and $W$; these apparent symbol slips remain unchanged. The typed ``elementary normalizer'' on that page likewise remains source text. The scan's baseline slash in $M/_{b}$ on p.~34 is preserved, and its reference reads ``section 28'', not the witness's page reference. The broad emphasis of the witness also differs from the selective underlining and marginal double bar of p.~33.

\subsection*{30. From the model-category obstruction to asphericity (pp. 36--38)}
The attempt to treat every monomorphism of categories as a cofibration now runs into the two elementary horns drawn on p.~37. Lifting against both would force the arrows of a fibrant category to have left and right inverses, leaving groupoids, whose homotopy types cannot exhaust those of arbitrary categories. The objection concerns that proposed choice of cofibrations; the source does not prove the nonexistence of every possible model structure on categories. Grothendieck nevertheless keeps the expectation that elementary modelizers should admit closed model structures, and turns instead to the asphericity criterion introduced in section~25. A morphism of topoi is required to remain a weak equivalence after every localization over the target; when the topoi come from categories, generators reduce the condition to the asphericity of each slice $C'_{/a}$. Applying this to the counit $i_Aj_A(B)\to B$ from section~29 gives three equivalent conditions: the strengthened test condition, asphericity of every counit, and asphericity of $i_Aj_A(B)$ whenever $B$ has a final element. That last formulation opens the explicit interval argument of section~31. The scan's four bracketed ``[small]'' restrictions on p.~38 are genuine hand additions, including two absent from the witness. Its caret under the $A$ in $\operatorname{Hom}(A_{/a},B)$ is retained in its actual position without asserting a mathematical interpretation. The apparent inline hat on $f$ reported by the initial comparison on p.~37 belongs instead to the display beneath it; the inline symbol has been corrected to plain $f$. The word ``closed'' on p.~37 and the replacement ``aspheric'' on p.~38 are interlinear typing, although their placing strokes are handwritten.

\subsection*{31. An interval proof of the test-category criterion (pp. 38--41)}
The search for a direct proof that the simplex category is a test category yields three sufficient conditions. T1 asks that $A$ be aspheric; T2 asks that $A_{/a\mathrm{x}b}$ be aspheric for any two objects, with their product interpreted in the presheaf topos when necessary; T3 supplies an aspheric interval object $I$ with disjoint final subobjects $e_o,e_1$. For $F=j_A(B)$ with $B$ admitting a final element, the aim is a homotopy from the identity to the constant section. Diagram (D) on p.~39 expresses the cylinder $F\mathrm{x}I$ and its two endpoint inclusions. Given $u:a\to I$, the inverse image of $e_o$ determines a sieve $C_u$ in $A_{/a}$; a functor $f$ is extended by its own values on $C_u$ and by the final object of $B$ on the complementary objects. The sieve property makes the remaining arrows uniquely definable. This constructs the homotopy without T1, T2 or the asphericity of $I$. Those hypotheses enter the second step: the slice criterion of section~30 reduces the asphericity of $F\mathrm{x}I\to F$ first to $a\mathrm{x}I$ and then to $a\mathrm{x}b$, precisely T2. T1 then converts the weak equivalence to the final presheaf into absolute asphericity. Page~40 separately justifies the retract step through the Artin--Mazur criterion rather than claiming it is formal for an arbitrary arrow class. The scan gives genuinely hooked tails to both arrows into the cylinder in (D), contrary to the initial transcription; all diagram arrows are handwritten and its object strings are typed. The p.~39 reference is ``cf théorème p 46'', while the marginal comparison involving T1 and T2 retains an uncertain short connector. The live punctuation after ``checked'' on p.~41 is a colon, and the source's ``asperic'' remains uncorrected.

\subsection*{32. A revised programme after the interval criterion (p. 42)}
Dated 15.3, the new planning page treats T1--T3 as convincing candidates rather than a finished characterization. Its first three tasks are to examine necessity, construct the familiar simplicial and cubical examples, and prove stability under products so that different elementary models can be compared. The fourth returns to section~30's unresolved closed-model question; the fifth restores the larger programme of modelizers attached to algebraic structure species and the possible uniqueness of their weak equivalences. The author explicitly notices that stacks have receded into the background. A sixth task broadens the inquiry to a satisfactory replacement for Verdier's provisional triangulated-category notion, including the exactness, generator and cogenerator questions for $(\mathrm{Hot})$. Section~33 begins carrying out the first task and introduces the distinction between test and strict test categories. The scan matters to the direction of the programme: item 5 says categories $M$, not the witness's $\widehat A$. Its call (9) is present without the witness's explanatory footnote. ``Gave to this notion'' and ``general'' are typed interlinear revisions, while ``ways of'' and the displaced ``of'' are handwritten. The cancelled ``topos'' in item 4 is legible beneath both machine and pen deletion; the cropped completion of ``which'' and the old word after ``constructing'' remain doubtful.

\subsection*{33. Strict test categories, the Lawvere interval and a qualified theorem (pp. 43--47)}
The provisional terminology first separates test categories from ``strict'' test categories satisfying T1--T3. T1 is necessary by the final-category case of the criterion on p.~38. T3 admits a canonical candidate: applying $j_A$ to the two-object ordered interval produces a presheaf whose sections are sieves, hence the ``Lawvere element'' representing subobjects. This ties the explicit interval construction of section~31 to an intrinsic object of the presheaf topos. A proposed concluding statement about an intrinsic T3$'$ is completely cancelled on p.~43. The remaining issue is T2. Pages~44--45 connect it with preservation of binary products by the localization functor: the comparison $A_{/F\mathrm{x}G}\to A_{/F}\mathrm{x}A_{/G}$ is a weak equivalence for all $F,G$ exactly when the representable-product condition holds, and the slice criterion in fact supplies asphericity. The blanket claim about finite direct and inverse limits in $(\mathrm{Hot})$ on p.~44 is stronger than the finite sums and products used here; its marginal (!35) must not be expanded into an unqualified theorem in the guide. The result assembled on p.~46 compares product-preserving localization, the localized adjunction $i_A\dashv j_A$, and T1--T3. Its later margin dated 9.4 is essential: condition (i) is said to be ``not strong enough'', and the left continuation proposes adding that $j_A=i_A^*$ is model preserving, with a reference to remark~3 on p.~174. The left note is readable but crossed by broad pencil hatching, so its cancellation status remains uncertain. The source says ``model preserving'', not the witness's ``equally modelizing''. The author also explains how to avoid assuming that $(\mathrm{Cat})\to(\mathrm{Hot})$ preserves binary products: formulate the condition directly for $i_A$ up to weak equivalence. Page~47 provisionally reserves ``test category'' for the strict case, with call (38), and characterizes elementary modelizers through presheaf-topos structure, localization and finite products; a wider modelizing-topos notion is deferred. The scan retains ``asphéric'' on p.~43, three rounded intersection caps rather than the transcription's wedges, typed interlinear additions to the theorem, and the actual oval marginal annotations. The two opposite bent arrows on p.~47 remain graphic annotations, without an invented reordering instruction.

\subsection*{34. Product-stable examples and the homotopy circularity (pp. 47--50)}
The first examples strengthen the sufficient hypotheses again. A final object gives T1$'$, and an actual object $I$ of $A$ with two disjoint sections gives T3$'$; if $A$ is also stable under binary products, T2 follows because each product is representable and its slice has a final object. This constructs test categories inside any suitable ambient category, even one that is not small, by selecting a small product-stable subcategory. The economical construction uses powers of an interval and only projections and constant endpoint maps as components. Its cubical category includes permutations and diagonals, so p.~49 distinguishes it from the more familiar semicubical category. Nonempty finite sets give a corresponding full simplicial construction, with affine simplices as its geometric interpretation. The notation on p.~48 literally says to include $I$ and $\Delta_o,\Delta_1$; these triangular symbols are retained even though the preceding discussion specifies endpoint maps $\delta_o,\delta_1$. Returning to the ordinary simplex and cube categories removes the convenient product-closure hypothesis. Grothendieck therefore tries to prove T2 by making every representable deform to the final presheaf. The decisive revision on p.~50 is his own detection of a circular argument: section~31's homotopy proof used $F\mathrm{x}I\to F$ as a weak equivalence, and that was established from T2. A separate condition, called (2), that every $a\mathrm{x}I$ be aspheric must therefore accompany the deformation-retract condition TH1. Section~35 postpones this unfinished verification for the intrinsic theory of modelizing topoi. The formula $u_a=p_a c_a$ on p.~50 is retained in that visible order; with $c_a:e\to a$ and $p_a:a\to e$, the surrounding right-to-left composition convention would instead require $c_a p_a$ for an endomorphism of $a$. The scan also confirms the literal spellings ``axists'', ``oddurs'' and ``consequance''. The extra presheaf caret on the full simplicial category on p.~49 is present despite the witness omission, and the stepped mark on p.~48 is retained graphically in the diplomatic reading and applied as a paragraph break only in the master reading.

\subsection*{35. Local and total asphericity of topoi (pp. 50--55)}
The detour from section~34's unfinished simplex proof translates its conditions into the language of topoi. An aspheric topos is required to be $0$-connected and to have trivial cohomology with constant coefficients in every defined positive degree, including the noncommutative first degree. For an object or arrow, the induced topos supplies the corresponding notion. A locally aspheric topos has a generating family of aspheric objects; these generators suffice for testing whether an arrow remains a weak equivalence after every localization (p.~51). Proposition~1 identifies product stability of aspheric objects with their being aspheric over the final object. Total asphericity combines that property with local asphericity, and the interlinear ``non empty'' inserted on p.~52 is essential: without it, the preceding assertion that T2 already implies T1 would admit the empty case. The ensuing \v{C}ech argument gives absolute asphericity. For topological spaces the much stronger total condition becomes irreducibility, so a Hausdorff example must be a point. This sharpens section~33's distinction between a local supply of cells and the global product condition. Proposition~2 then separates the interval requirement from total asphericity. Its original ``totally aspheric'' hypothesis is visibly cancelled, and p.~53 explicitly says that no local-asphericity assumption is being made. The remaining equivalence is between the Lawvere element being aspheric over the final object and the existence of an object with that property and two disjoint sections. The homotopy argument of section~31 supplies the converse, and localization upgrades the resulting weak equivalence to asphericity. A connected topological space provides a useful obstruction: the proof with direct summands of its Lawvere element shows that this element is disconnected precisely in the irreducible case, where it has two connected components. Consequently no topological space can simultaneously meet total asphericity and the interval requirement (pp.~53--54). The proposed term ``modelizing topos'' packages both, while a second, still conjectural characterization asks for a saturated localization equivalent to $(\mathrm{Hot})$, compatible with finite sums and products and with the Verdier--Artin--Mazur construction. The closing inquiry about morphisms $X\to\operatorname{Top}(B)$ seeks a categorical description closer to the given indexing category than a large inverse system of semi-simplicial objects (pp.~54--55). The source retains ``0-connexity'', ``independant'', ``Corollay'' and ``monomorphisme''. Its upper margin on p.~54 reads ``cf (34)'' beside a body call (39); the lower strictness-related note remains partly doubtful. These are retained source features, not expanded into the arXiv witness's retrospective commentary.

\subsection*{36. Shuffles, nerves and comparison with categories (pp. 55--58)}
The simplex proof postponed in section~34 is completed by separating its two tasks. The ordered map $\Delta_n\mathrm{x}\Delta_1\to\Delta_n$ contracts toward the last element, giving TH1. For TH2, the maximal flags of the product order---the ``shuffles''---are attached successively so that each new intersection is a subflag. A Mayer--Vietoris-type lemma then propagates asphericity from two subobjects and their intersection to their union (p.~55). The author's inserted ``and covering e'' is indispensable to that lemma. He immediately warns against replacing this argument by the observation that the product ordered category has a final element: that shortcut assumes the still-unproved compatibility between the category's homotopy type and that of its nerve. The literal typed replacement ``apshericity'' remains uncorrected in the reader. The adjacent handwritten call is (41), not the asterisk proposed by the first transcription. The comparison proposition of pp.~56--57 addresses the missing compatibility in a more general form. A functor $a\mapsto\Delta_a$ from $A$ to categories defines $\alpha(C)(a)=\operatorname{Hom}(\Delta_a,C)$. If each $\Delta_a$ has a final element, $A$ is aspheric, and $\alpha(\Delta_1)$ is aspheric over the final presheaf, choosing final elements defines an aspheric functor $\varphi:A_{/C}\to C$, natural in $C$. The choice of final rather than initial elements makes the extension to arrows possible. Its slices identify with the corresponding construction over $C_{/x}$, so it suffices to treat $C$ with a final element; applying $\alpha$ to the ordinary categorical contraction supplies the required interval homotopy. Thus $i_A\alpha$ preserves weak equivalences and induces the identity on $(\mathrm{Hot})$. The simplex case now needs only the verified TH2, while the category of nonordered simplices is sent to the ordered set of nonempty faces of each finite set. A proposed sieve criterion for the interval condition follows on p.~58, but the entire corollary and its proof are diagonally cancelled before the new date 21.3. Its uncancelled presentation in arXiv must not conceal that revision. The source diagram on p.~56 literally labels its upper objects $a',a$ rather than $\Delta_{a'},\Delta_a$; the reader preserves this abbreviation or slip. The cancelled formula on p.~58 has $C_u\cap C_o$, not the witness's composition sign, and its first $F\to e$ occurrence is actually drawn with a mapsto tail. Section~37 turns the repeated argument into a new systematic account.

\subsection*{37. Homotopy intervals and the key comparison theorem (pp. 58--62)}
The new account dated 21.3 begins by identifying the gap in section~36's cancelled sieve criterion: the choices $C_u$ must be functorial in $u$, so that they actually define a map to the proposed interval. The replacement of ``automatic'' by a claim restricted to the intended examples makes that qualification explicit (p.~59). Grothendieck now formulates the argument in a category with finite products and a class $W$ satisfying isomorphism, two-out-of-three and retraction conditions. The last condition is stronger than the earlier mild saturation convention. An interval has two disjoint sections and an aspheric projection, with the required base changes understood; the margin calls the terminology ``provisional'' and adds ``separating''. The first homotopy lemma contracts an object to a section and concludes universal asphericity, while explicitly observing that disjointness was not used. The comparison lemma then transfers this property along an endpoint-preserving map to an object equipped with a unit-and-zero multiplication (p.~60). The revised formulas consistently make endpoint 1 the unit and endpoint 0 the zero, although the earlier ``respectively'' sentence is not fully adjusted. The Lawvere element supplies the universal example. In its proof on p.~61, the pullback along $e_0\to L$ is still called a subobject of $e_1$; the edition preserves this apparent index slip. The resulting theorem separates two strengths of comparison for a functor $i:A\to(\mathrm{Cat})$ whose values have chosen final objects. The choices need not be preserved by the arrows of $A$. Part (a) equates asphericity of the evaluation $A_{/C}\to C$, its being a weak equivalence, and the resulting identity on localized homotopy types, with equivalent tests on categories having a final element. Part (b) asks for asphericity relative to $A\times C$, equivalently for the appropriate representable products, and reduces this stronger requirement to universal asphericity of $i^*(\Delta_1)$. Together with asphericity of $A$, part (b) implies part (a); the first remark explicitly doubts the converse (pp.~61--62). The proof again transports the canonical categorical contraction through the left-exact functor $i^*$, while nonemptiness of each $i(a)$ ensures separation of the two endpoints. A superseded proposed condition in part (b) is cancelled by both machine and pen. The marginal reference to the sequel in volume~2 has a still-uncertain first page numeral, tentatively 64--65. Section~38 extracts the consequences for the Lawvere criterion and asks how far the final-object hypothesis can be removed.

\subsection*{38. The Lawvere criterion and generalized nerves (pp. 63--66)}
The theorem's main gain is the reduction to one interval condition: universal asphericity of $i^*(\Delta_1)$, together with asphericity of $A$, makes $i_Ai^*$ induce the identity on $(\mathrm{Hot})$. Both that interval condition and the localized conclusion make sense without chosen final objects in the categories $i(a)$, and p.~63 therefore anticipates a more general comparison theorem, while acknowledging that a direct evaluation map would then be unavailable. For the canonical choice $a\mapsto A_{/a}$, the right adjoint is $j_A$ and $j_A(\Delta_1)$ is the Lawvere element $L_A$. Corollary~1 consequently gives a sufficient criterion for a test category in the wide sense: $L_A$ is aspheric over the final presheaf and $A$ is aspheric (p.~64). The scan actually has a separate handwritten caret in ``unit element \ldots\ of $\widehat A$'' on p.~63; the initial transcription's bare $A$ was not a source slip. For general $i$, a crible in $i(a)$ determines the crible of arrows $b\to a$ whose images factor through it. This defines an endpoint-preserving map $i^*(\Delta_1)\to L_A$, compatible with intersection. Section~37's comparison lemma transfers the interval condition to $L_A$. Corollary~2 then makes both $i^*$ and $i_A$ modelizing and identifies the functor induced by $i^*$ with the canonical quasi-inverse induced by $j_A$ (pp.~64--65). Thus the localized functor does not depend on the admissible choice of $i$, although the resulting concrete presheaf models can be much simpler than those obtained from the often infinite slice categories $A_{/a}$. The ordinary nerve is the motivating example of this ``generalized nerve functor''. Corollary~3 makes the remaining verification practical: if $A$ has a final object $e$, $i(e)=\Delta_0$, and a homotopy interval $I$ admits an endpoint-preserving map $i_!(I)\to\Delta_1$, adjunction and the same comparison lemma give the desired asphericity of $i^*(\Delta_1)$ (p.~66). In the intended examples $I$ belongs to $A$ and $i(I)=\Delta_1$, so only asphericity of all products $I\mathrm{x}a$ remains to check. The reference to the cancelled corollary is revised to the end of section~36. Page~64 literally says ``for a in J'' where the construction indexes by $a\in A$, and retains a stray typed $u$ beneath a final-object caret; these visible irregularities remain in the readers. Section~39 returns to total asphericity and refines the terminology accordingly.

\subsection*{39. Strict, local and wide-sense test categories (pp. 66--69)}
The return to total asphericity first reuses the homotopy lemma: if every representable admits a contraction to a chosen section through a fixed homotopy interval, then every representable is aspheric over the final object and $\widehat A$ is totally aspheric (p.~67). The proof does not use separation of the endpoints. The following assertion that separation follows from the other assumptions is too broad if those assumptions are taken alone: the terminal category with $I=e$ supplies a counterexample. The boxed and diagonally cancelled paragraph had discussed this exceptional case; its removal is preserved, together with the still-live preceding assertion. This is one of the places where the revision marks matter to the logical reading. The terminology is now reorganized around the Lawvere element. A local test category has $L_A$ aspheric over the final presheaf; a test category is additionally aspheric; a strict test category additionally has a totally aspheric presheaf topos. Their presheaf categories are correspondingly local elementary modelizers, elementary modelizers and elementary strict modelizers. A general strict modelizer is characterized by its weak-equivalence class and a localization equivalent to $(\mathrm{Hot})$ that preserves finite products; further requirements concerning sums and infinite products are still being considered (pp.~67--68). The weaker adjunction criterion $i_Aj_A(C)\to C$ remains the definition of a test category ``in the wide sense''. The source's dismissal of the likely usefulness of that condition is an uncancelled parenthetical, while an adjacent faint marginal note is more favourable; its exact wording and final Delta superscript remain uncertain. The page's note call is 45. The later shorthand ``T1 + T3'' must be read with care: the earlier T3 on p.~39 asked only for an aspheric object with disjoint endpoints, and its margin used T2 to pass to asphericity over the final object. The operative interval condition in the present definition is explicitly relative. The terminology for general modelizing topoi is proposed with an express reservation that its connection to actual models has not yet been established (p.~68). The distinction between test and strict test categories is then made concrete. A first proposed small category is generated by interval powers together with $a_o\mathrm{x}I^n$; Grothendieck expects the missing product $a_o\mathrm{x}a_o$ in the presheaf category to fail asphericity. A more decisive construction takes any test category $A$ and a homotopy interval $I$: the slice $A_{/I}$ is again local test and aspheric, hence a test category, but is never strict because the two endpoint subobjects of its final object are nonempty and disjoint (p.~69). Total asphericity would force their intersection to be aspheric, so it cannot hold. The first construction remains an expectation, whereas the slice construction supplies the stated universal counterexample. The scan's opening insertion on p.~69 reads ``as'', giving ``and as maps \ldots\ all those'', rather than the witness-based ``those'' of the initial transcription. Section~40 next supplies the product-stability lemma that the preceding arguments have been using.

\subsection*{40. Products and relative weak equivalences (pp. 69--71)}
The product argument supplies a fact already used in section~37: if $A$ is aspheric, the projection $C\mathrm{x}A\to C$ is aspheric for every $C$. The slice criterion reduces first to $C$ with a final element and then to products $C\mathrm{x}A_{/a}$, which themselves have final elements (p.~69). Grothendieck strengthens this to stability of weak equivalences under cartesian products. It is enough to treat $f\mathrm{x}\mathrm{id}_C$ with $f:A'\to A$ a weak equivalence. Viewing this as a morphism over $C$ motivates weak equivalence relative to a base topos $Y$: the property must survive each localization by an object of $Y$, and testing on generators suffices. For categories over $C$, the test is the induced functor over every $c\in C$. Thus the product statement reduces to $C$ with a final element, where the projection square and two-out-of-three prove it (p.~70). The displayed localization arrow literally has target $U$ in $Y_{/U}\to U$; its described role requires target $Y$, so the reader preserves the visible letter and flags the mismatch here. The final square's upper handwritten label has lowercase $c$, another source detail retained as drawn. Universal weak equivalences strengthen this relative notion by requiring stability under every base change in $(\mathrm{Cat})$, not only localizations (p.~71). Their test on base-change categories with a final element suggests an analogy with trivial fibrations. Grothendieck proposes that fibering or cofibering functors with aspheric fibers, and more generally such functors whose transition functors are weak equivalences, might supply the categorical counterparts of Quillen's fibrations. The inserted qualification ``in the purely categorical sense'' limits the proposal. Cofibrations remain unidentified, but monomorphisms alone appear insufficient, prompting the revised references to section~30 and (9). These are prospective comparisons rather than an established model structure on $(\mathrm{Cat})$. The margin labels the terminology ``awkward'' and carries call (46). The section then deliberately returns to test categories, leaving section~41 to clarify exactly which assumptions the preceding proofs have used.

\subsection*{41. The guiding assumption and weakened saturation (pp. 71--72)}
The conjectural triviality of the category of autoequivalences of $(\mathrm{Hot})$ had encouraged defining a modelizer through its weak-equivalence class, with the functor to homotopy types following essentially uniquely. Grothendieck now emphasizes, by changing the final period to an exclamation mark, that none of the proved statements has used that guiding assumption (p.~71). The operative definition of modelizer has also changed: the initial requirement of strong saturation, meaning that $W$ consists of precisely the arrows inverted by localization, has been replaced by the three weaker conditions written in section~37. The authorial transposition under ``wrote explicitly down'' supplies ``wrote down explicitly'' in the applied reading; the diplomatic text keeps the original order and the pen curve. The source spelling ``drapping'' is left visible where the earlier transcription had conjectured an $o$ correction. Strong saturation can nevertheless be transferred between $(\mathrm{Cat})$ and an elementary modelizer through the established comparison functors, so one example suffices for all (p.~72). Grothendieck invokes the known closed model structure on simplicial sets and his cited proposition~1 in Quillen's I~5.5 for that example. His remaining question concerns the proof of that model structure: can the present methods replace the specialized simplicial argument by one applicable to general test categories, or at least strict ones? He explicitly postpones this question and returns to comparing test categories and completing the inventory of standard examples. The handwritten call following ``strongly saturated'' is (47), while the separate margin has an asterisk and a reference to underlined VIII whose short abbreviation remains uncertain. This transition motivates section~42's matching revision of the definition of a model-preserving functor.

\subsection*{42. Model-preserving functors and full subcategories (pp. 72--74)}
The weakened saturation convention of section~41 requires a corresponding change to the definition of model-preserving functor: in addition to inducing an equivalence after localization, $f:M\to M'$ must satisfy the equality $W=f^{-1}(W')$, so that it reflects as well as preserves the specified weak equivalences (p.~72). The revised text first returns to the unresolved modelizing status of ordered sets and preorders. It then distinguishes the equivalence of localized categories from a canonical comparison with the chosen category of homotopy types. The earlier rigidity assumption would supply such a comparison, but has not been established. A subsidiary question asks whether an endofunctor of $(\mathrm{Hot})$ that is the identity on objects must be naturally isomorphic to the identity (p.~73). These questions remain conjectural; they do not become consequences merely from the existence of the concrete comparison functors already constructed. The constructive result uses a full subcategory containing the models produced by a test category. In the preliminary abstract argument an endofunctor $h$ of a modelizer factors through $M'\subset M$ and induces an autoequivalence on the localization; the text then describes the induced inclusion and factor as quasi-inverse (pp.~73--74). That exact formulation is stronger than its stated hypothesis: their composite is the localized $h$, and an autoequivalence need not be isomorphic to the identity. For example, taking $M'=M$ and inclusion equal to the identity still leaves that distinction. In the intended application $h=i_Ai^*$, the key comparison theorem does provide the required natural comparison to the identity. The resulting proposition says that a full subcategory of $(\mathrm{Cat})$ containing all categories of elements $A_{/F}$ for a test category $A$ inherits a modelizer structure and has model-preserving inclusion (p.~74). Categories locally isomorphic to $A$ supply a useful reformulation of the required objects. This reopens the ordered-set problem: a suitable test category defined by an ordered set would put all its categories of elements inside $(\mathrm{Ord})$. Such a category could not be strict, by the earlier obstruction for topoi attached to topological spaces in section~35. The last remarks contrast localization stability for test categories with the stricter total-asphericity condition, and lead directly to the standard examples of section~43. The source retains its literal ``T1 + T3'' shorthand. Its interlinear ``inclusion'', ``modelizer'', ``the test category'', ``algebraic'' and ``a suitable'' are typed revisions; the cancelled locally-ordered qualification is a genuine pen deletion.

\subsection*{43. Ordered, unordered and face-only simplices (pp. 75--78)}
The return to standard examples applies section~38's test-functor criterion directly. For ordered simplices, the inclusion $i:\Delta\to(\mathrm{Cat})$ is fully faithful and carries the interval to the ordinary two-object ordered category. Thus $i^*(\Delta_1)$ is the representable interval, and the previously proved asphericity of its products with simplices gives the test condition. Contractions to final vertices then give total asphericity and hence strictness (p.~75). The cancelled preliminary proof is replaced by this more precise argument. For unordered simplices, a small skeleton of nonempty finite sets has finite products and a two-point interval. Sending a finite set to its ordered set of nonempty subsets supplies the test functor. On the two-point set this produces the span of the two singleton subsets into the full set; the endpoint-preserving map to $\Delta_1$ sends $\{0\}$ to 0 and both other subsets to 1 (pp.~75--76). The comparison square with ordered sets identifies the resulting functor on ordered simplices with the flag, or barycentric-subdivision, construction. The single untilded $\Delta$ in the unordered-category conclusion on p.~75 is a visible source omission; the later functor-domain notation has an additional uncertain wavy stroke, rather than a securely readable presheaf hat. Removing degeneracies gives a substantially different answer. Grothendieck first emphasizes that face operators already determine homology, cohomology and the fundamental groupoid, and hence the cohomological test for a weak equivalence of the associated complexes. Nevertheless, in the category $\Delta^f$ of strictly increasing simplex maps, the product of the zero-simplex with an $n$-simplex has a discrete category of elements with $n+1$ objects. It is therefore not aspheric when $n>0$ (p.~76). Removing the zero-simplex does not supply a representable interval with sections, since arbitrarily large simplices cannot map into a fixed smaller one. The more flexible attempt using the presheaf $i^*(\Delta_1)$ also fails the required relative interval test. The canonical Lawvere element gives the same two-point obstruction over the zero-simplex (p.~77). Thus $\Delta^f$ is not a local test category, while the possibility that it is a test category in the wide sense remains open at this point. The added name ``weak test category'' anticipates the next section's systematic terminology. The remaining weak condition asks for absolute asphericity of $i_Aj_A(C)$ for every category $C$ with a final object. Its first tests are asphericity of $A$ itself and of the Lawvere element, expressly distinguished from asphericity over the final presheaf (p.~77). A positive answer would let section~42's proposition construct further full-subcategory modelizers. To pursue the ordered-set case, Grothendieck considers the poset $J$ of nonempty finite subsets of an infinite vertex set. It has the desired local simplex structure and, for countable vertices, is a category of elements of an infinite face simplex. Its filtering property suggests absolute asphericity. The local test condition still fails because a minimal vertex has a one-point slice, which cannot be a test category (p.~78). The closing calculation identifies the category of elements of $L_J$ with pairs $K\subset T$, with $T$ finite and nonempty, ordered by $T'\subset T$ and $K'=K\cap T'$. It begins an explicit test of the surviving weak condition rather than asserting its conclusion. Section~44 opens the following day by announcing that the face-only simplex category does satisfy that weaker notion. The readers retain the source's shifted apostrophes, the cancelled product/intersection argument on p.~78, and the distinction between the two inline typed round-letter indices and the displayed handwritten infinity.

\subsection*{\S67. Aspheric maps replace contractibility as the organizing notion (pp.\ 188--191).}
Grothendieck reorients the modelizer program: the crucial objects are aspheric objects and weak equivalences, and its natural comparison maps are aspheric functors. He explicitly judges the earlier restriction to contractibility and special test functors to have been an unnecessarily long detour (p.\ 188), while retaining its possible relevance for ``canonical modelizers.'' The discussion now allows arbitrary small categories, rather than only test categories. For a basic localizer $\underline W$ he introduces $(\mathrm{Hot}^{W}_{A})$, also written $(\mathrm{Hot}_{A})$, as the localization of the presheaf category by $\underline W_A$; the construction is then organized by inverse image $i^*$ and the comparison morphism $\lambda_i(F):A_{/i^*F}\longrightarrow B_{/F}$, interpreted through the corresponding induced topoi (pp.\ 189--190). This prepares the promised description of asphericity structures and testing functors at the end of the section (p.\ 191); the next section interrupts that development. The dated marginal correction of 3 May on p.\ 189 expressly retracts the claim that Quillen's weak equivalences are precisely $\underline W_A$, saying that Quillen's are stronger. It is therefore essential to read the original assertion together with the margin. The witness changes handwritten ``or'' into equality in (1), drops the subscript $A$ on $\underline W_A$, omits the handwritten domain and codomain after the action on $(a,p)$ on p.\ 190, and prints the struck word ``real'' on p.\ 188. These differences concern notation, functorial context, or the author's revision; none licenses an unrecorded mathematical repair of the typescript.

\subsection*{\S68. Mathematical work, the Teichmüller tower, and the necessity of groupoids (pp.\ 191--195).}
Dated 27 May, the section explains the six-week interruption in the asphericity discussion through a retrospective of the author's changing mathematical interests. Its mathematical center is the ``Lego-Teichmüller construction game'': cutting and gluing surfaces, the tower of groupoids $T_{g,\nu}$, and the arithmetic action on their profinite completions (pp.\ 191--192, 194--195). The account of \emph{La Longue Marche}, tame topology and ``unfolding'' of stratified structures supplies the context for the anabelian project; the intervening reflection on creative interest, ambition and alternating periods of work belongs to the argument's historical presentation and must not be omitted as an aside (p.\ 193). The final page gives the technical lesson: reducing several natural base points to one can destroy visible symmetries and obscure a presentation by generators and relations, whereas groupoids preserve the geometry. Van Kampen situations supply the comparison, and the modular ``multiplicities'' are identified as topoi whose points have finite automorphism groups (p.\ 195). This is a return to the text's groupoid motivation, rather than a proof of the asphericity claims left pending in \S67; \S69 resumes a retrospective of the technical scratchwork. The section explicitly revises an earlier preference for groups or ``exterior groups'' in favor of groupoids (p.\ 195). The witness has several ordinary wording errors here---including ``case'' for ``care'' on p.\ 195---while the underlining and handwritten exclamation in ``topoi!'' preserve the emphasis of the typescript.

\subsection*{\S69. Derivators and integration of homotopy types (pp.\ 196--202).}
The technical retrospective explains why an ordinary derived or homotopy category retains too little structure: traces, determinants and the compatibility of iterated constructions motivate the system of localized diagram categories $D_A(I)$, contravariant in the indexing category $I$, with inverse-image functors expected to admit both adjoints (pp.\ 196--198). This develops the modelizer program recalled in \S67 and connects it with the geometric motivations of \S68. For categories, a covariant diagram can be integrated by forming its associated cofibered category; the comparison with localized slice categories yields the model for ``integration of homotopy types'' (pp.\ 198--199). The name ``derivator'' is introduced on p.\ 200, together with the internal stability expressed by $D(I\times J)$ and the aspiration to a corresponding homotopy theory of derivators. The geometric example is the system of tubes around strata of Mumford--Deligne compactifications: their $K(\pi,1)$ homotopy types, with Teichmüller-type groups, suggest an algebraic description from a diagram of fundamental groupoids (p.\ 201). A diagram in $(\mathrm{Hot})$ alone is explicitly said to be insufficient; a diagram in a model category is required. Grothendieck retracts his recent expectation that finite limits exist in $(\mathrm{Hot})$ (p.\ 201), and \S70 immediately examines the extra factorization condition proposed for the two adjoints (pp.\ 200, 202). Significant witness differences include $T$ replaced by $\mathscr D$ in the statement about inverse limits (p.\ 198), $q=f\circ p$ replaced by an equation beginning with $g$, and ``basic modelizer'' replaced by ``basic localizer'' (p.\ 199). These readings are clear on the images and are preserved rather than repaired from the witness.

\subsection*{\S70. Base change, duality, and the limitations of $(\mathrm{Cat})$ (pp.\ 202--206).}
Starting from the desired two-stage factorization $X\to Z\to Y$, Grothendieck asks for preservation of weak equivalences by base change along the second map and co-base change along the first, with either factor optionally a weak equivalence (p.\ 202). The goal supplies the extra condition proposed for derivator operations in \S69. The investigation of the base-change side leads to cohomological smoothness and properness of functors, characterized through asphericity; fibered categories are smooth, cofibered categories proper, and passage to opposite categories interchanges the two conditions (p.\ 203). The pairing of presheaves and copresheaves suggests a broader duality of topoi, but the claim that cosheaves on an arbitrary topos form a topos remains a question (pp.\ 203--204). The co-base-change side is less tractable: open or closed immersions in $(\mathrm{Cat})$ fail the hoped-for general assertion, whereas pushouts along monomorphisms in $A^{\wedge}$ have the required homotopy meaning through the integral of the associated span (pp.\ 204--205). The stated smooth/proper exceptions and the possible special case of an immersion that is also a weak equivalence must be distinguished from that rejected general assertion. This forces a return to elementary modelizers and Quillen factorization questions; the author finally accepts the Kan condition after describing his resistance to it, while retaining dissatisfaction with the reduction to semisimplicial methods (p.\ 206). These questions become part B of the program in \S71. The witness omits the monomorphism hook in the p.\ 204 square and replaces the actual objects $Y,X,Y'$ of the p.\ 205 integral diagram by categories of elements. The scan-confirmed diagrams preserve the distinction between these two levels of the construction.

\subsection*{\S71. A provisional program: modelizers, derivators, and topoi (pp.\ 206--212).}
The program dated 30 May organizes the preceding scratchwork into four tasks. Part A returns to the asphericity structures and canonical modelizers suspended in \S67; part B develops the base-change, co-base-change and factorization questions of \S70, with homotopy integrals, amalgamated sums and Mayer--Vietoris sequences as the connecting constructions (pp.\ 206--207). A handwritten transposition on p.\ 207 makes strong saturation one of the properties expected to follow from the first three basic-localizer axioms. Part C develops the derivator proposal of \S69 while asking what additional structure is needed for external Hom's and basic invariants. It then distinguishes group objects modeling pointed connected homotopy types from the truncated output of group objects in $(\mathrm{Cat})$, and asks which modelizers satisfy a Dold--Puppe theorem (pp.\ 207--209). The proposed trilogy consists of semisimplicial, cubical and hemispherical complexes, the last also called globular complexes by Brown; their boundary and degeneracy operations are related to regular cellular subdivisions of spheres. Multicomplexes give multiple chain complexes, so the simple Dold--Puppe statement does not extend unchanged. The classification of the three series and the need for enough degeneracy maps are presented as questions to investigate, not proved theorems (pp.\ 209--210). Part D, ``Back to topoi,'' seeks relative homotopy categories $D(T)$ and their possible opposite counterparts, extending the $f^*,f_!,f_*$ formalism from indexed categories to topoi. The discussion keeps existence of derived adjoints tentative and moves from objects to proobjects when a topos fails to determine an essentially constant homotopy type (pp.\ 210--211). Local asphericity then reconnects the program with \S35; the revised distinction between modelizing and strictly modelizing topoi is recalled from p.\ 68, and the section ends with the Lawvere-element condition for a finitely triangulated contractible space (p.\ 212). The next section begins the promised return to the asphericity game. The witness changes $f_*$ to $f^*$ and $Lf_!$ to a subscript-$1$ expression on p.\ 210, and applies a handwritten transposition there without marking it. The duplicated target-only display on p.\ 211 was an error of the draft transcription, absent from the scan; its removal restores the actual three-display sequence.

\subsection*{\S72. Criteria for aspheric functors and the first stability results (pp.\ 213--219).}
Returning to \S67 after the program of \S71, Grothendieck isolates the localizer assumptions needed for asphericity: mild saturation, the condition on categories with final object, and the locality condition Loc~3 (pp.\ 213--214). Proposition~1 identifies the comparison maps on categories of elements, their restrictions over objects, and the corresponding slice criteria; it then distinguishes these equivalent asphericity conditions from preservation of weak equivalences and equivalence after localization (pp.\ 214--216). The extra hypotheses in the implication diagram matter: the discrete-target example separates aspheric fibers from merely nonempty fibers, and the product-projection example shows that inducing the localized equivalence need not determine an aspheric map without the stated assumptions (p.\ 216). Under asphericity, the localized triangle commutes up to the canonical comparison isomorphism; this reconnects the discussion with the earlier ``inspiring assumption'' of \S28 without using that assumption as a proof (pp.\ 216--217). Proposition~2 transfers total asphericity along an aspheric functor and relates the inverse image of a separating homotopy interval to asphericity of the functor. Its handwritten addition makes the source category aspheric, and the long interlinear insertion specifies the homotopy condition (TH1) of p.\ 50; the corollary deduces strict test categories from a totally aspheric source and a local test target (pp.\ 217--218). Proposition~3 supplies composition and right cancellation, base change along a fibering map, and products (pp.\ 218--219). The base-change proof uses the fibration condition L5, whose derivation from the stated localizer assumptions is expressly postponed; Grothendieck explains that only part~(a) is needed for the immediate asphericity-structure formalism, preventing the announced application from relying on that deferred argument (p.\ 219). Remark~4 denies the opposite cancellation assertion and leads directly to the examples in \S73. The visual collation restores the proof-diagram arrow from (i'') to (iii), rather than to (i'), on p.\ 215; it also identifies the apparent deletion of the t in Hot on p.\ 214 as pen retracing. On p.\ 218 the comparison sign is a handwritten tilde over a bar, with no arrowhead, contrary to the draft transcription and its diff description. These are distinctions between what the image shows and what a mathematically plausible reconstruction might supply.

\subsection*{\S73. Adjoints, coasphericity, and cofinality (pp.\ 220--221).}
The continuation dated 11 June begins by suppressing the qualifying $\underline W$ when one basic localizer is understood. It sharpens Remark~4 of \S72: a functor from the one-point category into an aspheric category is aspheric exactly when its chosen object is initial, illustrating how much stronger a map condition is than asphericity of its source and target (p.\ 220). Proposition~4 enlarges the preceding isomorphism example: every equivalence is aspheric, and more generally a left adjoint is aspheric because its slice categories identify with slices having final objects. Passage to opposite categories then motivates the name ``coaspheric'' for the dual condition, satisfied by the right adjoint. The corresponding undercategories and their initial objects connect coasphericity with a promised all-dimensional cofinality statement, whose dimension-zero consequence is invariance of direct limits (pp.\ 220--221). The analogy from \S70 returns: asphericity resembles properness with aspheric fibers, and coasphericity resembles smoothness with aspheric fibers. The comparison is explicitly weakened because arbitrary base change does not preserve either condition; the applicable base changes are fibering/smooth or cofibering/proper, respectively. Even an equivalence of categories can have an empty fiber, so its simultaneous asphericity and coasphericity do not supply arbitrary base-change stability. Universal weak equivalences, called ``trivial Serre fibrations'' here, are proposed as a stronger common condition and assigned to a later systematic investigation (p.\ 221). The section then turns to the overdue product discussion of \S74. On p.\ 220 the crossed handwritten heading above Indeed is Remark~5, not the draft's tentative Proof; it is written again at the subsequent paragraph. Whether the placement strokes also delete its typed initial The remains unresolved and is marked as such. The p.\ 221 witness changes singular condition to plural conditions and supplies a title for \S74 absent from the typescript.

\subsection*{\S74. Products, diagonal functors, and canonical comparison of models (pp.\ 221--230).}
The overdue product discussion supplies concrete constructions for the abstract stability results of \S72. The external product of presheaves has a category of elements identified with the product of the two categories of elements; it preserves aspheric objects and weak equivalences in the stated variables, and yields the product assertions for locally aspheric and totally aspheric categories (pp.\ 221--222). The corollaries distinguish local test, test, and strict test categories: the hypothesis on the other factor changes from arbitrary to aspheric to totally aspheric. Mixed multicomplexes illustrate the construction (p.\ 223), reconnecting it with the trilogy in \S71. The comparison functor associated with an aspheric map between modelizing categories leads to a functor from $\mathrm{Asph}(A,B)$ into a category of comparison functors and factors through its fundamental groupoid. Nonemptiness and connectedness suffice for the existence of an isomorphism class of comparison, while a canonical transitive system requires the stronger simple-connectedness conclusion; the text carefully separates these claims (pp.\ 224--225). The diagrams on p.\ 226 distinguish categories, presheaf modelizers, and localized homotopy categories, including the middle presheaf triangle omitted from the witness.
The diagonal criterion then identifies total asphericity of $A$ with asphericity of $A\to A\times A$, and permits simultaneous comparison through products of strict test categories (p.\ 227). Proposition~4 treats pointwise products of functors into a category $B$ with binary products and aspheric presheaves $\underline{\mathrm{Hom}}(b,y)$: the total-asphericity assumption on $A$, moved by a handwritten transposition into the hypotheses, is essential to the stated construction (p.\ 228). The external-product corollary is sharpened by revision to an equivalence: when $i$ is aspheric, $i\boxtimes i'$ is aspheric if and only if $A'$ is aspheric. Its proof uses a lemma on a presheaf over $A\times A'$ whose restrictions to $A$ are aspheric. The resulting category-of-elements projection is a weak equivalence by the fibration condition L5, explicitly recalling the claim and deferred justification on p.\ 219; this dependency should remain visible (p.\ 229). Contractibility of the objects of $B$ is discussed as a stronger practical hypothesis. Finally, closure of $\mathrm{Asph}(A,B)$ under binary products makes its nonempty case connected with the needed homotopy strength, giving the canonical transitive system of isomorphisms between the localized functors and the canonical functor announced on p.\ 225 (p.\ 230). A handwritten margin points forward to the generalization in \S82, p.\ 272. The image-based edition retains the unusual $X\longrightarrow B_{/X}$ display on p.\ 224, the missing slash in the final category-of-elements notation on p.\ 226, and the visibly typed product x and lowercase subscript o on p.\ 228. The hemisphere-like symbol on pp.\ 223 and 229 remains a doubtful approximate glyph; neither mathematical expectation nor the witness settles its exact drawn form.

\subsection*{\S75. Asphericity structures, testing functors, and weak equivalences (pp.\ 230--237).}
The generalities of \S\S72--74 now become a structure on a possibly large $\underline U$-category $M$. After making the universe convention explicit, Grothendieck defines a class $M_{as}$ of aspheric objects through a small category $A$ and a functor $i:A\to M$ that lands in $M_{as}$ and detects it by the presheaf $i^*(x)=\mathrm{Hom}(i(-),x)$ (pp.\ 230--231). Such an $i$ is a “testing functor.” Every small full subcategory supplies an asphericity structure, and Proposition~1 plus its corollary shows that any structure can be generated this way; enlarging the generating full subcategory within $M_{as}$ preserves the structure (p.\ 232). The following terminology discussion distinguishes testing functors from the earlier test functors of pp.\ 175--176. The essential condition depends on the possibly large category $M_{as}$ itself, rather than on its embedding in $M$, so the author tentatively adopts $M_{as}$-aspheric, or simply aspheric, as a synonym. His use of both “aspherical” and “aspheric” in this discussion is preserved (p.\ 233).
Proposition~2 defines $W=(i^*)^{-1}(\underline W_A)$ and proves independence of the testing functor by passing to a common small full subcategory. Mild saturation is immediate; strong saturation remains conditional on the corresponding result for the basic localizer (p.\ 234). When $M$ has a final object, calling the structure aspheric means that this final object belongs to $M_{as}$; the resulting criterion recovers aspheric objects from the maps to the final object. The small-generator descriptions then characterize both the admissible weak-equivalence sets and the asphericity classes. Without the aspheric final-object hypothesis, a weak-equivalence set need not determine a unique asphericity structure. The analogy with contractibility and homotopism structures in \S\S51--52 motivates a promised construction from contractibility to asphericity and a uniqueness assertion for the intended canonical examples (p.\ 235). Proposition~4 recasts the criteria of p.\ 214 as six conditions: three about aspheric objects and three about weak equivalences. The object conditions are equivalent and imply the others; the revised converse requires both $A$ and the asphericity structure to be aspheric. All six conditions become equivalent when the specified generating category $B$ contains the final object. If no $B$ is fixed, enlarging one to contain $i(A)$ and the final object yields equivalence of the first five conditions (pp.\ 236--237). These added hypotheses, some typed interlinearly and some handwritten, are central to the proposition and were partly concealed by malformed insertion macros in the draft. Section~76 immediately tests the formalism on examples.
Two visible mathematical slips are retained for the reader to assess: the universe paragraph says “two objects of $U$” where the witness substitutes $M$ (p.\ 230), and Proposition~2 says “Then $M$ is” a mildly saturated subset of $\mathrm{Fl}(M)$ after defining $W$, where the witness prints $W$ (p.\ 234). Proposition~3 and its corollary also call the relevant map “aspheric,” while the explicit revisions of Proposition~4 replace that word by “a weak equivalence” (pp.\ 234, 236); the edition does not silently harmonize the earlier wording. Several correction layers remain uncertain, including the usual-localizer subscript at the preceding section boundary and overtyped letters in these definitions and proofs.

\subsection*{\S76. Basic examples and total asphericity (pp.\ 237--239).}
Three examples test the definitions of \S75. For $M=(\mathrm{Cat})$, the $\underline W$-aspheric categories form the basic asphericity structure, generated through $a\mapsto A_{/a}$ for a weak test category $A$. The shift from the basic modelizer to this structure puts aspheric objects before weak-equivalence arrows. The distinction matters for functors into $(\mathrm{Cat})$: asphericity supplies the induced localized functor, whereas the weak-test requirement asks for a homotopy-model comparison. Theorem~1 of p.\ 176 suggests equivalence of these conditions when the values are contractible, and the author identifies a comprehensive answer to the “silly question” of \S46, p.\ 95, as the motivation for the present formalism (pp.\ 237--238). The second example takes topological spaces with the weakly contractible spaces as aspheric objects; the geometric realization of simplices supplies the testing functor, and the text appeals to Gabriel--Zisman for the stronger modelizer statement. Grothendieck adds the previously omitted assertion that the first two structures are aspheric and later notes that the resulting weak equivalences are the usual ones (p.\ 238). The third example is $A^\wedge$ with its usual aspheric presheaves, generated by the canonical fully faithful inclusion of any small $A$. Its weak equivalences are $\underline W_A$, while its asphericity is equivalent to asphericity of $A$; the modelizing property is not asserted for arbitrary $A$. A functor $A'\to A$ is aspheric precisely when its composite into this presheaf structure is $M_{as}$-aspheric (pp.\ 238--239). This leads to “totally aspheric” structures: $M$ admits finite products and finite products of aspheric objects, including the empty product, are aspheric. Proposition~5 identifies this condition with the existence of a small generating subcategory stable under finite products, and with a testing functor from a category whose presheaf category is totally aspheric. It reconnects the product arguments of \S74 with the general formalism and prepares the localized canonical functor of \S77 (p.\ 239). The visual check restores a missing typed a before the repaired l in “realization,” retains the literal thatnthis and appheric typing, and distinguishes typewritten interlinear additions from pen text. The upward margin arrow at the end of the examples had been misrepresented by three brackets in the draft; its direction is restored, without attributing a mathematical meaning to the layout mark.

\subsection*{\S77. The canonical homotopy functor and the meaning of modelizing (pp.\ 239--241).}
An asphericity structure now determines more than its weak-equivalence class: it supplies a canonical functor from $\mathrm{Hot}_M=W^{-1}M$ to $(\mathrm{Hot})$. Choosing a testing functor $i:A\to M$, Grothendieck uses the composite through $A^\wedge$ and its category-of-elements functor; the inverse-image descriptions of the weak equivalences permit localization, while conservativity is stated conditionally on strong saturation (pp.\ 239--240). To compare two choices, he places their images in one small full subcategory $B$ of $M_{as}$, obtains the comparison isomorphisms from aspheric functors, and then enlarges $B$ to prove independence. A common $B$ for three choices gives transitivity, completing the canonical construction promised in \S74. Formula (13) on p.\ 240 actually names its composite $i_M$, while the following localized functor carries the extra choice index; the witness's harmonization to $i_{M,i}$ in (13) is not adopted. The two sloping comparison signs there are hand-drawn bars with tildes, without arrowheads. The new definition calls the structure modelizing when this particular canonical functor is an equivalence; this is stronger than merely requiring an abstract equivalence between $W^{-1}M$ and $(\mathrm{Hot})$. In the presheaf example it says exactly that $A$ is a pseudo-test category (p.\ 241). The author immediately hesitates over the terminology and proposes a parallel to the hierarchy of pseudo-test, weak test, test, strict test categories and contractors, but leaves the final qualification as question marks. His deletion of “progressively stronger” in that discussion remains visible; the witness prints the deleted words as live text. Section~78 therefore begins with the modelizing name expressly provisional, and uses the canonical construction to revisit the observation and silly question of \S46 rather than treating the proposed hierarchy as already settled.

\subsection*{\S78. Test functors separated into asphericity and modelizing conditions (pp.\ 241--246).}
The provisional definition of modelizing from \S77 supplies the setting for a reformulation of the observation and silly question of \S46. Theorem~1 assumes a modelizing structure $(M,M_{as})$, a pseudo-test category $A$, a functor $i:A\to M$ with aspheric values, and the existence of a final object of $M$; the last hypothesis is a typed interlinear addition, not an original underlined parenthetical sentence. It equates the asphericity and weak-equivalence criteria of Proposition~4, p.\ 236, with the condition that $i^*$ induce an equivalence between the localized categories (p.\ 242). A lemma says that a final object in a modelizing structure is aspheric. Its supporting corollary characterizes $M_{as}$ as the objects whose images under the canonical map to $(\mathrm{Hot})$ are final. The proof uses strong saturation of the basic localizer, announced as a consequence of Loc~1--3 but deferred to part~V; the author explicitly offers the additional asphericity assumptions on the two final objects as a way of avoiding the feared circularity before that proof is supplied (pp.\ 242--243). The canonical triangle from \S77 then proves the localized equivalence. The proof's literal reference to conditions through (ii') is retained even though the statement speaks of six conditions through (ii'').
The theorem prompts a revision of the terminology inherited from \S65. The evolving pseudo-test-functor definition first describes model-preserving composites, then separates the intended conditions into: $A$ is pseudo-test, $(M,M_{as})$ is modelizing, and $i$ is $M_{as}$-aspheric. The discussion revises both its hypotheses and its formulation on the page; these are stages of the argument rather than interchangeable completed definitions (pp.\ 243--244). Passing to weak test functors strengthens the condition on $A$ to weak test. Passing to test functors requires a test category and asphericity both of $i$ and of all its slice-induced functors. Definition~2 calls the latter family of slice conditions locally aspheric, and the conjunction with asphericity of $i$ totally aspheric. The local criterion additionally requires that $i$ factor through $M_{as}$ and that $i^*(x)$ be aspheric over the final presheaf, with $x$ restricted to a generating subcategory containing $i(A)$ when desired (p.\ 245). The Remarks paragraph draws the conceptual conclusion: the independent asphericity conditions have wider relevance than the combined test-functor terminology. Yet the earlier theorem on $(\mathrm{Cat})$ still has one advantage, namely its criterion involving only $i^*(\Delta_1)$ under the contractibility hypothesis on the values of $i$. Recovering such a minimal criterion requires relating asphericity structures to the contractibility structures of \S51, which becomes the subject of \S79 (p.\ 246). The typed minimal is inserted before criterion without deleting generating in the preceding line. Most of the final sentence on p.\ 246 is physically covered by a laid-over sheet; its tentative transcription remains doubtful and is not treated as fully recovered evidence.

\subsection*{\S79. From contractibility to asphericity and the general test-functor criterion (pp.\ 247--255).}
The entry dated 20 June supplies the missing bridge identified at the end of \S78. Proposition~1 removes a smallness assumption on a generating full subcategory $N$: the existence of a small generating part, a small inclusion detecting asphericity on $N$, and the assertion that $(N,N)$ is an asphericity structure are equivalent. The name “aspherator” emphasizes that the condition belongs to $N$ itself, while its dependence on the chosen universe and basic localizer remains explicit (p.\ 247). The author then strengthens the standing convention for contractibility structures to require a small generating full subcategory and recalls the finite-product assumption on $M$. Proposition~2 makes $M_c$ generate a totally aspheric structure $M_{as}$; its clause about a chosen generating category $C$ acquires the additional handwritten hypothesis that $C^\wedge$ be totally aspheric (p.\ 248). The proof uses finite-product closure, contractibility of the induced presheaves, and Proposition~5 of p.\ 239. The following remark distinguishes this construction from arbitrary asphericity structures and asks about a canonical largest contractibility structure generating a given one. Its claim about every previously encountered modelizing structure is immediately retracted using nonstrict test categories, then narrowed to modelizing totally aspheric structures (pp.\ 249--250). The second part recalls locally and totally aspheric presheaves and the homotopy interval structure generated by the locally aspheric ones. The source's reference to part IV, \S\S54--55, is retained as written. Six equivalent compatibility conditions H1--H4' express preservation of homotopies, homotopisms, intervals, or contractible generators by a product-preserving functor; Remark~2 explains exactly what generation by $C$ means (pp.\ 250--252). Theorem~1 applies this to $i^*:M\to A^\wedge$ when $i:A\to M$ lands in $M_c$: compatibility with the homotopy structures, local $M_{as}$-asphericity, and the same presheaf criterion restricted to $C$ are equivalent. Its corollaries add asphericity of $A$ for total asphericity, and the test-category and modelizing requirements for a test functor. A separating multiplicative interval supplies a sufficient criterion forcing $A$ to be a local test category; later margin notes question smallness of $C$ and say that multiplicativity is unnecessary, pointing to p.\ 348, Remark~1 (pp.\ 252--253). The proof abandons a boxed localization argument in favor of a lemma and corollary about testing on generators; the blank relation positions in the deleted draft are not completed editorially (p.\ 254). Remark~3 specializes to totally aspheric $A$ and obtains a strict test category under the interval hypothesis, closing the promised restatement of the earlier result. Its corollary reference is corrected by hand from 2 to 3, and a separate margin note points to a stronger version on p.\ 347 (p.\ 255). Section~80 then reverses the direction of inquiry, asking what homotopy structure is recovered from the associated weak equivalences.

\subsection*{\S80. Recovering homotopy from weak equivalences and the canonical-modelizer proposal (pp.\ 255--258).}
The construction $M_c\mapsto M_{as}$ from \S79 is set beside the reverse-looking association $W_a\mapsto h_{W_a}$ of \S54. The first direction distinguishes contractibility from the weaker homotopy interval and homotopism structures; it also makes clear why the associated weak equivalences $W_a$ should not be confused with the generally smaller class $W_c$ of homotopisms. The author recalls how the aspheric-object viewpoint emerged after the earlier “awkward approach,” and corrects and supplements the references to that development (p.\ 255). For an asphericity structure with aspheric final object, $W_a$ and $M_{as}$ determine each other. The intervals generating $h_{W_a}$ are those for which every projection $x\times I\to x$ lies in $W_a$. Under total asphericity this is equivalent to asphericity of $I$, an observation promoted from a tentative expectation to an immediate statement. Two questions then ask whether passing from $M_c$ to $M_{as}$ and back recovers the same homotopy structure, and which extra conditions on $M_{as}$ ensure that the recovered structure comes from a contractibility structure generating it (p.\ 256). A third question, now redundant with the preceding observation, is boxed and deleted; its formula reads $UW_a$, with both letters visible (p.\ 257). Before answering, Grothendieck resumes the “idyllic picture” of canonical modelizers from \S50, incorporating \S\S57 and 59. A finite-product category $M$ with a functor $\pi_o:M\to(\mathrm{Sets})$ receives an interval structure from projections inducing bijections under $\pi_o$. If this structure comes from $M_c$, then $M_{as}$, $W_a$, and the canonical functor $W_a^{-1}M\to(\mathrm{Hot})$ follow. The further compatibility between $\pi_o$ and the composite through $(\mathrm{Hot})$ remains to be imposed; the mild localizer condition L~a) on p.\ 165 is needed to define the final connected-components functor (pp.\ 257--258). When $M$ is totally 0-connected, $\pi_o$ is its intrinsic connected-components functor, and the resulting asphericity structure is modelizing, the construction deserves the name canonical modelizer because it depends on the category itself. The closing Remark calls the definition essentially complete while explicitly retaining the compatibility question. This unresolved point, and the two recovery questions, lead into \S81. The visual collation distinguishes the typed interlinear hypotheses from their pen locators and preserves the literal “homotopisms into homotopisms” on p.\ 257, where the witness substitutes “isomorphisms.”

\subsection*{\S81. Recovering contractibility and comparing the three families of notions (pp.\ 259--266).}
The entry dated 25 June announces the positive answer to the first recovery question of \S80: a contractibility structure $M_c$ on a category with finite products is recovered from its associated asphericity structure by taking the objects contractible for the homotopy interval structure generated by all aspheric intervals. Thus $\mathrm{Cont}(M)\to\underline{W}\text{-}\mathrm{Asph}(M)$ is injective when the basic localizer satisfies Loc~4; the proof is deferred to \S82. The four homotopy notions recalled in (2)--(4) place this assertion beside the earlier inclusions from contractibility to interval structures and homotopism structures, with the latter equivalent to homotopy relations (p.\ 259). Changing the basic localizer from $\underline{W}$ to a coarser $\underline{W}'$ induces a canonical change of asphericity structure, and the triangle on p.\ 260 makes this construction compatible with the inclusion of contractibility structures. Loc~4 requires weak equivalences of small categories to induce bijections on connected components. Among localizers satisfying it, $\underline{W}_o$ is the coarsest; the finest basic localizer $\underline{W}_{\mathrm{oo}}$ is introduced as their intersection and its identification with the usual weak equivalences is promised for part V. The strongest recovery statement is therefore the elementary case associated with $\underline{W}_o$ (pp.\ 260--261). Writing $M_o$ for the 0-connected objects and $W_o$ for the corresponding maps, (17)--(20) show how either determines the other and how both are recovered from the connected-components functor $\pi_o:M\to(\mathrm{Sets})$. The first line of (19) literally omits “is” before “a one-point set.” A handwritten marginal note objects to the name 0-connected maps and proposes 0-équivalences, but its crowded continuation remains partly unreadable; the later asterisk on p.\ 266 refers back to this objection (p.\ 262). The retrospective discussion explains how the special canonical connected-components functor delayed recognition of this general construction and why much of part III is superseded by part IV, while its homotopy notions remain essential (pp.\ 263--264). The three families are test notions, homotopy notions, and asphericity notions. Test and asphericity notions depend on the basic localizer; homotopy notions are absolute. Their order in a deductive account would reverse much of their order of discovery. Contractibility becomes the hinge between homotopy and asphericity: it is both the strongest of the four homotopy structures and a specially constrained asphericity structure (pp.\ 264--265). The revised displays then enumerate ten mutually determining data, including $M_c,J_c,W_c,R_c$, the 0-connected and $\mathrm{oo}$-aspheric objects and maps, and the two canonical functors to sets and homotopy types. This count is reached through visible revisions from four to three asphericity-flavored data and from eight through nine to ten structures. Two further relations $R_o,R_{\mathrm{oo}}$, and more generally $R_{\underline{W}}$, compare maps having equal images under the corresponding functors; whether such a relation alone recovers $M_c$ is explicitly left unclear. The inclusions and commutative comparison in (26), followed by the localization chain in (27), complete the section (pp.\ 265--266). The checked source preserves the typed “oo” notation, identifies the underlined $\pi_o$ in (26), and distinguishes pen arrows from typed diagram objects. Its revision record also retains the five-dot reinstatement of “being” on p.\ 259, the phrase transposition on p.\ 260, and the cancelled draft displays on p.\ 265.

\subsection*{\S82. Proof of recovery and consequences for internal Hom objects (pp.\ 267--272).}
Theorem~1 supplies the proof promised in \S81. Under Loc~4, with finite products and the stipulated small full generating category, the contractible objects for the homotopy interval structure generated by all aspheric intervals are exactly the original $M_c$. The argument compares $M_c$, the recovered contractible objects, and $M_a$, then reduces the essential point to Lemma~1: any two sections of an aspheric object are already $M_c$-homotopic (pp.\ 267--268). Passing to the elementary localizer and a small full generating category $C$ reduces the problem to its presheaf topos. Lemma~2 restates the finite-chain argument from \S60, p.\ 149: in a totally 0-connected topos with the stated section and generator assumptions, two sections of a 0-connected object can be joined through a finite chain of objects of $C$. The proof of this lemma is left to the reader. The displayed inclusion $\underline W_o\subset\underline W$ on p.\ 268 is retained as printed, although the witness reverses it; it should be read alongside the coarseness discussion on pp.\ 260--261.
The next application assumes the relevant internal Hom objects exist. For a pointed object $X$, the endomorphism object $\underline{\mathrm{Hom}}(X,X)$ has the identity and the constant map as endpoints and supplies a multiplicative interval. The resulting equivalences connect contractibility of $X$, contractibility or 0-connectedness of internal Hom objects, and homotopy equivalence of the constant-map inclusion $Y\to\underline{\mathrm{Hom}}(X,Y)$ (pp.\ 269--270). These are consequences of the recovered homotopy structure, not new definitions of the asphericity class. The source uses $\pi_o,\pi_1$ for the endpoints in one display; that notation is preserved. Lemma~3 compares homotopy classes of maps from $Y\in C$ with the corresponding classes after applying $i^*$, referring back to p.\ 144. Together with the availability of sections of contractible generators, this gives Proposition~1: homotopy classes of sections map bijectively to connected components. Its internal-Hom corollary identifies homotopy classes of maps $X\to Y$ with the connected components of $\underline{\mathrm{Hom}}(X,Y)$ when that object is representable (pp.\ 270--271).
Proposition~2 generalizes the product argument of \S74, p.\ 228, under internal-Hom hypotheses. The revisions add total asphericity of the indexing category and alter several occurrences of the generating category and object variables; the cancelled final criterion remains visible in the diplomatic account (p.\ 271). Setting $B=C=M_c$ yields the corollary that the product of an aspheric functor with a functor valued in $M_c$ is aspheric. The author explicitly wonders whether the internal-Hom existence assumption for pairs of contractible objects is superfluous. Hence the category of aspheric functors valued in $M_c$ is stable under binary products. If nonempty, it gives the canonical functor $\mathrm{Hot}_{M,\underline W}\to\mathrm{Hot}_{A,\underline W}$, defined up to unique isomorphism through the transitive system recalled from p.\ 230 (p.\ 272). This closes the recovery proof and its product application; the next section asks for a converse characterization of the asphericity structures to which it applies.

\subsection*{\S83. A tautological characterization and an unresolved functor question (pp.\ 272--273).}
Grothendieck returns to the question omitted from the preceding theorem: which asphericity structures arise from contractibility structures? He first imposes total asphericity, then forms the interval structure $J$ generated by intervals whose objects lie in $M_a$ and calls its contractible objects $M_c$, warning that this set need not itself be a contractibility structure (pp.\ 272--273). The two conditions are that $M_c$ generate $J$, equivalently that two sections of any object of $M_a$ be $M_c$-homotopic, and that $M_c$ generate the given asphericity structure through a small full subcategory with aspheric inclusion. The first condition ensures that $J$ does arise from a contractibility structure, necessarily this $M_c$. The characterization is expressly called a tautology. A parallel request asks when a functor $\phi:M\to\mathrm{Hot}_{\underline W}$ is isomorphic to one associated with a contractibility structure, especially when it reduces to $\pi_o:M\to(\mathrm{Sets})$. The text offers no useful necessary-and-sufficient criterion and leaves open whether suitable exactness and internal-Hom assumptions would produce one. Its final underlined question asks whether such a functor can have nontrivial automorphisms, connecting this inquiry to the “inspiring assumption” of \S28 (p.\ 273). The checked text preserves the visible overtyped ending of “generated,” the parentheses around Sets, and “writing it down”; it does not replace these with the witness's readings.

\subsection*{\S84. A provisional notion of morphism and the direct-image problem (pp.\ 274--280).}
The entry dated 29 June begins with letters to Faltings and Porter, then describes the unexpectedly rapid emergence of a proposed morphism notion after prolonged scratchwork (p.\ 274). For two asphericity structures and a functor $f^*:M\to M'$ with left adjoint $f_!$, the proposition parallels the six conditions of \S75, p.\ 236: preservation and detection of aspheric objects, and corresponding assertions about weak equivalences. Its extra assumption (Awk) requires an aspheric testing functor into $M'$ whose composite with $f_!$ lands in the chosen generating part of $M$. The page begins with a generating category $B$ but uses $C$ in that auxiliary clause; the source wording is retained. A triangle through $A^\wedge$ reduces the equivalences to the earlier proposition. The revised conclusion first equates five conditions under the aspheric-final-object hypotheses, then all six if the final object belongs to the chosen generating category (p.\ 275). The canonical localization functors give a second triangle, and the corollary proves that a morphism between modelizing structures induces an equivalence of localized categories. The author calls this the answer to the “silly question,” with the source reference \S45, p.\ 95 (p.\ 276).
The rest of the section examines the restrictions rather than treating the proposed definition as settled. The left-adjoint requirement is compared with preservation of inverse limits. Condition (Awk) is reformulated through a small generating subcategory of $M'_a$ taken by $f_!$ into $M_a$, or through the generating property of $(f_!)^{-1}(M_a)\cap M'_a$. Replacing it by the stronger condition $f_!(M'_a)\subset M_a$ would pair it neatly with $f^*(M_a)\subset M'_a$, but its reasonableness is unresolved (pp.\ 276--277). Question~1 asks whether the left adjoint $i_!:A^\wedge\to M$ attached to an aspheric functor preserves aspheric objects. In the presheaf case this becomes Question~1': when does $i_!:A^\wedge\to B^\wedge$ preserve them? Dropping asphericity of $i$ reduces the issue to the colimit of the representables $i(a)$ over an aspheric category. A cofibering example with 0-connected fibers shows why asphericity of the indexing category alone does not suffice; asphericity of $B$ must at least enter the discussion (pp.\ 277--278).
The proposed counterexample on p.\ 279 concerns an amalgamated sum $a\amalg_c b$. The monomorphism case recalls \S70, but the author explicitly has not completed a general counterexample. Neither the extra initial-object condition nor total asphericity appears convincing to him; imposing a strict test category leads to a further question he does not know how to handle. The checked edition preserves the inward arrow directions of diagram (9), despite the surrounding reference to a dual category, and the capital $B$ in the inline map $c\to B$. Finally, for $M=(\mathrm{Cat})$, the canonical slice functor $a\mapsto A_{/a}$ supplies the positive example when $A$ is a weak test category, with both associated functors modelizing and the left adjoint preserving aspheric objects. The standard inclusion $\Delta\hookrightarrow(\mathrm{Cat})$ is the decisive unresolved test: its left adjoint is not generally modelizing, and preservation of aspheric objects still has to be checked (p.\ 280). At the opening of \S85 on the same page, Grothendieck retracts the provisional morphism notion because stability under composition has not been established; this later objection governs how the preceding proposition and corollary should be read.

\subsection*{110. Schematization through Postnikov fibrations, and its obstruction (pp. 443--445; 134-7/1--3).}
Grothendieck resumes the question of homotopy types over a ground ring by comparing an abelian group $M$ with the functor $W(M)(k)=M\otimes_{\underline Z}k$: the discrete object acquires a structure analogous to a vector bundle, suggesting an analogous construction for semisimplicial models and successive Postnikov fibers. The proposed ``schematization above level n'' would leave a lower truncation discrete and require the comparison from schematic cohomology with quasi-coherent coefficients to ordinary cohomology to be bijective. The reference to section 94 situates the hope of replacing $\underline Z$ by $\underline Q$ when nilpotence replaces an abelian fundamental group. On p.~444 he immediately questions this construction: torsion and the universal coefficients formula prevent the homotopy groups after base change from being obtained by tensoring the individual homotopy groups. The alternative is to reconstruct a homotopy type from a suitably enriched ``abelianization'' or ``linearization''; section 111 then seeks a different construction without explicit Postnikov decomposition. The handwritten replacement of ``(= homotopy)'' by ``(or rather, homotopy)'' on p.~444 matters here: the arXiv witness retains the cancelled equality wording. Its mapsto in the displayed base-change formula also differs from the scan's ordinary arrow; the scan alternates star subscripts with the typed notation X. Page 445 closes with a methodological reflection on the tension between detailed preparatory work and the desire to move quickly. The rust-obscured words on p.~443 remain marked rather than restored from the mathematical argument.

\subsection*{111. Unipotent bundles, divided powers, and a revised notion of equivalence (pp. 446--459; 134-7/4--17).}
The starting problem is to construct categories $Hot_n(k)$ of schematized $n$-connected homotopy types, with extension and restriction of the ground ring and a comparison with the ordinary homotopy category over $\underline Z$ (pp.~446--449). Grothendieck builds $U(k)$ inductively from affine bundles $W(M)$ and torsors, first in the fpqc setting. He asks whether the successive fibration data belong to the object or merely witness its existence, then proposes semisimplicial objects subject to connectivity and filtration conditions. The proposed weak equivalences are initially detected by taking sections, and the hoped-for $k$-module structures on homotopy groups lead to a search for a linearization into $D.(Ab_k)$. On p.~448 the model-category symbol has a typed subscript $m$, although the surrounding development uses $n$; this apparent index slip is retained in the text. Pages 450--455 test divided powers as the required linearization. The candidate universal map $x\mapsto\sum x^{(i)}$ immediately forces a completion, and the completed divided-power coalgebra is considered through its group-like elements. After the Eilenberg--Maclane test, the discussion returns to pointed objects and uncompleted divided powers. The comparison (w) on p.~454 places ordinary free abelian linearization, $L$, and the pointed version $L_{pt}$ in one diagram. The marginal correction on p.~452 is consequential: ``One has to take k=Z, and components $X_n$ projective.'' Both its initial words and the subscript are legible on the scan. The subsequent distinction between quasi-isomorphisms and chain homotopisms supplies the projectivity restriction. The final reference in the sentence on p.~455 is to (x''), after a preceding reference to (x'''); the arXiv witness repeats (x''') and loses this distinction. The argument is substantially revised on pp.~456--459. The author retracts the unrestricted derived-functor claim of p.~454, requires projective or flat components, and remarks that the Postnikov approach must again be set aside. The plain typed $I_k(M_*')$ on p.~456 is retained as printed, although the neighboring formulas suggest a notation slip. Page 457 then exposes the problem with ordinary homology as a $k$-module: the exterior-power calculation for $K(M,1)$, applied to $k=\underline Z[T]/(T^2+1)$, yields an obstruction. Page 458 accordingly defines total linearized homology $LH.(X_*)=L_{pt}(X_*)$ and considers weak equivalences detected by its homology invariants. The handwritten revision says it is becoming \emph{dubious} whether equivalences should be defined through $X_*(k)$; the arXiv witness reads \emph{obvious}, reversing the author's change of judgment. At the end of the section, Question 2 asks whether the two definitions coincide for projective pointed unipotent bundles satisfying $X_0=X_1=e$ (p.~459). The weaker condition $X_0=e$ raises an additional issue about abelian fundamental groups. These pages formulate questions and revised expectations, rather than establish an unrestricted equivalence theorem; section 112 resumes from the projective $\underline Z$ case.

\subsection*{112. Kan conditions and the scope of the linearization program (pp. 460--465; 134-7/18--23).}
The entry dated 26.8 takes the projective $\underline Z$ case of the comparison maps (w) as its working foothold. Grothendieck sketches a passage from the Eilenberg--Maclane case to finite chain complexes, then through finite and ``way-out'' infinite successive fibrations, using multiplicativity and hoped-for spectral-sequence or derived-category transitivity statements (pp.~460--461). The proposed Kan--Postnikov complexes use shifted projective resolutions of homotopy groups. To transfer the construction from discrete to schematic objects, he returns explicitly to the cohomology obstruction of p.~444 and hopes to dualize the homology comparison. This is an outline requiring substantial verification, as the author himself stresses. Constant complexes then show why $X_0=e$ is necessary, and the monoid discussion turns the issue into comparison of polynomial with arbitrary cochains. The multiplicative-monoid example on p.~462 is offered as a counterexample to the assumptions previously used; the circled marginal cross is part of the page's evidence and is absent from the arXiv witness. The response is to ``drop Postnikow and keep Kan'' (p.~462). Since $U(k)$ is not closed under arbitrary fiber products, a schematic horn condition must ensure both representability of the relevant finite limits and a suitably strong surjectivity property (p.~463). The tentative name ``submersions'' covers alternatives ranging from trivial bundles to formally smooth maps admitting sections. The first nontrivial condition asks for the three maps $X_2\to X_1\times X_1$ to be submersions; for a monoid nerve this is meant to force a group. Page 464 lists the desired consequences: the linearization theorem over $\underline Z$, a homology criterion for homotopisms of connected Kan complexes over a general ring, homotopy fibers and Postnikov decomposition, and eventually equivalence of $Hot_n(\underline Z)$ with ordinary $Hot_n$. Page 465 distinguishes restriction of rings from extension, noting that restriction does not preserve projective components without an additional assumption and suggesting resolutions to remove that restriction. Its last sentence proposes compatibility of ring extension with the individual homotopy invariants; the next entry immediately revises that proposal. This transition is essential to reading the section as a developing argument rather than a list of established properties.

\subsection*{113. Soft Postnikov decomposition and a schematic fundamental group (pp. 466--472; 134-7/24--30).}
The entry dated 27.8 corrects the previous evening's assertion. Formula (A) concerns total homology and derived extension of scalars; the analogous termwise assertion (B) for homotopy modules fails without flatness hypotheses (p.~466). The revised arrow in (B) points left, a direction separately emphasized by a circled marginal arrow omitted by the arXiv witness. The reason is the use of shifted projective resolutions in a Postnikov fiber: after tensoring, higher homology can appear, so the base-changed successive fibration need not remain a Postnikov decomposition (p.~467). This explicitly answers the puzzle of p.~444 and motivates the judgment that ordinary Postnikov decomposition is not ``geometric,'' meaning compatible with extension of the ground ring. The proposed alternative keeps the entire abelianization complex. Using the completed divided-power linearization and the embedding $X\to W(L(X))$, Grothendieck proposes repeatedly taking the homotopy fiber of $X_*\to W(L(X_*))$, without first extracting the first nonzero homology module (pp.~467--468). He names this ``soft Postnikov d\'evissage.'' For a discrete $K(G,1)$ the analogy leads to iterated commutators, and he distinguishes the versions associated with solvability and nilpotence. Formula (D) contains an index problem: the visible inner exponent is $n+1$, like the outer exponent, while the arXiv witness gives $n-1$. The checked text preserves the scan's reinforced plus and does not silently substitute a mathematically expected recurrence. The next page adds the issue of the action of $\pi_1$ on higher homotopy groups; trivial or unipotent action appears relevant to unipotent schematization (p.~469). The second part of the section revisits complexes with $X_0=e$ but possibly nontrivial $X_1$. A group object gives a fundamental group that need not be abelian, and finite-presentation and smoothness assumptions lead to questions about nilpotence (p.~469). The Kan relations then motivate the presheaf $k'\mapsto\pi_1(X_*(k'))$. The desired invariant is a group object of $U(k)$ representing this sheaf, with an underlying parametrization by a module and a polynomial group law (pp.~470--471). This changes what the notation $\pi_1$ denotes: the group of sections alone is no longer the complete object. Formulas (G) and (G') express evaluation and base-change compatibility for the proposed schematic group. The following duplicated equation on p.~471 is nevertheless literally tautological in the scan, with $\pi_1(X_*(k))$ on both sides; the witness edits one side. The blank between $k$ and $k'$ in the preceding prose likewise contains no arrow on the page. To admit arbitrary quotient modules, the notion of submersion is reconsidered once more: maps isomorphic to those induced by epimorphisms of modules are proposed as the most convincing of the three possibilities considered so far (p.~471). The printed ``unless \ldots{} does'nt give rise'' contains an apparent extra negation, retained despite the witness's change to ``gives rise.'' The stronger requirement that the schematic fundamental group come from addition on a module is distinguished from commutativity; the discussion ends with tests involving the linearization theorem, a fibration over $K(G,1)$, and the relationship between nilpotent discrete groups and group objects of $U(k)$ (p.~472, referring back to section 94). The handwritten addition about unipotent action on the higher $\pi_i$ belongs here and is not merely incidental: it qualifies the contemplated range of homotopy types. Section 114 begins on the same page by describing the accumulated picture as provisional and preparing a review of the program.

\subsection*{114. A provisional model category and the completion gap (pp. 472--480; 134-7/30--38).}
This review arranges the preceding experiments into a prospective construction. The category $U(k)$ consists of unipotent bundles built from modules; smooth objects correspond to projective modules, submersions to epimorphisms, and smooth submersions to maps with projective kernel (pp.~472--473). The completed divided-power coalgebra describes the unpointed linearization, while the pointed version removes the completion in the basic construction (pp.~474--475). Completion is a substantive distinction: p.~474 explicitly says that the map (C) of p.~467 ``does'nt quite exist'' with $W$ as written, and introduces the completed version $W^{\wedge}$ instead. Base change for the completed functor consequently uses completed tensor products. The subsequent model categories impose $X_i=e$ for $i\leq n$, take sections to pointed semisimplicial sets, and localize at the proposed weak equivalences. The hoped-for comparison with ordinary homotopy types is restricted to the appropriate connected range over $\underline Z$ (p.~476). On p.~477 the first expression for the category of models visibly uses $\pi_n(k)$, although the surrounding formulas use the handwritten model-category symbol $\mathcal M_n(k)$; this notation anomaly is retained rather than silently regularized. Smooth Kan resolutions are then proposed to derive extension of scalars and the total linearization $LH.(X_*)=L_{pt}(X_*)$ (p.~477). The next two pages ask whether this homology functor is conservative and whether pointed linearization and completed linearization are quasi-isomorphic. They also reconsider comparison with ordinary free-module homology. The second question undergoes an immediate correction: on p.~479 the author says the proposed comparison with $k^{(X_*(k))}$ is generally false away from $k=\underline Z$. In the first nonvanishing degree, the expected $k$-module invariant is identified with integral homology of the section complex; tensoring that homology with $k$ produces a different candidate. The word \emph{integral} is explicit in the scan, whereas the arXiv witness gives \emph{internal}. The correction renews the need to construct homotopy modules intrinsically and to obtain them through a Postnikov argument, rather than transfer all ordinary homology assertions indiscriminately. The final page makes the outstanding construction precise. The exponential map gives $X_*\to W^{\wedge}(L(X_*))$, but the required homotopy-class map is $X_*\to W(L_{pt}(X_*))$ (p.~480). Even a quasi-isomorphism between the two linearizations does not by itself remove the difference between $W^{\wedge}$ and $W$; the handwritten phrase ``flaw in this `argument''' marks exactly this gap. The target of the completed map is a semisimplicial sheaf and need not belong to the proposed model category. Hard or soft Postnikov decomposition therefore still requires this factorization, a suitable notion of homotopy fiber, and homotopy constructions internal to the model category. The failure of $U(k)$ to have fiber products explains why Quillen's standard machinery cannot simply be applied as stated. These are explicit unresolved tasks, and section115 returns the following morning to the meaning of the coefficient sheaf itself.

\subsection*{115. Quasicoherent coefficients and prorepresentable linearization (pp. 481--489; 134-7/39--47).}
The entry dated28.8 restarts the formula numbering and identifies the source of the preceding confusion. The appropriate free linearization is over the quasicoherent ring sheaf $\underline O_k$ on the fpqc site, not over a constant ring (pp.~481--482). Grothendieck explains that the obstruction groups in schematic Postnikov decomposition use the quasicoherent sheaf associated with a homotopy module, rather than the same module regarded as a constant coefficient group. The circled marginal clarification of the argument $X(n)_*,\pi_{n+1}$ belongs to this distinction; the overwritten version in the running formula remains marked as doubtful. He connects the earlier mistake to having translated geometric language into commutative algebra too thoroughly, losing sight of which coefficient objects were intended. The universal free Module $\underline O_k^{(X)}$ gives the Hom, Ext and derived-Hom expressions (4)--(5') (p.~483), but is generally not quasicoherent. The proposed replacement would represent the same functor on quasicoherent coefficients. For $X=W(M)$, the continuous-Hom calculation through truncated divided-power modules shows that the correct object is a proobject, rather than the completed module regarded as an ordinary module (pp.~484--485). The obstruction is again scalar extension: inverse limits do not generally commute with tensoring by a new ground ring. Formula(13) packages the finite products of divided powers into a strict projective system; ordinary representability would require it to become essentially constant. The added phrase ``but with slightly different meaning'' therefore changes the interpretation of $L$ in (14), now valued in pro-modules. The pointed functor still takes values in ordinary modules, embedded as constant proobjects. Formula(18), on p.~486, supplies the exponential map into $WL(X)$ in the procategory and distinguishes this abelianization from its richer linear structure. The proposed derived comparison (20) and its Ext form are then tested on $X=W(M)$ (pp.~486--487). Formula(21) contains a meaningful handwritten revision from a limit expression to a direct sum of Ext groups of divided powers. Projectivity of $M$ gives the expected vanishing in positive degree, with a passage to finite-type submodules sketched for the infinite case. For nonprojective $M$, the potential nonzero Ext$^1$ term challenges the author's previous expectation that quasicoherent cohomology of these bundles should vanish. His reference to Larry Breen is a question about this expectation, not a proof of the general formula. The text visibly says ``quasiunipotent'' where the witness shortens it to ``unipotent'' (p.~487). The next page restricts attention to the smooth case and postpones the projective-resolution argument for full generality. For a semisimplicial complex, (22) uses Hyperext rather than a degreewise Ext construction; (23) compares the resulting pro-linearization with pointed linearization (p.~488). The comparison map(24) revisits the map from p.~454 with its proobject interpretation made explicit. Its relationship to a quasi-isomorphism, and then to the componentwise inverse limit used on p.~478, still requires foundational work. Page489 plainly admits that no argument has yet been found for the desired assertion under the proposed smoothness, connectivity and possibly Kan hypotheses. The second comparison is reformulated as $\underline O_k^{(X_*)}\to L(X_*)$ in (25). The author does not confidently assert that this is a quasi-isomorphism: quasicoherence of the first homology sheaves and essential constancy of the second remain in doubt. The practical hope is instead that it gives the needed comparisons against quasicoherent coefficients through(20). Section116 subsequently reports the collapse of the earlier general abelianization theorem, so these provisional claims must be read with that later revision in view.

\subsection*{116. The failed abelianization theorem and the interruption of the notes (pp. 490--491; 134-7/48--49).}
The entry dated 7.9 reports the decisive revision of the preceding program. After finding a simple description of schematic homotopy groups, Grothendieck learned in a telephone conversation with Illusie that the asserted comparison between discrete and schematic homology or cohomology was false. The anticipated equivalence between schematic and ordinary homotopy types therefore fails (p.~490). This is an explicit retraction of a central assumption underlying sections111--115; it does not cause him to abandon schematic homotopy types, whose formal properties he now wants to study in their own right, possibly through different models. The remainder of the section explains the intervening pause in writing through fatigue and the need to attend to the death of his granddaughter Ella. Page491 distinguishes work allowed to take place through attention from a compulsive forward movement that leaves perceptions disconnected from one's picture of reality. The revisions matter to the voice of this passage: the hand moves ``who'' after ``finally,'' retains ``his language'' for the body's intervention, and changes ``cause'' to ``caused.'' The witness substitutes ``which'' and ``its'' and loses the last correction. These pages provide both the mathematical warning needed to read the preceding conjectures and the author's own account of why the notebook paused before the technical review resumes in section117.

\subsection*{117. Size, exactness, and the distinction between schemes and formal schemes (pp. 492--496; 134-7/50--54).}
The entry dated 10.9 begins with two technical qualifications. The affine site used earlier is not automatically small in the working universe; a possible remedy is to bound the sizes of the modules and the rings used as arguments, using the fppf site for finite-type bundles (p.~492). More consequentially, $M\mapsto W(M)$ is right exact but need not preserve kernels. Thus a free resolution of $M$ need not become a resolution of $W(M)$. The derived-Hom comparison used to interpret formula(21) of p.~487 is therefore questioned in general and retained for the flat case (p.~493). This is a specific dependency of the cohomology calculation, not merely a set-theoretic caveat. The finite projective case clarifies the geometry through duality. Polynomial maps can be described by maps of symmetric algebras or by continuous maps of the dual completed coalgebras; preserving the origin translates into augmentation compatibility (pp.~493--494). The uncompleted divided-power coalgebra is instead dual to the adically completed symmetric algebra. Consequently, the two algebraic descriptions lead to schemes and formal completions along their zero sections, respectively. The scan's formulas require care: the prose on p.~493 includes a dual check on $M$, but formulas(2), (3) and (4) omit dual signs that the arXiv witness supplies. Those omissions remain in the checked text. The upward-arrow-shaped completion mark in(3') on p.~494 is likewise recorded as drawn, and a separate unreadable pen mark above the star in(5) remains marked. This geometric interpretation overturns the earlier hoped-for equivalence between the pointed and nonpointed linearizations (pp.~495--496). For the additive group in characteristic $p$, the degree-one comparison distinguishes additive polynomials from formal series $\sum_{i\geq0}c_i t^{p^i}$: the schematic side allows only finitely many nonzero coefficients. Grothendieck acknowledges that this example is neither simply connected nor over $\underline Z$, but states that he also no longer expects the comparison to be a quasi-isomorphism for the basic connected examples previously contemplated. He reports that Illusie rejected the vanishing assumptions behind his attempted ``way-out'' argument. The closing underline therefore warns explicitly against substituting $L_{pt}(W(M))$ for $L(W(M))$ when computing schematic homology or cohomology. Formal invariants have their own meaning and are distinct. Section118 begins on the same page with a different and deliberately striking claim: homotopy invariants themselves should nevertheless be recoverable from the first-order tangent object.

\subsection*{118. The tangent complex as a proposed homotopy invariant (pp. 496--500; 134-7/54--58).}
Against the distinction just established between schematic and formal cohomology, Grothendieck proposes that homotopy groups should depend only on the first-order tangent object. The Lie functor sends a smooth pointed bundle $W(M)$ to $M$, and formula(9) identifies the homotopy groups of $X_*(k)$ with those of the semisimplicial module $\mathrm{Lie}(X_*)$, hence with homology of its associated chain complex (p.~496). The next page sharply qualifies the claim: he has no direct description of a natural isomorphism under the stated hypotheses. It is tautological for complexes arising from chain complexes and is expected to extend through Postnikov decomposition, if every relevant model can be replaced by one admitting such a decomposition; naturality would still require work (p.~497). The typed revision makes this formula a ``cornerstone'' of a theory of schematic homotopy types and a proposed intrinsic definition of the invariants, rather than subordinating them from the outset to discrete homotopy groups. Weak equivalences would then be detected by the tangent complexes. Both subsequent comparison displays map an $X_*$ object to a $Y_*$ object; the arXiv witness repeats $X_*$ as target in each. The contemplated application is to explicit models of spheres (pp.~498--499). To make the Lie formula useful for ordinary homotopy groups, one must first realize the discrete type schematically over $\underline Z$ and obtain a workable model. Successive suspension from $S^1$ suggests the need for suspension constructions, complementing the earlier need for fibers and loops; the proposed derivator would organize the finite integration and cointegration operations. These are motivations and difficult construction problems, not completed calculations of sphere homotopy groups. The discussion of $S^2$ as a homogeneous space under $S^3$, and of the corresponding group/flag-variety geometry, asks for a model of the entire relationship: a group, a subgroup playing the role of a torus or Borel subgroup, and a quotient. Several overprinted group-name letters remain uncertain in the text, while the star on $\underline C^*$ is visibly an asterisk rather than the witness's multiplication cross. Hemispherical complexes supply a second route (pp.~499--500). Their smaller list of boundary and degeneracy operations, together with a prospective direct computation of homotopy in the discrete setting, is compared with the tangent-complex method. The marginal exclamation and vertical rule emphasize this alternative. The author immediately admits that he has not checked the proposed computation for arbitrary hemispherical sets that do not arise from an infinity-groupoid or infinity-stack. Comparing the two descriptions is therefore offered as a possible route toward understanding formula(9), with that unresolved premise preserved. Section119 then resumes on26.9 after another interruption and returns to the failure of the earlier homology comparison.

\subsection*{119. Integer-valued polynomials and a different comparison problem (pp. 500--505; 134-7/58--63).}
The entry dated 26.9 returns to the failed comparison of discrete and schematic linearization. Grothendieck reports Illusie's objection from Breen's work on stable cohomology, while admitting that the precise counterexample has not yet been made explicit to him (p.~500). He then gives a concrete degree-two obstruction using central extensions of the additive group (p.~501). The multiplicative group of truncated series $1+at+bt^2$ defines the extension $W_2$; a polynomial splitting requires $P(a+a')=P(a)+P(a')+aa'$. The solutions displayed over $\underline Q$ involve $t(t-1)/2$, which has no integral-coefficient representative after addition of the displayed linear term. This separates polynomial cochains over $\underline Z$ from arbitrary cochains on the discrete group. The final splitting display literally has $t$, without the square appearing in the preceding general homomorphism display; the checked text preserves that visible slip rather than adopting the arXiv witness's $t^2$. The main argument should therefore be read through the preceding coefficient equation, not by treating the last display as an independently correct formula. The next proposal replaces polynomial algebras by ``binomial algebras,'' built from the integer-valued expressions $P_n(t)=t(t-1)\cdots(t-n+1)/n!$ (p.~502). Grothendieck tentatively requires the ground ring to carry the corresponding binomial operations and compares this enrichment with divided powers in crystalline cohomology. The marginal reference to pp.~506--507 anticipates the next section's reconsideration of these structures and of the restrictions on the ground ring. On p.~503 the finite-type qualification on the free module $M$ is an above-line typed insertion. The algebra $\mathrm{Symbin}_{\underline Z}(M)$ is defined through rational polynomial functions integral-valued on the dual lattice, and the map to all functions on a semisimplicial free abelian group is posed as a possible quasi-isomorphism in the connected range. Reduction to $K(\underline Z,n)$ is proposed by d\'evissage, not proved. The broader question asks whether Eilenberg--Maclane homology or cohomology is derived from a suitable nonadditive functor $B$, leading to ``B-schematic homotopy types'' and unipotent $B$-bundles. The dark superscript on the first Symbin display remains unresolved and is not silently discarded. The author declines to pursue that alternative immediately, because the ordinary schematic program still promises a simple Lie-functor description of homotopy invariants (pp.~503--504). Loops, suspensions, and compatibility with passage to discrete types remain the practical tests. Reading Bousfield--Kan then changes the comparison: their construction is identified as $k$-completion, distinct from a fully developed category of types over $k$. The marginal ``definitely not, in general!'' rejects the tentative identification of completion with extension of the ground ring; the faint underline beneath ``ring extension'' identifies the disputed phrase (p.~504). Page~505 links this reading to the integration and cointegration operations of section~69 and to Don Anderson's existence theorem in closed model categories. A pen transposition moves the closed-model-category qualification directly after ``operations,'' making the setting of that theorem clearer. The section closes by recognizing Bousfield--Kan as foundational reading while postponing it long enough to record the recent schematic reflections. Section~120 starts on 28.9 by abstracting the extra algebraic operations just proposed.

\subsection*{120. Extra algebraic operations and polyhedral models (pp. 506--507; 134-7/64--65).}
The entry dated 28.9 abstracts the previous section's binomial proposal together with divided powers, lambda operations, and Adams-type operations. Grothendieck uses the notion of an ``analyzer'' and asks for an extension of the analyzer of commutative rings in which unary operations determine higher components by tensor powers. The handwritten changes from $n$ to $n+1$ tensor factors and from the older bounds to $n\geq-1$ and $n\geq1$ affect the indexing of the proposal (p.~506). The live corrected bounds are clear, although overwritten earlier digits remain doubtful. The essential datum is described as a ground ring $k_0$ and an algebra $\Omega_0$ equipped with composition. An $\Omega$-structure on a ring is then an action of these operations satisfying compatibility with addition, multiplication, and constants. Page~507 displays the two diagonals that encode the first two conditions. This is an outline of the desired axioms; the author expressly declines to write all of them down. The circled instructions ``\`a la ligne'' and ``id'' on p.~506 are retained as typist instructions, whose precise targets are not sufficiently explicit to justify speculative paragraph restructuring. The second half of p.~507 turns from these operations to David W.~Jones's work on polyhedral cells and poly-T-complexes, supplied by Ronnie Brown. Grothendieck relates it to the questions of sections~91 and~93: when does a category of cells form a weak test category, and when does the proposed standard chain complex compute the homology of its presheaves? He finds that Jones's emphasis is instead on generalizing thin structures and a corresponding equivalence of categories. The original questions therefore remain open for this notebook. Both appearances of the presheaf-category notation have a hand-drawn upward caret in the scan, rather than the downward vee supplied by the witness. Section~121 next returns to the formal role that linearization should play across the various possible schematic approaches.

\subsection*{121. Abelianization, adjunction, and the pointed correction (pp. 508--510; 134-7/66--68).}
Grothendieck begins the formal review by defining $\mathrm{Hotab}=D.(Ab)$ and considering abel\-ian\-iza\-tion $LH.$ together with the Eilenberg--Maclane functor $K_\pi$ (p.~508). He relates their adjunction to the free-abelian-group and forgetful functors on semisimplicial objects, then interprets ordinary cohomology through the representing objects $K(\pi,n)$. The notation is visibly under revision: a blacked-out earlier functor label is replaced by $K_\pi$, a handwritten $K_\pi$ is itself inserted and cancelled before the typed phrase ``defined via,'' and several category-name marks remain doubtful. The checked text does not reconstruct those older labels from the mathematical context. Page~509 explains why $K_\pi$ is conservative using $\pi_i(K_\pi(L.))\simeq H_i(L.)$, while explicitly denying full faithfulness and the assertion that the counit $LH.(K_\pi(L.))\to L.$ is an isomorphism. The typed letter $E$ in the preceding adjunction sentence is retained even though the surrounding notation uses $K_\pi$. The unit $X\to K_\pi(LH.(X))$ leads to the Hurewicz maps and a proposed homotopy-fiber interpretation of the Whitehead exact sequence (p.~509). The reduced homology in that sequence immediately calls for care with base points. Diagrams(8) and(10) factor the construction through pointed objects, using free pointing $X\amalg e$, the forgetful functor, and reduced free abelianization (pp.~509--510). The quotient description is replaced by an exact triangle, and the parenthetical typed ``one hopes!'' qualifies the continuing claim of adjunction. At the end of p.~510 the author explicitly corrects the earlier passage: the sequence and its fiber should come from the pointed unit(6'), since the unpointed expression lacks the canonical base points needed to define the indicated homotopy groups and fibers. The scan's(6') uses a tilde on both $K$ and $LH$, with no $\pi$ subscript on $K$; the witness regularizes this differently. Section~122 then revisits reduced abelianization itself, seeking a definition that no longer requires the chosen point.

\subsection*{122. Reduced abelianization without a point, and the failed homotopy-linearization functor (pp. 511--517; 134-7/69--75).}
The entry dated 2.10 revises the pointed construction of section~121. Instead of taking the quotient by the chosen base point, Grothendieck defines $\widetilde W(X)$ as the kernel of the augmentation $\underline Z^{(X)}\to\underline Z$ (p.~511). The fiber $W(X)(1)$ is the associated torsor and receives the universal map from $X$ to an abelian torsor. A point splits the extension, recovering the earlier reduced construction and giving the direct-sum description of homology; without a chosen point the natural exact triangle is arranged differently from the triangle on p.~510 (p.~512). The comparison is more informative over an object $Y$, because a section need not exist and the torsor class $c(X/Y)$ can be nontrivial. This supplies the connection with Postnikov invariants and motivates a more general Kan--Dold--Puppe correspondence for coefficients in the relevant slice topos (p.~513). The generalization to local test categories is announced as an intuition to revisit, not established here. The paper's notation is also simplified: the old $K_\pi$ and its pointed variant become $K$ and $\widetilde K$. A proposed partner $S$ or $\widetilde S$ would realize prescribed reduced homology, in contrast with the homotopy groups prescribed by $\widetilde K$ (pp.~514--515). The coefficient letter in the sphere-like objects $S(H,n)$ is changed in pen from $\pi$ to $H$, and circled marginal $H$ notes emphasize the change. Restriction to connected complexes addresses the immediate obstruction that degree-zero homology of a space is free. The author admits that he has neither constructed the desired functor nor established its existence or uniqueness. In the schematic setting even the relevant low-dimensional sphere objects have not yet been constructed. The possible value of this program depends on the tentative Lie-functor description from section~118, not on an available computation of sphere homotopy groups. The proposed fourth functor $L\pi.$ is more sharply rejected. It would turn homotopy groups into homology of a chain complex and satisfy $L\pi.(\widetilde K(L.))\simeq L.$. The continuation dated 3.10 states that these properties cannot all hold (p.~515). Applying the functor to the abelianization unit would provide a canonical section of the counit. On $L.=\pi[n]$, this would split the universal-coefficient sequence naturally in the groups involved (p.~516). The corrected indices matter: formula(18) has $\mathrm{Ext}^1$, $H^{n+1}$ and $H_{n+1}$, where the witness gives different indices. The text then specializes to degree one over $\underline F_2$ and the sequence $0\to V'\to\mathrm{Sym}^2(V')\to\bigwedge^2V'\to0$ (p.~517). An editorial mathematical qualification is needed at the last step. The manuscript claims equivariant nonsplitting for $\dim V\geq2$; the literal bound is retained. If $\mathrm{Aut}(V)$ means the finite group $\mathrm{Aut}_{\underline F_2}(V)$, as the preceding discussion of automorphisms of the abstract abelian group suggests, dimension two is an exception: $q=x^2+xy+y^2$ is fixed by all six invertible linear substitutions over $\underline F_2$, while its polar form maps to the generator of $\bigwedge^2V'$. It therefore supplies an equivariant section for that finite-group action. The qualification is specific to this interpretation: the calculation does not give equivariance for the algebraic group scheme $\mathrm{GL}(V)$ or for all matrices after scalar extension. This small calculation identifies a limitation of the asserted bound on the finite-group reading and does not repair the manuscript's text. The ensuing remark is conditional in a second way: it assumes that an analogous nonsplitting example can be found for $n\geq2$ before drawing the conclusion for simply connected homotopy types. Grothendieck presents that expected example as an alternative route to the failure of the earlier discrete/schematic comparison, avoiding reliance on the reported Breen results. Section~123 opens by acknowledging that, even after these discarded expectations, the basic schematic homotopy operations and the correct weak equivalences are still not in hand.

\subsection*{123. Sphere complexes and scalar operations as tests of the proposed theory (pp. 518--520; 134-7/76--78).}
The entry dated 4.10 takes stock after the preceding counterarguments. Even the equivalence of the proposed criteria for weak equivalence, and the construction of fibers and cofibers, are not yet settled (p.~518). Grothendieck therefore treats the realization of the ordinary two-sphere over $\underline Z$ as a basic test. Suspension should then produce the higher spheres; a larger category admitting the connected complexes needed for $S(1,k)$ might supply an alternative starting point (pp.~518--519). These are requirements for a workable theory, not constructions already obtained. If schematic sphere objects and the Lie description of their homotopy invariants exist, their homotopy groups should be the homology of canonical complexes $L\pi.(S^n_{\underline Z})$. Maps between spheres would then induce maps of those complexes, compatibly with composition (p.~519). The qualification from section~122 is crucial: the author has just rejected such a functor on arbitrary ordinary homotopy types. He now asks for a special structure on the sphere subcategory and explicitly allows the possibility that this structure is impossible, which would refute the intended schematic theory. The notation is still rough: the defining display for ${}^n\Pi.$ literally writes $L\,.(S^n)$, omitting the $\pi$ present in the surrounding discussion. The checked text retains that omission and the hand-corrected exponent and punctuation in the preceding displays. The first operation to examine is the action of the multiplicative monoid of integers, obtained by taking maps $S^n\to S^n$ (p.~520). The author relates the resulting actions on finite homotopy groups to the profinite monoid $\widehat{\underline Z}^{\mathrm{mult}}$ and proposes its lift to the derived base extension ${}^n\Pi.\otimes^{L}_{\underline Z}\widehat{\underline Z}$. The derived tensor symbol in the scan is not a direct-sum symbol, and no parentheses surround the tensor expression. For odd-dimensional spheres, restriction to the units is interpreted through the cyclotomic Galois action on the profinite homotopy type of punctured affine space. The pen changes from ``maximal'' to ``largest'' subgroup, adds ``(complex),'' and brackets the exponent $2m-1$; the above-line ``profinite'' remains typed. The point of the passage is the additional structure carried by these operations and a hoped-for lift to complexes, rather than a calculation of the sphere groups. Section~124 next separates the distinctive features of this schematic approach from the two alternatives considered earlier.

\subsection*{124. Four functors and a provisional formal calculus (pp. 521--530; 134-7/79--88).}
Grothendieck contrasts the present unipotent-bundle models with the earlier de Rham and diagonal-category approaches of sections~94 and~109. All seek a homology functor, but the present construction also paraphrases the Eilenberg--Maclane functor $\widetilde K$ compatibly with extension of scalars (p.~521). Simply forgetting a chain complex's $k$-module structure and taking a de Rham complex would not supply that compatibility. Moreover, the present schematic types map to ordinary types, whereas the earlier approaches had the reverse direction of comparison. The notation $\mathrm{Hotab}_0(k)$ deliberately leaves its algebraic realization open: section~115 required complexes of proobjects, and taking their inverse limits would lose the asserted base-change behavior. Consequently $\widetilde K$ may require pro-unipotent models on the nonlinear side as well. Pointed linearization offers a different formal theory, whose invariants section~117 has already distinguished from schematic ones (p.~522). The continuation dated 5.10 organizes four expected functors between $\mathrm{Hotab}_0(k)$ and $\mathrm{Hot}_0(k)$: total homology $L\widetilde H.$, total homotopy $L\pi.$, and their partners $\widetilde K$ and $\widetilde S$ (p.~523). Even the zero-connected convention is provisional: the author lacks an intrinsic satisfactory formulation of abelian fundamental group with trivial higher action, and briefly considers restricting to simply connected objects. Property~A posits the adjoint pairs $(L\widetilde H.,\widetilde K)$ and $(\widetilde S,L\pi.)$. Their units and counits give linearization and colinearization. The blank after $L$ in formula(5)b remains unfilled in the scan; the checked text does not silently add $\pi$. The Hom formulas interpret $K(\pi,n)$ as representing cohomology and $S(H,n)$ as representing mixed homotopy with a ``co-coefficient'' $H$. Taking $H=k$ recovers the expected ordinary homotopy-module functor (p.~524). This is the stated motivation for the spherical adjunction, whose existence has not been established. The circled marginal $H$ and $\widetilde W(X)$ record changes to the coefficient and linearization notation. Property~B gives the inversion identities $L\pi.(\widetilde K(L.))\simeq L.$ and $L\widetilde H.(\widetilde S(L.))\simeq L.$ (p.~525). Property~C obtains a Hurewicz map by either applying total homotopy to the linearization map or total homology to the colinearization map. The desired compatibility is that these constructions agree. The first nonvanishing homology objects should occur in the same dimension, with an isomorphism there; a pen addition also demands an epimorphism in the next dimension. The lengthy reduction to Eilenberg--Maclane homology in the preceding paragraph is present in the typescript even though the witness omits it. Property~D states preservation of cofibration sequences by the left adjoints and of fibration sequences by the right adjoints. The two exact triangles spell out the homology and homotopy cases (pp.~525--526). In the second triangle the scan's corrected top vertex is $L\Pi.(Y)$ and its bottom row is $L\Pi.(Z)\to L\Pi.(X)$; the witness gives a different cyclic arrangement. Both oblique labels are the numeral~1. The exactness discussion reveals an unwanted asymmetry. Suspension on the connected linear category is a shift and is fully faithful, whereas loops require truncation and lose information (p.~526). The first truncation bound has been corrected from~0 to~1. Remark~1 explains that the loop and fiber constructions therefore mean the neutral or marked connected component of the unrestricted construction. Grothendieck says that he has only just become aware of this restriction while writing, questions the practical value of imposing $X_0=e$, and contemplates enlarging both categories (p.~527). The wished-for trivial action of $\pi_1$ on higher groups is a different condition and is not supplied by $X_0=e$. Remark~2 recalls the Leray spectral sequence for homology of fibrations and asks whether a corresponding mechanism exists for homotopy of cofibers; it does not supply such a mechanism. Remark~3 starts from the evident exchange symmetry of the four-functor formulas, but then distinguishes autodual axioms from an actual autoduality of the categories (pp.~527--529). Full faithfulness of suspension must be excluded from the symmetric axioms. Derived duality on perfect complexes sends chains to cochains, prompting a tentative appeal to stabilization. The subsequent recollection of geometric duality is explicitly qualified by ``as far as I remember'' and ``if I got it right''; the former qualification, the word ``discovered,'' and the shift qualification are typed insertions, placed with pen brackets (p.~529). That discussion pertains to homology and still does not produce a duality exchanging homology and homotopy. The literal sentence ``homotopy type of X'' omits the vee used on the complement elsewhere; the checked text preserves the omission. Properties~E--G return to the operational requirements (pp.~529--530). Homology and homotopy should be conservative in the indicated connected setting, and all four functors together with their structural maps should commute with extension of the ground ring. Restriction of scalars is mentioned without a developed list of formal properties. Comparison with ordinary types combines forgetting $k$ on the linear side with taking sections of schematic models. It should respect $\widetilde K$, but the already disproved homology comparison is explicitly excluded, even for $k=\underline Z$. Compatibility with $\widetilde S$ is still hoped for over $\underline Z$, conditional on constructing the discrete spherical functor. There is no ordinary total-homotopy functor to compare with $L\pi.$; only the separate homotopy invariants are available there. The closing expectation that the comparison respects fiber and cofiber structures is another requirement of the program. Section~125 begins on13.10 with Baues's proposed axiomatic setting for making those operations precise.

\subsection*{125. From Baues's categories to weak equivalences and derivators (pp. 530--532; 134-7/88--90).}
The entry dated13.10 records a letter from H.~J.~Baues and a reading of his preprint on the homotopy classification problem. Grothendieck recognizes the relevance of the dual notions of cofibration and fibration category to his integration and cointegration operations (p.~530). The present semisimplicial unipotent models, however, do not admit the pushouts he would need for a cofibration structure, except in the stated trivial case. He has begun to consider a Kan-fibration notion but reports no definite conclusion. He therefore proposes to leave the schematic construction temporarily and return after further work on these categorical methods (p.~531). The more general objection concerns the scope of the axioms. Grothendieck questions whether every meaningful homotopy universe must arise from Quillen's or Baues's extra structures on a chosen category of models. He invokes semisimplicial objects in a topos, Illusie's weak equivalences, and Verdier's derived categories as reasons to seek a wider formulation (p.~531). These are the author's contemporary judgments and questions, not proofs of general nonexistence of model structures. The qualifications ``it seems,'' ``I would think,'' and ``but (it seems)'' are significant: the last is a pen revision to the claim about enough projectives or injectives and bounded derived categories (p.~532). The empty reference brackets after Illusie and the marginal cross remain empty in the text. The circled $C^{-}(A)$ and $D^{-}(A)$ correct the signs used in the discussion of the bounded categories. The proposed replacement starts with a category $M$ and weak equivalences $W$, without treating a selected fibration or cofibration class as the fundamental datum (p.~532). Different compatible choices should yield the same integration and cointegration operations. Section~69 supplies the earlier point of reference; a derivator is now identified as the structure that should express those operations independently of a particular model presentation. Constructing and characterizing the schematic derivator remain hopes. The passage also frames the derivator as the further structure missing when one stops at a localized or triangulated category, in both abelian and nonabelian settings. The typed insertions ``specific,'' ``on it,'' and ``make any mention of models'' sharpen that distinction, while the correction from ``is'' to the typed ``are'' refers to the two operations. Section~126 resumes the schematic discussion by testing the formal symmetry against Postnikov decomposition.

\subsection*{126. Postnikov truncation and the failure of a literal dual filtration (pp. 533--535; 134-7/91--93).}
The section tests the four-functor symmetry on the construction that originally motivated the schematic program. An $n$-co-connected object has no homotopy groups above degree $n$, and the proposed truncation $X\to X_n$ is the universal map to that full subcategory (p.~533). Equivalently, its inclusion should admit a left adjoint. This universal property produces the familiar inverse tower, with $X$ viewed as an appropriate inverse limit (p.~534). The discussion is axiomatic: it asks what should follow from the basic functors and their structural maps. The double underline on ``co'' distinguishes the upper-degree vanishing condition from connectedness. A literal dual would be a canonical increasing filtration ${}_nX\to X$, inducing the required homology isomorphisms through degree $n$ and vanishing above that degree. The ordinary skeletal filtration misses the endpoint of the desired isomorphism range. After trying to modify it, Grothendieck reports the feeling that no such canonical filtration exists, while explicitly admitting that he has neither formulated nor proved a nonexistence statement (p.~534). The decisive structural objection is that both proposed truncation subcategories are inverse images of the same linear subcategory under $L\pi.$ or $LH.$. That linear subcategory is stable under loops but not under suspension (p.~535). Even if the doubtful right adjoints existed on the homology side, their successive terms would not have the literally dual Postnikov relationship. A different suspension-stable family is contemplated and then abandoned for lack of a reasonable candidate. The conclusion qualifies the formal symmetry of section~124: homology and homotopy retain materially asymmetric properties. Several visible details prevent the reader from mistaking the formulas for a fully regularized exposition. On p.~533 the object named $Y$ is tested by a formula using $X$; on p.~534 the skeletal vanishing display writes $H_i(X_n)$ with a right subscript despite the surrounding left-subscript filtration. Formula labels are changed by hand, while some references to(5) remain on p.~535 after the relevant subcategory has been renumbered. Those literal features are retained. Section~127 then changes the question from symmetry of operations to the choice of components allowed in a schematic model.

\subsection*{127. Which objects should serve as schematic models? (pp. 535--538; 134-7/93--96).}
The section keeps the usual simplicial indexing category and changes the category of allowed components. The broad ambient setting is the fpqc topos of sheaves over the ground ring, with Illusie's weak equivalences of semisimplicial sheaves (pp.~535--536). Such a weak equivalence need not induce a weak equivalence on ordinary sections. Grothendieck therefore distinguishes the section functor from the canonical cointegration operation that should recover a discrete homotopy type. The above-line words ``canonical'' and ``(= weak equivalence)'' are typed revisions. He hopes that a suitably restricted class of models, especially the earlier unipotent bundles, will make the naive section construction adequate. The handwritten sentence opening p.~537 makes the base-change requirement explicit: a schematic homotopy type over $k$ should give a functor from $k$-algebras to ordinary homotopy types. The relative homotopy sheaves are defined first after choosing a section. The proposed ``pseudo-abelian'' condition combines relative zero-connectedness, abelian fundamental group, and trivial action on the higher homotopy sheaves (p.~537). Under a local formulation of that condition, the abelian sheaves should be independent of the chosen section and exist even without a global section. Their desired quasicoherent module structures are an additional issue. The comparison from the homotopy of sections to sections of the homotopy sheaves should be an isomorphism under mild hypotheses, for example $X_0=X_1=e$; the author expressly says he has not attempted a proof for unipotent bundles. He still calls the module structure a mystery without the Lie functor of section~118. The final remark proposes a spectral sequence from cohomology of the homotopy sheaves to the homotopy of cointegration, suggesting that acyclicity might imply the comparison (pp.~537--538). This remains a suggested argument, marked by ``maybe'' and ``If so,'' rather than a constructed spectral sequence. Section~128 next imposes a separate closure requirement: linearization should itself remain in the chosen category of bundles.

\subsection*{128. Linearization and the competing $W$- and $V$-bundles (pp. 538--541; 134-7/96--99).}
The new requirement is that linearization of an allowed bundle should itself be an allowed bundle, carrying the additional structure of an $\underline O_k$-module (p.~538). A forgetful functor $K:Uab(k)\to U(k)$ should therefore have a left adjoint $L$. The two natural choices have opposite variance: $W(M)(k')=M\otimes_k k'$ gives quasicoherent modules, while $V(M)(k')=\operatorname{Hom}_k(M,k')$ gives the opposite of the category of modules. The typed wording does not introduce an extra $X$ after ``every bundle'' or ``for a given bundle,'' despite the witness's additions, and formula(7) has a semicolon after $K$. On passing to semisimplicial objects, these functors are intended to produce the earlier $\widetilde K$ and $LH.$. The final dot of $LH.$ on p.~539 is baseline punctuation in the source, not a supplied subscript. Section~115 had already shown that the $W$-case gives only a proadjoint. Grothendieck now draws out the cost: one must work with prounipotent bundles and semisimplicial prosheaves outside the sheaf-theoretic foundations he is using (p.~539). A diagonal marginal addition offers working with the inverse limit of the corresponding sheaves; its wording is ``we may however,'' and ``corr.'' remains abbreviated. The witness omits this qualification. The $V$-case is more tractable because $V(M)=\operatorname{Spec}(\operatorname{Sym}_kM)$. Pullback--pushforward adjunction and quasicoherence of the direct image for a quasicompact, quasiseparated scheme give $L_V(X)=V(\Gamma(X,\underline O_X))$, retaining only the underlying module of the global function algebra (pp.~539--540). The Hom subscript in the first display on p.~540 is $\underline O_S$, whereas the witness prints $\underline O_X$. Two circled marginal formulas emphasize the global-sections description. The replacement of ``inclusion'' by the typed ``forgetful'' makes the functor under discussion precise. This advantage raises a larger question: why restrict components to schemes isomorphic to vector bundles instead of allowing all affine schemes? In the affine setting, semisimplicial schemes correspond to cosemisimplicial commutative algebras, and linearization simply forgets multiplication (p.~540). The answer returns to the original ambition over $\underline Z$, where the schematic category should remain close to ordinary homotopy types. Grothendieck calls the initial choice of $W$ rather than $V$ almost a matter of chance, but formulates two substantive tests for either choice (p.~541). First, the homotopy-sheaf comparison with sections should hold, possibly under $X_0=e$ or $X_0=X_1=e$. Second, under flatness assumptions, the Lie functor should describe the homotopy sheaves. Postnikov d\'evissage supplies only the tentative special case, without the required naturality. He regards recovery from a first-order neighborhood as plausible for affine spaces and implausible for general affine schemes, citing an elliptic curve minus a point. Thus these tests motivate the restricted components; they are not established by the easier linearization adjunction. The next day's section~129 adds the further problem of descending that adjunction to localized categories.

\subsection*{129. The missing descent of linearization and q.c.h. equivalences (pp. 542--543; 134-7/100--101).}
The entry dated14.10 corrects an omission in the preceding day's argument: existence of $L:U(k)\to Uab(k)$ does not itself imply a functor on the corresponding localized categories. One must specify weak equivalences on the linear objects and prove that $L$ preserves them (p.~542). Grothendieck expects preservation and reflection for the $W$- and $V$-models under suitable connectedness assumptions, but explicitly says that a proof has not even been attempted in the unipotent case. This is the third test, after the two tests on p.~541. He compares it with the Whitehead problem already considered in section~92 for replacing the simplex category by other indexing categories. The first live initial before Whitehead in the scan is a corrected typed $J$, over a doubtful older impression. The raised qualifications ``maps between'' and ``total'' are typed additions, and the filled versus open round indices remain graphically distinguished. For broader classes of semisimplicial schemes, the section proposes defining ``quasi-coherent homological quasi-isomorphisms'' by their image under linearization (p.~543). Ordinary weak equivalence need not yet be known to imply this condition, and a homological criterion can be weaker without connectedness restrictions. The proposed refinement allows twisted quasicoherent coefficients: a sheaf $F_*$ must be carried isomorphically to itself by the structural pullback maps. With a section and relative zero-connectedness, this is related to a fundamental-group-sheaf action on coefficients from the base. The circled marginal abbreviation q.c.h. is present in the manuscript and absent from the witness. Grothendieck then hopes that localization at these maps could make a mild category of general semisimplicial schemes equivalent to the category obtained from bundles, producing a subtler connection with ordinary homotopy types than the known \`etale pro-homotopy connection. He immediately suspends the suggestion as possibly out of reach or mistaken. The typed insertions ``very mild'' and ``only'' emphasize the proposed enlargement and comparison, not a proved equivalence. Section~130 returns to a more limited enlargement by allowing disconnected bundles.

\subsection*{130. Disconnected components as a possible aid to integration (pp. 544--545; 134-7/102--103).}
The limited enlargement now proposed allows direct sums of unipotent bundles as the components of a semisimplicial model. Deleting ``finite'' and replacing the following sentence make the role of infinite sums explicit: linearization still makes sense if the linear category is closed under them; otherwise one must retain a finiteness restriction (p.~544). Descent would also call for twisted finite sums, but the author postpones those technicalities. Applying connected components degree by degree produces a discrete semisimplicial object $E_* = \pi_{o/k}(X_*)$ and a map $X_*\to E_*$. Here $\operatorname{Spec}(k)$ is assumed connected and the sums untwisted. The ``strict'' schematic part lies over this discrete homotopy type, whose own invariants are not themselves schematic. To retain the original strict theory, $E_*$ should be aspheric, explicitly meaning isomorphic to a point in $\mathrm{Hot}$. Equivalence of the resulting localized category with the earlier connected-model category is still only presumed. The motivation is operational: the ordinary construction of integration from section~69 introduces many connected components, and standard simplices would become available among the models (p.~545). For an aspheric indexing category, Grothendieck hopes to preserve the asphericity of $E_*$. He proposes performing integration and cointegration first among semisimplicial sheaves, then checking whether the results lie in the restricted bundle class. He reports no clear-cut result. Moreover, disconnected components threaten the tangent-sheaf description of section~118, so the enlargement would be justified only after proving that it gives the same localized category as the connected models; no attempt at that proof has yet been made. The typed insertions ``components of the,'' ``sheaf,'' and ``would'' refine this cautious formulation. The cancelled prospective number131 belongs to the continued warning, while the actual section131 starts below a literal date16.9. That date is retained despite the preceding entry dated14.10. Its opening personal interlude records the kimchi instructions and a letter to Larry Breen before the next page sharpens the mathematical objection to the enlargement.

\subsection*{131. Infinitesimal spheres and the failure of the naive Lie formula (pp. 545--551; 134-7/103--109).}
After the personal interlude opening the section, Grothendieck rejects the disconnected-model proposal as largely a symptom of his inability to construct the required ``left'' homotopy operations within connected unipotent bundles (pp.~545--546). The implicit condition $X_o=e$, or even $X_o=X_1=e$, is essential to the expected Lie description. Without it, tangent invariants at different sections need not be canonically identified, already for a constant semisimplicial bundle. Reimposing it would undermine the ordinary integration construction that motivated disconnected components in the first place. The author therefore distinguishes the restricted models for schematic homotopy types from more general semisimplicial sheaves that are useful auxiliary objects. The added pen $o$ in ``loosing track'' is clear and is retained, although the witness prints ``losing.'' A different enlargement is suggested by the still hypothetical spherical functor $\widetilde S$ of section~124, which should be left adjoint to total homotopy and right inverse to reduced total homology (pp.~547--548). For flat augmented symmetric algebras representing pointed $V$-bundles, the augmentation ideal modulo its square gives the co-Lie module. A module $M$ also defines the square-zero algebra $D(M)=k\oplus M$ and its spectrum $I(M)$. Applied degreewise, this construction gives the formal candidate $\widetilde S(L_*)=I(L^*)$ for $L_*=V(L^*)$. The resulting Hom identity expresses ordinary representability of the tangent functor by first-order infinitesimal schemes. Grothendieck explicitly marks two gaps: it is an identity between model objects, not yet an adjunction of derived categories, and the resulting infinitesimal schemes are outside the existing bundle category. The pen correction of ``adjunction'' and the deleted claim that the first Hom is taken in $\mathrm{Hotab}$ make the first gap especially visible on p.~548. Page~547 also literally introduces ``L'' as the arbitrary module immediately before writing $D(M)$; that variable mismatch is retained. The inversion identity with reduced homology is equally tautological for square-zero algebras, motivating a new definition of bundles as nilpotent thickenings of $V$-bundles (p.~549). The chosen nilpotent ideal is not part of the structure, and further flatness conditions remain possible. The subscheme in the displayed inclusion is $X_o$, whereas the witness prints $X_*$. Three of the functors now have simple model descriptions: augmentation ideals for reduced homology, symmetric algebras for $\widetilde K$, and their first-order truncations for $\widetilde S$. The canonical closed immersion $\widetilde S(L_*)\to\widetilde K(L_*)$ reveals a missing compatibility axiom in section~124 (p.~550). Its two adjuncts should be the inversion isomorphisms, and the maps obtained from those two isomorphisms must agree. Taking this natural transformation as primary data would also give a single Hurewicz map. The literal condition near the foot of p.~550 has right-hand side $k$, with a doubtful round left index; the witness replaces it by $X_0=e$. No such mathematical correction is made in the text. The apparent simplicity breaks down at the interpretation of total homotopy (p.~551). On the algebra side, naive $L\Pi.$ takes an augmentation ideal to its quotient by the square; dualizing gives the inclusion form of the Hurewicz map. But for a first-order thickening the square is already zero, so $L\Pi.(\widetilde S(L_*))\simeq L_*$. If these were the desired spheres, their naive Lie homotopy would therefore equal their homology, which the author emphatically rejects. His remaining hope is restricted to actual flat vector-bundle components, the basepoint condition, and possibly unspecified Kan hypotheses. The triple exclamation and the two marginal corner marks emphasize the failure and the beginning of section~132. That next section replaces the naive tangent functor by the Andr\'e--Quillen--Illusie tangent and cotangent complexes.

\subsection*{132. Derived tangent complexes and the wrong-quadrant obstruction (pp. 551--554; 134-7/109--112).}
The failure of the naive Lie functor in section~131 prompts replacement of ordinary differentials and tangent spaces by the Andr\'e--Quillen--Illusie complexes. Grothendieck emphasizes that the naive objects behave as the desired homological objects only in the smooth case, which here means the flat vector-bundle components under consideration (pp.~551--552). Formula(15) defines the superscript complex as the derived dual of the subscript complex, and formula(16) gives the corresponding co-Lie and Lie objects along the marked section by derived pullback. The prose first lists ``tangent or cotangent'' before those two notations; that visible order is retained, while the displayed duality and the subsequent co-Lie/Lie terminology specify the roles used in the argument. A free semisimplicial model of the co-Lie complex, followed by its dual, leads to the proposed formula(18) for total homotopy. Its handwritten question mark is essential: the formula is a candidate, not a theorem. Applying this construction to each component of a semisimplicial model introduces two indices of opposite variance. Under Kan--Dold--Puppe it becomes a bicomplex in the quadrant $n\leq0$, $p\geq0$ (p.~553). Taking its associated simple complex is tempting, but each total degree involves infinitely many terms. The manuscript therefore leaves open whether totalization should use a direct sum or a direct product. It also observes that the result can have both positive and negative degrees, so it need not lie in the desired chain-complex subcategory of the derived category. Truncation is proposed as a possible further correction if positive cohomology does not vanish. A separate difficulty is the nonquasicoherence introduced by duals without finiteness assumptions. Returning to the cotangent rather than tangent complexes would address that motivation but still leaves a mixed complex in a wrong quadrant (pp.~553--554). The proposed sphere $S(n,k)=\widetilde S(k[n])$ would be a concrete test of(18), especially over $\underline Z$, where its homotopy modules are hoped to recover the ordinary sphere groups. The author has not begun the necessary computation of the Andr\'e--Quillen Lie complex of the square-zero algebra $D(M)$ (p.~554). More generally, even applying the earlier flat-vector-bundle formula requires replacements of arbitrary models by models with flat components. Polynomial resolutions of the associated augmented algebras again introduce a second index of the opposite variance. The closing question is therefore whether these mixed wrong-quadrant objects could themselves form a workable category of models. Grothendieck expressly admits that his formal intuition breaks down and that the suggestion may be meaningless. The stepped double pen stroke before this question is a paragraph-separation cue whose upper arm passes beneath ``spectra''; it is not an independent underline emphasizing that word. The next supplied page is substantially masked, so the visible mathematical argument ends here without a resolution.

\subsection*{133 (witness numbering). A masked transition and the visible seminar interlude (pp. 555--557; 134-8/1--3).}
The supplied images of p.~555 expose the date22.10. and page number but cover the entire body with an opaque rectangle. The upper body of p.~556 is likewise masked, and two smaller rectangles hide the end of the personal passage on p.~557. The separate arXiv witness supplies a section133 heading and narrative in these positions, but neither that hidden wording nor the heading can be visually collated from the supplied scans. The checked text therefore retains only visible matter and records the masks as limits of the evidence; it does not silently fill them from the witness. The visible continuation concerns teaching for the concours d'agr\'egation and a small seminar on Teichm\"uller groupoids (p.~556). The typed revisions replace ``get involved in'' by ``start'' and ``great'' by ``special,'' while the pen addition describes the relaxed and friendly atmosphere of the teaching session. Grothendieck anticipates learning mathematics again through teaching, but worries that the renewed fascination of the Teichm\"uller structures will compete with completion of the homotopical-algebra notes. The continuation on p.~557 explicitly compares that fascination with motives; a pencilled marginal qualification reads ``it's very close anyhow.'' The visible text then breaks at the masks. Section~134 begins below them with the literal typed date26.8., whereas the witness changes it to26.10. The guide and checked text preserve the source date and make no inference about the hidden transition.

\subsection*{134. Linearization through integrators and two stronger comparison questions (pp. 557--560; 134-8/3--6).}
The return to ordinary homotopy types contrasts the earlier modelizer $(Cat)$ with the schematic discussion, where semisimplicial models and Kan--Dold--Puppe had dominated. For a general category $A$ and its terminal map $p$, the left adjoint $p^{ab}_{!}$ to constant abelian presheaves has a left derived functor $Lp^{ab}_{!}$, interpreted as total homology (pp.~557--558). Restricting it to individual abelian presheaves can be computed by pairing the presheaf with an integrator $L.^B$ for the opposite category $B=A^o$, then passing from chain complexes to the derived category. The text's later description on p.~558 calls this an integrator for $A$, or cointegrator for $B$; the varying terminology is retained, and the explicit pairing and variance supply the reliable description. Abelian weak equivalences are defined by this homology functor, producing the localization $\mathrm{Hotab}_A$. They must not be confused with quasi-isomorphisms of complexes of presheaves, which belong to a different input category. Even for $A=\Delta$, the unnormalized chain-complex functor in diagram(7) is not the Kan--Dold--Puppe equivalence: normalization is still required (p.~559). Grothendieck therefore places the hoped-for general equivalence at the lower, localized level $\mathrm{Hotab}_A\to D.(Ab)$. He suspects this for test categories, possibly even weak test categories. The first additional question asks whether a map of abelian presheaves is an abelian weak equivalence exactly when its underlying map is an ordinary weak equivalence. The proposed route is a functorial identification of the homotopy groups of the pointed category $A_{/F}$ with $H_i(A,F)$. The marginal note weakens the assumed final object in $A$ to1-connectedness; it is absent from the witness. The author explicitly revises the convention of section~92: if the two classes of maps disagree, the homology-defined localization is now regarded as the correct $\mathrm{Hotab}_A$. Neither their equality for general $A$ nor the proposed homotopy--homology identifications are established here. The second additional question concerns homotopisms (p.~560). In the simplicial abelian case, the Kan property relates ordinary weak equivalence to ordinary homotopy equivalence. For objects with projective module values, Kan--Dold--Puppe also identifies weak equivalences with homotopisms in the stricter abelian sense and identifies equality in the localized category with strict abelian homotopy of maps. The projectivity condition belongs to this stronger statement. Grothendieck recalls the analogous formulations over a ground ring and in an abelian category, then asks for suitable interval structures on a general presheaf category. Sections~51,~54,~79,~95 and~97 provide the dependencies: interval and contractibility structures, the weak-equivalence construction, and contractors. He hopes that the three earlier interval structures coincide for a contractor and support these extra features; for a mere test category he suggests the structure defined from $W_A$. The qualifications ``it would seem,'' ``we may hope,'' and ``possibly'' prevent these proposed extensions from being read as results. The two marginal corner strokes and handwritten135 separate the next section's attempt to recover an earlier abelianization proof.

\subsection*{135. Why integrators give the absolute Whitehead functor (pp. 560--564; 134-8/6--10).}
Grothendieck sets aside the stronger comparison questions of section~134, suggesting that closed model structures belong in a later chapter, and undertakes the promised result from sections~99--103: integrators are abelianizators (pp.~560--561). He first defines the absolute Whitehead functor without simplicial machinery. For a category $A$, total homology is $LH.(A)=Lp^{ab}_{A!}(\underline Z_A)$, and the desired functor on $\mathrm{Hot}$ follows once this construction is functorial and inverts weak equivalences. The handwritten parenthesis adds the name ``absolute Whitehead functor''; the opening ``First of all'' is instead a typed interlinear revision. For general coefficients the morphism data are a category map $f:A\to A'$ and a coefficient map to the pullback of $F'$. Formula(4) literally prints $u^*(F')$ after introducing $f$ as the category map. That visible variable mismatch is retained; the adjunction argument on the next page uses the pullback corresponding to $f$. Functoriality follows by composing the canonical augmentation $Lf^{ab}_{!}(F)\to f^{ab}_{!}(F)$ with the map to $F'$ adjunct to the coefficient map (p.~562). Composition of left adjoints gives the comparison of total homology, and preservation of projectives permits the derived composition used there. To show invariance under weak equivalence, the author tests the induced map by derived Hom into bounded-below complexes and identifies the result with the corresponding map in cohomology with constant coefficients. He then invokes the earlier definition of weak equivalence in $(Cat)$ through cohomology. The proof is presented as an argument, unlike the speculative extensions immediately preceding it. Several formula underlayers remain unreadable, but the supposed extra comma after $K_{A'}$ is not a separate source mark: the scan has the prime belonging to the subscript, which the initial transcription counted twice. The remaining comparison is the natural isomorphism $LH.(A_{/X},\underline Z)\simeq LH.(A,\underline Z^{(X)})$ (p.~563). With $f:A_{/X}\to A$, it follows from $Lf^{ab}_{!}(\underline Z_{A_{/X}})\simeq\underline Z^{(X)}$. The ordinary direct image of the constant presheaf already gives the right-hand side, and the key strengthening is that this particular $f^{ab}_{!}$ is exact. Page~564 gives the reason: $f$ is a fibration with discrete fibers, and its abelian left adjoint is the direct sum, over $u\in X(a)$, of the corresponding coefficient fibers. The formula is left to the reader to verify, after which exactness supplies the vanishing of all positive derived functors. Thus the square relating presheaf abelianization, the category of elements, total homology and the absolute Whitehead functor commutes canonically. This result does not by itself settle the stronger equivalence and homotopy questions of section~134. The closing reflection explains the delay as the author's habit of recasting homology in cohomological terms; the direct-sum argument finally provides the straightforward homological proof he had postponed.

\subsection*{136. Intrinsic homology and the corrected duality analogy (pp. 565--570; 134-8/11--16).}
The entry dated27.10 rereads the abelianization arguments of sections~92 and~100--103 from the viewpoint of the unusual direct image. Grothendieck now regards the initial detour through simplicial models as an obstacle to recognizing the intrinsic construction $Lp^{ab}_{!}(\underline Z_A)$ (pp.~565--566). The comparison proved in section~135 supplies the missing compatibility with presheaf abelianization; for $A=\Delta$ it recovers the description through Kan--Dold--Puppe and localization. The new interlinear restriction to $\widehat\Delta_{ab}$ makes clear which domain is being compared with the classical functor. He draws a methodological conclusion: use the description native to the chosen model category when proving formal properties. The discussion of Thomason's homotopy-colimit construction is another example, contrasting the total category of a fibered category with the Bousfield--Kan presentation. This is the author's assessment, accompanied by an explicit promise to check the latter book, rather than a completed comparison proof. The references to page569 on pp.~565--566 are the literal source numbers; the witness instead supplies359 and369 respectively. The compatibility argument extends to $k$-linearization for an arbitrary ring $k$, explicitly without assuming commutativity (p.~567). Formula(1) proposes obtaining it from the integral case by derived extension of scalars, $Lp^k_{!}(k_A)\simeq Lp^{\underline Z}_{!}(\underline Z_A)\otimes^{L}_{\underline Z}k$. This motivates a formulaire for $Lf_{!},f^*,Rf_*$,including a projection formula and the two internal operations derived tensor product and derived Hom. The analogy with six-operation duality is initially pushed too far: p.~567 calls $f^!$ right adjoint to $Rf_*$, while a pencil note corrects that expression to $Rf_!$. Page~568 nevertheless continues the speculation that suitable finite-type maps in $(Cat)$ might give four successive adjoints $Lf_!,f^*,Rf_*,f^!$. The text states uncertainty about whether there is a deeper relation with geometric duality. It does not construct the proposed extra adjoint. The typed interlinear replacement of ``identity'' by ``relationship'' is one visible softening of the claim. Page~569 explicitly diagnoses the confusion. For a proper map in the geometric duality contexts,$Rf_*$ agrees with $Rf_!$; in general the duality adjunction is $Rf_!\dashv f^!$, alongside $f^*\dashv Rf_*$. Even in the proper case,$Rf_!$ is not generally left adjoint to $f^*$,so the earlier chain of four adjoints does not follow. The pen replacement of a doubtful capital by $R$ in the first displayed pair, and the added explanation that $f^*$ is not isomorphic to $f^!$, belong to this correction. Discrete-fiber fibrations $A_{/X}\to A$ are proposed as the categorical analogue of the exceptional \`etale case. Grothendieck consequently allows that a future categorical duality theory might require a different $Rf_!$ from the $Lf_!$ currently expressing homology. The section therefore preserves an exploratory analogy and its immediate correction, rather than asserting an identification of those operations. The remaining proposal concerns classes of maps and factorization (pp.~569--570). The previously suggested proper and smooth maps are interchanged by passage to opposite categories,while immersions are self-dual. Proper maps include categorical cofibrations and smooth maps include fibrations; the author explicitly distinguishes those cofibrations from the Kan--Quillen notion. He asks how restrictive it is to factor a functor as an immersion followed by a map that is both proper and smooth,perhaps a bifibration. Here an immersion identifies a full subcategory containing every object lying between two of its objects in a composable pair of arrows. No factorization theorem is proved. The closing paragraph sets the larger duality project aside as unnecessary for the present homotopy-model program. Section~137 then immediately revisits the confidence that these definitions of properness and smoothness were the only reasonable ones.

\subsection*{137. Properness tested on stratified realizations (pp. 570--580; 134-8/16--26).}
The immediate reconsideration of section~136 separates two demands that the earlier proper and smooth terminology had conflated. Cohomological base change alone makes every object proper and smooth over the final category, since that projection is a bifibration. Such a weak condition cannot guarantee a right adjoint to direct image: even an infinite discrete category supplies an obstruction (p.~571). Grothendieck therefore asks for compatibility with filtered direct limits and a finite cohomological-dimension condition. Finite groups do not automatically supply the latter, while infinite groups with Poincar\'e-type duality suggest another source of examples. These are tests for a definition, not a completed characterization. The section's frequent cancellations and replacements record the changing scope of the proposed finiteness hypotheses. The provisional six-operation picture assumes a factorization of a suitable map as an immersion followed by a proper map (p.~572). It proposes defining $Rf_!$ by a composite $Rg_*Li_!$, and defining the exceptional inverse image through the corresponding adjoints. The warning added to the page is essential: the notation $Rf_!$ is not meant to assert that this operation is the right derived functor of the already available $f_!$. For a smooth map the proposed comparison involves a rank-one orientation sheaf $T_f$ and a locally constant dimension $d_f$, with $Rf_!$ expressed through $Lf_!$, the inverse orientation twist and the shift $[-d_f]$. Formula(6) supplies the particular relationship recalled at the start of section~138. The factorization independence, the formalism and the smooth comparison remain proposed work. The typed additions ``partial'' and ``here'' qualify the discussion of the earlier confusion; the duplicated definition of $f^!$ is enclosed in a marginal oval. Finite ordered sets provide the testing ground (pp.~573--575). A presheaf on an ordered set can be interpreted as a sheaf constant along the strata of its simplicial realization, with the order reversed relative to containment of strata. Passing to the opposite ordered set does not change the underlying realized homotopy type, but it changes the stratification data relevant to sheaves. This distinguishes the geometric symmetry of realizations from the proposed categorical interchange of properness and smoothness. Finite ordered sets should give proper objects; smoothness is tentatively related to spherical links. For a map of ordered sets, the first proposed condition asks whether every relation $f(x)<y'$ can be lifted to a relation $x<x'$ with $f(x')=y'$. Requiring a unique lift is recognized as too strong. The stronger replacement is that the ordered set of such lifts, written $I(f,x,y')$ in formula(13), should be aspheric (p.~576). Acyclicity may suffice for a version restricted to commutative coefficient groups, but the text does not conflate that weaker condition with asphericity. Formula(12) is meant to compare the relevant upper cones and thereby explain preservation of sheaves compatible with the stratifications. Its printed target is literally $C_y=|I_{\geq y}|$, although $y$ belongs to $J$; the category letter is retained in the transcription. The pencil note ``namely: aspheric fibers!'' sharpens the intended geometric criterion. Page~577 then distinguishes existence of the six operations and their boundedness restrictions from preservation of constructibility or finiteness. The latter is needed for the intended Grothendieck-group calculations but is not made a condition for every formula in the general duality formulaire. The first remark on p.~578 identifies biduality as a place where finiteness is needed. It asks for a dualizing complex and for its realization as $p^!(k_e)$, with uniqueness qualified by an added dimension shift and an invertible coefficient twist. The later running formula writes $p^!(e_k)$ instead; that index reversal is a visible source feature. For finite ordered sets, or locally finite ones of finite combinatorial dimension, the author expects the relevant dualizing statement to follow from the contemplated formalism. He also suggests an intersection-complex construction, while doubting its relevance when the strata are not even-dimensional and when the outcome depends on the stratification. These are explicit qualifications, not claims of a stratification-independent invariant. Two examples show why merely nonempty lifting sets are insufficient (pp.~578--580). First, adjoining a smallest element $x$ to an arbitrary ordered set $I_1$ gives a cone and a map to the two-element ordered set. With $I_o=f^{-1}(0)=\{x\}$, the lifting condition asks only that $I_1$ be nonempty, while the stronger condition asks for asphericity. The text later names1 as the exceptional point fiber and writes an unbarred $I_1$ where a realized fiber is discussed; p.~579 repeats the endpoint labels in the complements $|J|\setminus\{1\}$ and $|J|\setminus\{0\}$. Those labels are not harmonized with the earlier definition of $I_o$. The argument's intended contrast remains clear: a nonaspheric general fiber and a point fiber prevent the higher direct images of constant coefficients from remaining constant on the required stratum. The second example avoids relying solely on that two-stratum interval. Its target has three elements $y\to0$ and $y\to1$, giving the ordinary interval with its two endpoints. In a polygon with $2m\geq4$ vertices, $I_o$ and $I_1$ are the alternating vertex sets. The sentence specifying this polygon is present in the typescript and absent from the witness. The nonempty-lift condition holds, yet the relevant fibers are disconnected: at an endpoint there are $m$ points, away from the endpoints and barycenter there are $m$ disjoint segments, and at the barycenter those segments meet and form a contractible union. Thus even $R^0|f|_*$ of a constant sheaf fails the desired constancy on the open one-dimensional stratum. The closing paragraph proposes that the ordered lift sets have exactly the homotopy types of the geometric fibers of the maps between strata. Asphericity would then express the required contractibility and explain preservation of stratification-compatible complexes. The phrasing remains ``should''; no general comparison proof is supplied before section~138 begins on28.10.

\subsection*{138. Coefficient rings, localization and the existence of $f_!$ (pp. 580--587; 134-8/26--33).}
The entry dated28.10 first enlarges the preceding proposal from a fixed coefficient ring to categories equipped with sheaves of rings (p.~580). Commutativity is required when the internal tensor and Hom operations are considered. The author also insists that the homological operation $Lf_!$ is an additional operation whose relation with geometric $Rf_!$ must be understood; the smooth comparison on p.~572 is recalled as one such relation. The pencil exclamation and cross emphasize the addition. Yet the claim that $Lf_!$ exists for every map of ringed categories is explicitly withdrawn at the end of the section. The opening recollection on p.~581 itself writes $Rf_!$ before turning to the ordinary $f_!$ and its left derived functor. That visible notation is retained rather than silently made uniform. The first obstruction concerns restriction of operator rings (pp.~581--582). For a map of ringed categories with underlying functor $f_o$, the abelian sheaf $f^{ab}_{o!}(F)$ need not visibly acquire the required action of the target sheaf of rings. The source's test case takes $\underline O_A=\underline{End}_{\underline Z}(F)$ and $\underline O_B=f_*(\underline O_A)$. It supplies an action by global sections of $\underline O_B$, but not the desired sheafwise action. If the target ring is constant, that global action yields an ordinary coefficient action. Derived compatibility imposes a further acyclicity requirement on projective source modules, and the example with $B$ final shows that it can fail for a variable source ring. The square-zero extension $\underline Z_A\oplus M$ makes clear how strong universal vanishing would be: it would force homological dimension zero. The scan has a direct sum here; the witness's tensor-product sign is incorrect. A separate literal mismatch on p.~581 writes $f_!(A)$ after introducing the module $F$. For constant source and target rings, the argument reduces to the generators $M^{(a)}$ (p.~582). Resolving $M$ in abelian groups and applying the exact functor $M\mapsto M^{(a)}$ gives a resolution whose direct image is the corresponding complex $L.^{(b)}$, with $b=f_o(a)$. Exactness of $M\mapsto M^{(b)}$ then proves the vanishing in formula(4). Thus the source presents an argument for the constant-ring comparison and only suggests the locally constant extension. The sentence asserting that one cannot go beyond that assumption is cancelled; the interlinear ``rid of such an'' is also separately crossed out. The distinction between a proved calculation and a withdrawn restriction matters for the next discussion. The obstacle is recast as a variance issue (pp.~583--585). Ordinary and derived direct image commute with localization on the base, represented by $B_{/b}\to B$, because the relevant inverse image preserves injectives. The source extends this to discrete-fiber fibrations and then to fibrations more generally. One intermediate display literally reads $B_{/Y}\to Y$; the target letter is retained despite the surrounding discussion of maps to $B$. For $f_!$, the appropriate operation is the opposite construction ${}_{b\backslash}B\to B$, called colocalization. With fixed coefficients, the value of $f_!(F)$ at $b$ is the direct limit over ${}_{b\backslash}A$, dual to the inverse-limit formula for $f_*$. Passing to $Lf_!$ requires preservation of projectives by the corresponding inverse image. The pasted replacement strip on p.~584 supplies the revised proof setup; faint typing beneath it is not reconstructed. Page~585 finishes the argument by expressing the right adjoint along a discrete-fiber cofibration as a product over the fiber. The asserted exactness of those products gives preservation of epimorphisms and hence the required projective statement. The formula is explicitly compared with the direct-sum calculation of section~135. The first remark allows the injective/projective preservation statements with an induced sheaf of rings $\underline O_{A'}=h^*\underline O_A$, while still doubting that this alone suffices for the colocalization comparison (p.~585). The second remark corrects the more basic existence claim. Enough projectives permits deriving an additive functor that already exists; it does not prove that the inverse-image functor has a left adjoint. Even for one-object final categories, extension of scalars need not be left exact or commute with products (p.~586). The source identifies the combined condition with finite projectivity of the relevant coefficient modules, then adopts the simpler induced-ring condition for the intended overview. To recover restriction-of-scalars and colocalization compatibility, it returns to constant, or perhaps locally constant, coefficients. The concluding return to chapter~V therefore has a restricted scope: the author seeks a mutually dual account of $Lf_!,f^*,Rf_*$ and the two internal operations, without yet producing a complete six-operation theory. Section~139 asks whether resolving the integrator instead of the coefficients can actually supply that relative account.


\chapter*{Sources}
Grothendieck, A. (1983). \emph{Pursuing Stacks} [Typescript with handwritten revisions; cotes 134-2--134-8]. Grothendieck archive, Université de Montpellier. The supplied page images and their metadata are the primary witnesses for this edition.

Grothendieck, A. (2021). \emph{Pursuing Stacks} (M. Carmona, with the collaboration of U. Buchholtz, Eds.; arXiv:2111.01000, version 2). Extended from the 2015 edition by ``the scrivener''; CC0 1.0. This edition is consulted as a comparison witness, not reproduced as the source of record.

Brown, R. (Trans.). (n.d.). \emph{English translation of the mathematical part of two letters from Alexander Grothendieck to L. Breen, 1975} [Typescript; prepared with assistance from L. Breen and J-L. Loday]. Grothendieck archive, Université de Montpellier, cote 134-2, sites 115--143.
\end{document}
